\chapter{Algorithms and Complexity}
\label{ch:algorithms}

This chapter covers the foundations of algorithm design, computational complexity, and problem-solving techniques essential for government technical recruitment examinations. It synthesizes theoretical principles, asymptotic analyses, rigorous correctness proofs, and complete C++ implementations across sorting, searching, graph traversal, greedy strategies, recursion, dynamic programming, and high-probability interview problems.

The chapter is systematically organized into three interconnected core areas:
\begin{enumerate}[label=\bfseries Part \Roman*:]
    \item \textbf{Sorting, Searching, and Graph Traversal Algorithms:} Comparative analysis of comparison-based sorts (Bubble, Insertion, Merge, Quick), non-comparison sorts, two-pointer search, Binary Search, Breadth-First Search (BFS), Depth-First Search (DFS), and A* heuristic search.
    \item \textbf{Greedy Algorithms, Recursion, and Dynamic Programming:} Greedy-choice properties, Fractional Knapsack, Activity Selection, Huffman Coding, recursion trees, and dynamic programming formulations (0/1 Knapsack, Longest Common Subsequence, Matrix Chain Multiplication).
    \item \textbf{High-Value Algorithmic Problems and Solutions:} Number theory, prime generation, Euclidean GCD/LCM, series summation, two-sum variations, Bellman-Ford, Floyd-Warshall, Dijkstra, and advanced analytical viva problems.
\end{enumerate}


\begin{center}
\vspace{1em}
{\color{primary}\rule{0.5\textwidth}{1.5pt}}\\[6pt]
{\Large\bfseries\color{primary} Part I: Sorting, Searching, and Graph Algorithms}\\[3pt]
{\color{primary}\rule{0.5\textwidth}{1.5pt}}
\vspace{1em}
\end{center}

This chapter covers fundamental algorithms for technical recruitment examinations: Bubble Sort, Insertion Sort, Merge Sort, Quick Sort, and specialized sorting methods, followed by two-pointer search strategies, Binary Search, Breadth-First Search (BFS), Depth-First Search (DFS), and the informed A* Search algorithm.

The emphasis is on algorithmic design, step-by-step execution traces, pseudocode, complete production-ready C++ implementations, mathematical complexity analysis, heuristic optimality, and common examination pitfalls.

\begin{center}
\vspace{0.5em}
{\color{primary}\rule{0.45\textwidth}{1.2pt}}\\[4pt]
{\Large\bfseries\color{primary} Section I: Fundamentals of Sorting}\\[2pt]
{\color{primary}\rule{0.45\textwidth}{1.2pt}}
\vspace{0.5em}
\end{center}

\section{Fundamentals of Sorting}

\subsection{What is Sorting?}

\textbf{Sorting} is the process of arranging a collection of data in a specific sequential order according to a designated key or comparison criterion.

The two fundamental orderings are:
\begin{itemize}
    \item \textbf{Ascending order (Non-decreasing):} Elements increase from smallest to largest, e.g., $3, 7, 12, 19$.
    \item \textbf{Descending order (Non-increasing):} Elements decrease from largest to smallest, e.g., $19, 12, 7, 3$.
\end{itemize}

Sorting applies across fundamental data types: integers, floating-point numbers, strings, dates, employee records, database rows, and custom objects. For example, a government recruitment portal sorts applicants by merit marks, roll numbers, or alphabetical name order.

\subsection{Why is Sorting Important?}

Sorting dramatically enhances downstream computational efficiency:
\begin{itemize}
    \item \textbf{Binary Search:} Enables logarithmic $O(\log n)$ item location in arrays with random access.
    \item \textbf{Database Indexing:} B-trees and B+ trees maintain sorted key order to service range queries efficiently.
    \item \textbf{Duplicate Detection:} Identifies duplicate records by scanning contiguous neighbors in $O(n)$ time.
    \item \textbf{Algorithmic Prerequisite:} Greedy algorithms (Kruskal's MST, interval scheduling), divide-and-conquer routines, and convex hull algorithms rely directly on pre-sorted data.
\end{itemize}

\section{Important Sorting Terminology}

Technical recruitment examinations regularly test sorting properties beyond raw asymptotic time complexity:

\subsection{Stable Sorting}

A sorting algorithm is \textbf{stable} if elements with equal sorting keys preserve their original relative order in the sorted output.

\begin{table}[htbp]
\centering
\small
\begin{tabular}{c l c}
\toprule
\textbf{Original Index} & \textbf{Candidate Name} & \textbf{Exam Marks} \\
\midrule
1 & Rahim  & 85 \\
2 & Karim  & 92 \\
3 & Salma  & 85 \\
4 & Nabila & 78 \\
\bottomrule
\end{tabular}
\caption{Candidate records prior to sorting by marks.}
\end{table}

Sorting by marks in ascending order:
\begin{itemize}
    \item \textbf{Stable Result:} $\text{Nabila (78)} \to \textbf{Rahim (85)} \to \textbf{Salma (85)} \to \text{Karim (92)}$. Rahim appears before Salma because Rahim preceded Salma in the input.
    \item \textbf{Unstable Result:} $\text{Nabila (78)} \to \textbf{Salma (85)} \to \textbf{Rahim (85)} \to \text{Karim (92)}$. Salma is placed ahead of Rahim, reversing their input chronology.
\end{itemize}

\begin{quote}
\textbf{Engineering Significance:} Stability is vital in multi-key sorting (e.g., sorting first by Student Name, then stably sorting by Department ID, preserving alphabetical order within each department).
\end{quote}

\subsection{In-Place Sorting}

An algorithm is \textbf{in-place} if it transforms the input using $O(1)$ (constant) auxiliary memory beyond the original input container, allowing at most $O(\log n)$ stack space for recursion:
\begin{itemize}
    \item \textbf{In-place algorithms:} Bubble Sort, Insertion Sort, Selection Sort, Heap Sort, and array-based Quick Sort.
    \item \textbf{Not in-place algorithms:} Conventional array Merge Sort ($O(n)$ auxiliary array), Counting Sort, and Radix Sort.
\end{itemize}

\subsection{Adaptive Sorting}

An algorithm is \textbf{adaptive} if its execution time improves when the input sequence is already sorted or nearly sorted:
\begin{itemize}
    \item \textbf{Insertion Sort} is inherently adaptive, running in $O(n)$ linear time on already sorted inputs.
    \item \textbf{Optimized Bubble Sort} (with a swap-detection flag) is adaptive ($O(n)$ on sorted input).
    \item \textbf{Merge Sort} is non-adaptive in its standard top-down implementation, performing $\Theta(n \log n)$ operations regardless of initial order.
\end{itemize}

\subsection{Comparison-Based vs Non-Comparison-Based Sorting}

\begin{itemize}
    \item \textbf{Comparison-Based Sorts:} Determine relative order exclusively by pairwise comparisons ($A[i] \leq A[j]$). Examples: Bubble, Insertion, Selection, Merge, Quick, and Heap Sort. \textbf{Theoretical Lower Bound:} $\Omega(n \log n)$ worst-case time.
    \item \textbf{Non-Comparison-Based Sorts:} Exploit intrinsic mathematical properties or digit representations of keys (e.g., bounded integer ranges). Examples: Counting Sort, Radix Sort, Bucket Sort. Can achieve linear $O(n)$ time.
\end{itemize}

\subsection{Internal vs External Sorting}

\begin{itemize}
    \item \textbf{Internal Sorting:} All data resides entirely inside primary random-access memory (RAM).
    \item \textbf{External Sorting:} The dataset exceeds RAM capacity and resides on secondary external storage (SSD/HDD). Algorithms (e.g., External Merge Sort) optimize sequential block I/O.
\end{itemize}

\newpage
\begin{center}
\vspace{0.5em}
{\color{primary}\rule{0.45\textwidth}{1.2pt}}\\[4pt]
{\Large\bfseries\color{primary} Section II: Bubble Sort}\\[2pt]
{\color{primary}\rule{0.45\textwidth}{1.2pt}}
\vspace{0.5em}
\end{center}

\section{Bubble Sort}

\textbf{Bubble Sort} is a simple comparison-based sorting algorithm that repeatedly steps through the array, compares adjacent elements, and swaps them if they are in the wrong order. With each complete pass, the largest unsorted element "bubbles up" to its correct position at the end of the array.

\subsection{Algorithmic Steps}
\begin{enumerate}
    \item Traverse the array from index $0$ to $n-2$.
    \item Compare adjacent pair $(A[j], A[j+1])$. If $A[j] > A[j+1]$, swap them.
    \item At the end of pass $i$, the $(i+1)^{\text{th}}$ largest element is locked into position $n - 1 - i$.
    \item If a complete pass finishes without a single swap, terminate early (array is sorted).
\end{enumerate}

\subsection{Step-by-Step Example}
Sort $A = [5, 1, 4, 2, 8]$ in ascending order:
\begin{itemize}
    \item \textbf{Pass 1:}
    \begin{itemize}
        \item Compare $5, 1 \implies \text{swap} \implies [1, 5, 4, 2, 8]$
        \item Compare $5, 4 \implies \text{swap} \implies [1, 4, 5, 2, 8]$
        \item Compare $5, 2 \implies \text{swap} \implies [1, 4, 2, 5, 8]$
        \item Compare $5, 8 \implies \text{no swap} \implies [1, 4, 2, 5, 8]$
        \item Largest element $8$ is in its permanent position.
    \end{itemize}
    \item \textbf{Pass 2:}
    \begin{itemize}
        \item Compare $1, 4 \implies \text{no swap} \implies [1, 4, 2, 5, 8]$
        \item Compare $4, 2 \implies \text{swap} \implies [1, 2, 4, 5, 8]$
        \item Compare $4, 5 \implies \text{no swap} \implies [1, 2, 4, 5, 8]$
        \item Second largest element $5$ is locked into place.
    \end{itemize}
    \item \textbf{Pass 3:}
    \begin{itemize}
        \item Compare $1, 2 \implies \text{no swap}$; Compare $2, 4 \implies \text{no swap}$.
        \item No swaps occurred $\implies$ \textbf{Early termination triggered!}
    \end{itemize}
\end{itemize}
Final sorted array: $\mathbf{[1, 2, 4, 5, 8]}$.

\subsection{Pseudocode}

\begin{lstlisting}[style=pseudocode, caption={Optimized Bubble Sort Pseudocode.}]
BUBBLE_SORT(A, n):
    for i = 0 to n - 2:
        swapped = false
        for j = 0 to n - 2 - i:
            if A[j] > A[j + 1]:
                swap(A[j], A[j + 1])
                swapped = true
        if swapped == false:
            break
\end{lstlisting}

\subsection{C++ Implementation}

\begin{lstlisting}[style=cppstyle, caption={C++ Optimized Bubble Sort Implementation.}]
#include <iostream>
#include <vector>
#include <utility>
using namespace std;

void bubbleSort(vector<int>& a) {
    int n = static_cast<int>(a.size());
    for (int i = 0; i < n - 1; i++) {
        bool swapped = false;
        for (int j = 0; j < n - 1 - i; j++) {
            if (a[j] > a[j + 1]) {
                swap(a[j], a[j + 1]);
                swapped = true;
            }
        }
        if (!swapped) break; // Early termination
    }
}

int main() {
    vector<int> a = {5, 1, 4, 2, 8};
    bubbleSort(a);
    for (int x : a) cout << x << " ";
    cout << '\n';
    return 0;
}
\end{lstlisting}

\subsection{Complexity Analysis}
\begin{itemize}
    \item \textbf{Worst-Case Time:} Occurs when the array is in reverse sorted order. Total comparisons:
    \[
    (n-1) + (n-2) + \dots + 1 = \frac{n(n-1)}{2} \implies \boxed{O(n^2)}
    \]
    \item \textbf{Average-Case Time:} Quadratic pairwise inversions $\implies \boxed{O(n^2)}$.
    \item \textbf{Best-Case Time:} Occurs when array is already sorted. The \texttt{swapped} flag triggers exit after $n-1$ comparisons $\implies \boxed{O(n)}$.
    \item \textbf{Auxiliary Space:} Modifies the array in place $\implies \boxed{O(1)}$.
\end{itemize}

\subsection{Properties Summary}

\begin{table}[htbp]
\centering
\small
\begin{tabularx}{\textwidth}{l X}
\toprule
\textbf{Property} & \textbf{Characteristics} \\
\midrule
Best-Case Time    & $O(n)$ (with \texttt{swapped} flag optimization) \\
Average-Case Time & $O(n^2)$ \\
Worst-Case Time   & $O(n^2)$ \\
Auxiliary Space   & $O(1)$ (In-place) \\
Stability         & \textbf{Stable} (equal elements are never swapped when using $A[j] > A[j+1]$) \\
Adaptivity        & \textbf{Adaptive} (benefits from existing order) \\
Practical Usage   & Primarily educational; useful for tiny or nearly sorted lists \\
\bottomrule
\end{tabularx}
\caption{Properties of Bubble Sort.}
\end{table}

\newpage
\begin{center}
\vspace{0.5em}
{\color{primary}\rule{0.45\textwidth}{1.2pt}}\\[4pt]
{\Large\bfseries\color{primary} Section III: Insertion Sort}\\[2pt]
{\color{primary}\rule{0.45\textwidth}{1.2pt}}
\vspace{0.5em}
\end{center}

\section{Insertion Sort}

\textbf{Insertion Sort} incrementally constructs a sorted subarray at the front of the array. For each new element (the \texttt{key}), it scans backward through the sorted prefix, shifts larger elements one position to the right, and inserts the key into its correct relative position.

\subsection{Algorithmic Steps}
\begin{enumerate}
    \item The first element $A[0]$ is trivially sorted.
    \item For $i = 1$ to $n-1$: designate $\text{key} = A[i]$ and set $j = i - 1$.
    \item While $j \geq 0$ and $A[j] > \text{key}$: shift $A[j+1] = A[j]$, and decrement $j$.
    \item Place $\text{key}$ at the vacated slot: $A[j+1] = \text{key}$.
\end{enumerate}

\subsection{Step-by-Step Example}
Sort $A = [5, 1, 4, 2, 8]$:
\begin{itemize}
    \item Initial sorted prefix: $[5]$.
    \item \textbf{Pass 1 ($i = 1$, key = 1):} $5 > 1 \implies$ shift $5 \implies$ insert $1 \implies \mathbf{[1, 5}, 4, 2, 8]$.
    \item \textbf{Pass 2 ($i = 2$, key = 4):} $5 > 4 \implies$ shift $5$; $1 < 4 \implies$ stop $\implies$ insert $4 \implies \mathbf{[1, 4, 5}, 2, 8]$.
    \item \textbf{Pass 3 ($i = 3$, key = 2):} $5 > 2, 4 > 2 \implies$ shift $5, 4$; $1 < 2 \implies$ stop $\implies$ insert $2 \implies \mathbf{[1, 2, 4, 5}, 8]$.
    \item \textbf{Pass 4 ($i = 4$, key = 8):} $8 > 5 \implies$ no shifts required $\implies \mathbf{[1, 2, 4, 5, 8]}$.
\end{itemize}

\subsection{Pseudocode}

\begin{lstlisting}[style=pseudocode, caption={Insertion Sort Pseudocode.}]
INSERTION_SORT(A, n):
    for i = 1 to n - 1:
        key = A[i]
        j = i - 1
        while j >= 0 and A[j] > key:
            A[j + 1] = A[j]
            j = j - 1
        A[j + 1] = key
\end{lstlisting}

\begin{quote}
\textbf{Stability Condition:} Using strict inequality \texttt{A[j] > key} ensures stability. If non-strict \texttt{A[j] >= key} were used, equal keys would shift past each other, destroying stability!
\end{quote}

\subsection{C++ Implementation}

\begin{lstlisting}[style=cppstyle, caption={C++ Insertion Sort Implementation.}]
#include <iostream>
#include <vector>
using namespace std;

void insertionSort(vector<int>& a) {
    int n = static_cast<int>(a.size());
    for (int i = 1; i < n; i++) {
        int key = a[i];
        int j = i - 1;
        while (j >= 0 && a[j] > key) {
            a[j + 1] = a[j];
            j--;
        }
        a[j + 1] = key;
    }
}

int main() {
    vector<int> a = {5, 1, 4, 2, 8};
    insertionSort(a);
    for (int x : a) cout << x << " ";
    cout << '\n';
    return 0;
}
\end{lstlisting}

\subsection{Complexity and Properties}

\begin{table}[htbp]
\centering
\small
\begin{tabularx}{\textwidth}{l X}
\toprule
\textbf{Property} & \textbf{Characteristics} \\
\midrule
Best-Case Time    & $\boxed{O(n)}$ (Array already sorted; inner loop test fails immediately) \\
Average-Case Time & $\boxed{O(n^2)}$ (Average $n^2/4$ shifts) \\
Worst-Case Time   & $\boxed{O(n^2)}$ (Reverse sorted array; $\frac{n(n-1)}{2}$ shifts) \\
Auxiliary Space   & $\boxed{O(1)}$ (In-place) \\
Stability         & \textbf{Stable} \\
Online Property   & \textbf{Online} (can sort a live stream of data as elements arrive) \\
Practical Usage   & Base sort for small subarrays ($n \leq 16$--$32$) in hybrid algorithms (TimSort, Introsort) \\
\bottomrule
\end{tabularx}
\caption{Properties of Insertion Sort.}
\end{table}

\newpage
\begin{center}
\vspace{0.5em}
{\color{primary}\rule{0.45\textwidth}{1.2pt}}\\[4pt]
{\Large\bfseries\color{primary} Section IV: Merge Sort}\\[2pt]
{\color{primary}\rule{0.45\textwidth}{1.2pt}}
\vspace{0.5em}
\end{center}

\section{Merge Sort}

\textbf{Merge Sort} is an optimal divide-and-conquer sorting algorithm. It recursively bisects the array into halves until subarrays contain single elements, and then merges the sorted halves into a complete sorted array in linear time.

\subsection{Algorithmic Structure}
\begin{enumerate}
    \item \textbf{Divide:} Bisect the array into two halves at midpoint $\text{mid} = \text{left} + (\text{right} - \text{left})/2$.
    \item \textbf{Conquer:} Recursively invoke \texttt{mergeSort} on left half $[\text{left}, \text{mid}]$ and right half $[\text{mid}+1, \text{right}]$.
    \item \textbf{Combine:} Merge the two sorted subarrays using a linear two-pointer merge procedure.
\end{enumerate}

\subsection{Step-by-Step Example}
Sort $A = [5, 1, 4, 2, 8]$:
\begin{enumerate}
    \item \textbf{Divide Tree:}
    \[
    [5, 1, 4, 2, 8] \implies [5, 1] \text{ and } [4, 2, 8]
    \]
    \[
    [5, 1] \implies [5], [1]; \qquad [4, 2, 8] \implies [4], [2, 8] \implies [2], [8]
    \]
    \item \textbf{Merge Stage 1:}
    \begin{itemize}
        \item Merge $[5]$ and $[1] \implies [1, 5]$.
        \item Merge $[2]$ and $[8] \implies [2, 8]$.
    \end{itemize}
    \item \textbf{Merge Stage 2:}
    \begin{itemize}
        \item Merge $[4]$ with $[2, 8] \implies [2, 4, 8]$.
    \end{itemize}
    \item \textbf{Merge Stage 3 (Final Merge):}
    \begin{itemize}
        \item Merge $[1, 5]$ and $[2, 4, 8]$:
        \item Compare $1 < 2 \implies \text{take } 1$
        \item Compare $5 > 2 \implies \text{take } 2$
        \item Compare $5 > 4 \implies \text{take } 4$
        \item Compare $5 < 8 \implies \text{take } 5$
        \item Append remainder $8 \implies \mathbf{[1, 2, 4, 5, 8]}$.
    \end{itemize}
\end{enumerate}

\subsection{C++ Implementation}

\begin{lstlisting}[style=cppstyle, caption={C++ Merge Sort Implementation.}]
#include <iostream>
#include <vector>
using namespace std;

void merge(vector<int>& a, int left, int mid, int right) {
    vector<int> temp;
    temp.reserve(right - left + 1);

    int i = left;
    int j = mid + 1;

    while (i <= mid && j <= right) {
        if (a[i] <= a[j]) { // '<=' ensures stability
            temp.push_back(a[i++]);
        } else {
            temp.push_back(a[j++]);
        }
    }

    while (i <= mid)   temp.push_back(a[i++]);
    while (j <= right) temp.push_back(a[j++]);

    for (int k = 0; k < static_cast<int>(temp.size()); k++) {
        a[left + k] = temp[k];
    }
}

void mergeSort(vector<int>& a, int left, int right) {
    if (left >= right) return;
    int mid = left + (right - left) / 2;
    mergeSort(a, left, mid);
    mergeSort(a, mid + 1, right);
    merge(a, left, mid, right);
}

int main() {
    vector<int> a = {5, 1, 4, 2, 8};
    mergeSort(a, 0, static_cast<int>(a.size()) - 1);
    for (int x : a) cout << x << " ";
    cout << '\n';
    return 0;
}
\end{lstlisting}

\subsection{Complexity Analysis}
\begin{itemize}
    \item \textbf{Recurrence Relation:}
    \[
    \boxed{T(n) = 2T(n/2) + O(n)}
    \]
    By the Master Theorem (Case 2: $a=2, b=2, f(n)=n \implies n^{\log_2 2} = n$), the solution is:
    \[
    \boxed{T(n) = O(n \log n)}
    \]
    This bound holds strictly across \textbf{Best, Average, and Worst Cases}.
    \item \textbf{Space Complexity:} The merge procedure requires an auxiliary buffer of size $n$, yielding $\boxed{O(n)}$ auxiliary space for array sorting.
\end{itemize}

\subsection{Properties and Architectural Strengths}

\begin{table}[htbp]
\centering
\small
\begin{tabularx}{\textwidth}{l X}
\toprule
\textbf{Property} & \textbf{Characteristics} \\
\midrule
Best / Average / Worst Time & $\boxed{O(n \log n)}$ guaranteed \\
Auxiliary Space             & $O(n)$ for arrays; $O(\log n)$ call stack \\
Stability                   & \textbf{Stable} (preserves order when $A[i] \leq A[j]$) \\
In-Place?                   & \textbf{No} (requires $O(n)$ temporary buffer for arrays) \\
Linked Lists Suitability    & \textbf{Ideal}: Linked lists can be merged in $O(1)$ extra space by pointer redirection without array shifting! \\
External Sorting            & \textbf{Standard choice}: Excellent sequential streaming access patterns on secondary disks. \\
\bottomrule
\end{tabularx}
\caption{Properties of Merge Sort.}
\end{table}

\newpage
\begin{center}
\vspace{0.5em}
{\color{primary}\rule{0.45\textwidth}{1.2pt}}\\[4pt]
{\Large\bfseries\color{primary} Section V: Quick Sort}\\[2pt]
{\color{primary}\rule{0.45\textwidth}{1.2pt}}
\vspace{0.5em}
\end{center}

\section{Quick Sort}

\textbf{Quick Sort} is a divide-and-conquer algorithm that selects a \textbf{pivot element} and partitions the array such that elements smaller than the pivot are placed to its left, and elements greater are placed to its right. The pivot is locked into its final sorted position, and the partitions are sorted recursively.

\subsection{Pivot Selection Strategies}
\begin{itemize}
    \item \textbf{Last element / First element:} Simplest, but degrades to $O(n^2)$ on already sorted arrays.
    \item \textbf{Randomized Pivot:} Selects a random index; guarantees $O(n \log n)$ expected time regardless of input ordering.
    \item \textbf{Median-of-Three:} Selects the median of $\{\text{first}, \text{middle}, \text{last}\}$ elements.
\end{itemize}

\subsection{Lomuto Partition Scheme}
Using the last element as pivot:
\begin{lstlisting}[style=pseudocode, caption={Lomuto Partition Algorithm.}]
PARTITION(A, low, high):
    pivot = A[high]
    i = low - 1
    for j = low to high - 1:
        if A[j] <= pivot:
            i = i + 1
            swap(A[i], A[j])
    swap(A[i + 1], A[high])
    return i + 1
\end{lstlisting}

\subsection{C++ Implementation}

\begin{lstlisting}[style=cppstyle, caption={C++ Quick Sort Implementation.}]
#include <iostream>
#include <vector>
#include <utility>
using namespace std;

int partitionArray(vector<int>& a, int low, int high) {
    int pivot = a[high];
    int i = low - 1;
    for (int j = low; j < high; j++) {
        if (a[j] <= pivot) {
            i++;
            swap(a[i], a[j]);
        }
    }
    swap(a[i + 1], a[high]);
    return i + 1;
}

void quickSort(vector<int>& a, int low, int high) {
    if (low >= high) return;
    int pivotIndex = partitionArray(a, low, high);
    quickSort(a, low, pivotIndex - 1);
    quickSort(a, pivotIndex + 1, high);
}

int main() {
    vector<int> a = {5, 1, 4, 2, 8};
    quickSort(a, 0, static_cast<int>(a.size()) - 1);
    for (int x : a) cout << x << " ";
    cout << '\n';
    return 0;
}
\end{lstlisting}

\subsection{Complexity Analysis}
\begin{itemize}
    \item \textbf{Best Case ($O(n \log n)$):} Pivot divides array into balanced halves ($n/2, n/2$):
    \[
    T(n) = 2T(n/2) + O(n) \implies O(n \log n)
    \]
    \item \textbf{Average Case ($O(n \log n)$):} Even unbalanced splits (e.g., $90\% / 10\%$) yield logarithmic tree depth $\implies O(n \log n)$.
    \item \textbf{Worst Case ($O(n^2)$):} Occurs when pivot is always the smallest or largest element (e.g., already sorted array with last-element pivot):
    \[
    T(n) = T(n-1) + O(n) \implies O(n^2)
    \]
    \item \textbf{Space Complexity:} In-place array modification, but recursion stack uses:
    \[
    \boxed{O(\log n) \text{ average}, \quad O(n) \text{ worst-case}}
    \]
\end{itemize}

\subsection{Properties Summary}

\begin{table}[htbp]
\centering
\small
\begin{tabularx}{\textwidth}{l X}
\toprule
\textbf{Property} & \textbf{Characteristics} \\
\midrule
Best-Case Time    & $O(n \log n)$ \\
Average-Case Time & $O(n \log n)$ \\
Worst-Case Time   & $\boxed{O(n^2)}$ \\
Auxiliary Space   & $O(\log n)$ average recursion stack; $O(n)$ worst-case \\
Stability         & \textbf{Unstable} (long-distance swaps disrupt relative order of equal keys) \\
In-Place?         & \textbf{Yes} (partitions array in place) \\
Performance Note  & Best cache locality and lowest constant factors among $O(n \log n)$ comparison sorts \\
\bottomrule
\end{tabularx}
\caption{Properties of Quick Sort.}
\end{table}

\newpage
\begin{center}
\vspace{0.5em}
{\color{primary}\rule{0.45\textwidth}{1.2pt}}\\[4pt]
{\Large\bfseries\color{primary} Section VI: Comparative Analysis of Core Algorithms}\\[2pt]
{\color{primary}\rule{0.45\textwidth}{1.2pt}}
\vspace{0.5em}
\end{center}

\section{Comparative Analysis of Core Algorithms}

\subsection{Master Evaluation of the Four Fundamental Sorts}

\begin{table}[htbp]
\centering
\footnotesize
\begin{tabularx}{\textwidth}{p{2.6cm} X X X X}
\toprule
\textbf{Feature} & \textbf{Bubble Sort} & \textbf{Insertion Sort} & \textbf{Merge Sort} & \textbf{Quick Sort} \\
\midrule
Paradigm & Adjacent comparisons & Incremental sorted prefix insertion & Divide-and-Conquer & Partition around Pivot \\
Best Time & $O(n)$ (optimized) & $O(n)$ & $O(n \log n)$ & $O(n \log n)$ \\
Average Time & $O(n^2)$ & $O(n^2)$ & $O(n \log n)$ & $O(n \log n)$ \\
Worst Time & $O(n^2)$ & $O(n^2)$ & $O(n \log n)$ & $O(n^2)$ \\
Space & $O(1)$ & $O(1)$ & $O(n)$ (Buffer) & $O(\log n)$ (Stack) \\
Stability & \textbf{Stable} & \textbf{Stable} & \textbf{Stable} & \textbf{Unstable} \\
In-Place & \textbf{Yes} & \textbf{Yes} & \textbf{No} & \textbf{Yes} \\
Adaptive & \textbf{Yes} & \textbf{Yes} & \textbf{No} & \textbf{No} \\
Primary Strength & Simple to teach & Fastest on small / nearly sorted data & Predictable $O(n \log n)$ and stable & High cache speed, in-place \\
Primary Weakness & Sluggish $O(n^2)$ & Sluggish $O(n^2)$ on large data & Requires $O(n)$ extra memory & Unstable; $O(n^2)$ worst case \\
\bottomrule
\end{tabularx}
\caption{Comprehensive comparison of the four fundamental sorting algorithms.}
\end{table}

\subsection{Strategic Decision Criteria}
\begin{itemize}
    \item \textbf{Merge Sort vs Quick Sort:} Merge Sort guarantees $O(n \log n)$ and stability at the cost of $O(n)$ memory. Quick Sort avoids extra memory and provides superior cache performance, but can degrade to $O(n^2)$ and is unstable.
    \item \textbf{Insertion Sort vs Bubble Sort:} Both are $O(n^2)$ average and $O(n)$ best, but Insertion Sort performs far fewer writes and comparisons, making it universally preferred for small arrays.
\end{itemize}

\newpage
\begin{center}
\vspace{0.5em}
{\color{primary}\rule{0.45\textwidth}{1.2pt}}\\[4pt]
{\Large\bfseries\color{primary} Section VII: Other Important Sorting Algorithms}\\[2pt]
{\color{primary}\rule{0.45\textwidth}{1.2pt}}
\vspace{0.5em}
\end{center}

\section{Other Important Sorting Algorithms}

\subsection{Selection Sort}
Repeatedly identifies the minimum element from the unsorted section and swaps it with the first unsorted element.
\begin{itemize}
    \item \textbf{Time Complexity:} $\boxed{\Theta(n^2)}$ in Best, Average, and Worst cases (always performs $\frac{n(n-1)}{2}$ comparisons).
    \item \textbf{Auxiliary Space:} $O(1)$ (In-place).
    \item \textbf{Stability:} \textbf{Unstable} (e.g., sorting $[4a, 4b, 1]$ swaps $4a$ with $1$, placing $4a$ after $4b$).
    \item \textbf{Key Feature:} Performs at most $\boxed{n - 1 \text{ swaps}}$, making it optimal when \textbf{memory write operations are exceptionally expensive}.
\end{itemize}

\subsection{Heap Sort}
Builds a max-heap in $O(n)$ time, then repeatedly extracts the root maximum and sifts down:
\begin{itemize}
    \item \textbf{Time Complexity:} $\boxed{O(n \log n)}$ in Best, Average, and Worst cases.
    \item \textbf{Auxiliary Space:} $\boxed{O(1)}$ (In-place array manipulation).
    \item \textbf{Stability:} \textbf{Unstable}.
    \item \textbf{Comparison with Quick Sort:} Guarantees $O(n \log n)$ without extra memory, but exhibits poorer hardware cache locality due to jumping indices ($2i + 1, 2i + 2$).
\end{itemize}

\subsection{Shell Sort}
Diminishing increment sort extending Insertion Sort by comparing elements separated by a \textbf{gap} $h$, shrinking $h$ until $h = 1$:
\begin{itemize}
    \item \textbf{Complexity:} Strongly depends on the gap sequence (e.g., Shell's original $N/2^k$ is $O(n^2)$; Knuth's $(3^k - 1)/2$ is $O(n^{3/2})$; Sedgewick's sequence achieves $O(n^{4/3})$).
    \item \textbf{Space:} $O(1)$ in-place; \textbf{Unstable}.
\end{itemize}

\subsection{Counting Sort}
Non-comparison integer sorting that counts occurrences of each value in range $[0, k]$:
\begin{itemize}
    \item \textbf{Time Complexity:} $\boxed{O(n + k)}$. Linear when $k = O(n)$.
    \item \textbf{Auxiliary Space:} $\boxed{O(n + k)}$ using frequency count and output arrays.
    \item \textbf{Stability:} \textbf{Stable} (when iterating input backward during placement).
    \item \textbf{Constraint:} Inefficient when $k \gg n$ (e.g., sorting 10 keys in range $1$ to $1,000,000$).
\end{itemize}

\subsection{Radix Sort}
Processes keys digit-by-digit from Least Significant Digit (LSD) to Most Significant Digit (MSD) using a stable sub-sort (typically Counting Sort):
\begin{itemize}
    \item \textbf{Time Complexity:} $\boxed{O(d(n + b))}$, where $d$ is digit count and $b$ is base ($b = 10$).
    \item \textbf{Auxiliary Space:} $O(n + b)$; \textbf{Stable}.
    \item \textbf{Linear Time:} For fixed-width 32-bit or 64-bit integers, $d$ is constant, yielding linear $O(n)$ sorting!
\end{itemize}

\subsection{Bucket Sort}
Distributes $n$ uniformly distributed elements over $[0, 1)$ into $k$ buckets using index $\lfloor n \cdot x \rfloor$, sorts individual buckets, and concatenates them:
\begin{itemize}
    \item \textbf{Time Complexity:} Expected $\boxed{O(n + k)}$ on uniform inputs; Worst-case $\boxed{O(n^2)}$ under heavy clustering.
    \item \textbf{Space:} $O(n + k)$.
\end{itemize}

\subsection{TimSort}
State-of-the-art hybrid algorithm combining Insertion Sort and Merge Sort:
\begin{itemize}
    \item Identifies natural contiguous ascending/descending subsequences (\textbf{runs}).
    \item Extends short runs using Insertion Sort, then merges them with galloping mode.
    \item \textbf{Time Complexity:} Best $O(n)$ (strictly sorted data); Average/Worst $\boxed{O(n \log n)}$.
    \item \textbf{Space:} $O(n)$; \textbf{Stable}.
    \item \textbf{Standard Adoption:} Default sorting engine in Python (\texttt{list.sort}), Java (\texttt{Arrays.sort} for objects), and Android.
\end{itemize}

\subsection{Introsort (Introspective Sort)}
Hybrid algorithm providing Quick Sort's speed with guaranteed $O(n \log n)$ worst-case safety:
\begin{itemize}
    \item Begins with Quick Sort.
    \item Monitors recursion depth; if depth exceeds $2 \lfloor \log_2 n \rfloor$, switches automatically to \textbf{Heap Sort}.
    \item Switches to \textbf{Insertion Sort} for tiny partitions ($n < 16$).
    \item \textbf{Time Complexity:} Guaranteed $\boxed{O(n \log n)}$ worst-case; Space $O(\log n)$; \textbf{Unstable}.
    \item \textbf{Standard Adoption:} Standard implementation of C++ \texttt{std::sort()}.
\end{itemize}

\subsection{Tree Sort}
Inserts all $n$ items into a Binary Search Tree (BST), then traverses it \textbf{inorder}:
\begin{itemize}
    \item \textbf{Time Complexity:} $O(n \log n)$ average; degrades to $\boxed{O(n^2)}$ on skewed trees. With self-balancing trees (AVL/Red-Black), guarantees $O(n \log n)$.
    \item \textbf{Space:} $O(n)$ for tree nodes.
\end{itemize}

\subsection{External Merge Sort}
Designed for massive datasets exceeding physical RAM:
\begin{enumerate}
    \item \textbf{Run Generation:} Stream chunks of data into RAM, sort internally, and write sorted runs to disk.
    \item \textbf{Multi-Way Merge:} Concurrently merge $k$ sorted runs into a unified output stream using min-heaps.
\end{enumerate}

\newpage
\begin{center}
\vspace{0.5em}
{\color{primary}\rule{0.45\textwidth}{1.2pt}}\\[4pt]
{\Large\bfseries\color{primary} Section VIII: Master Comparative Reference Table}\\[2pt]
{\color{primary}\rule{0.45\textwidth}{1.2pt}}
\vspace{0.5em}
\end{center}

\section{Master Comparative Reference Table}

\begin{table}[htbp]
\centering
\footnotesize
\begin{tabularx}{\textwidth}{p{3.2cm} p{1.8cm} p{1.8cm} p{1.8cm} p{2.6cm} c c}
\toprule
\textbf{Algorithm} & \textbf{Best Time} & \textbf{Average Time} & \textbf{Worst Time} & \textbf{Auxiliary Space} & \textbf{Stable?} & \textbf{In-Place?} \\
\midrule
Bubble Sort (Opt) & $O(n)$ & $O(n^2)$ & $O(n^2)$ & $O(1)$ & Yes & Yes \\
Insertion Sort    & $O(n)$ & $O(n^2)$ & $O(n^2)$ & $O(1)$ & Yes & Yes \\
Selection Sort    & $O(n^2)$ & $O(n^2)$ & $O(n^2)$ & $O(1)$ & No & Yes \\
Merge Sort        & $O(n \log n)$ & $O(n \log n)$ & $O(n \log n)$ & $O(n)$ & Yes & No \\
Quick Sort        & $O(n \log n)$ & $O(n \log n)$ & $O(n^2)$ & $O(\log n)$ stack & No & Yes \\
Heap Sort         & $O(n \log n)$ & $O(n \log n)$ & $O(n \log n)$ & $O(1)$ & No & Yes \\
Shell Sort        & $O(n \log n)$ & $O(n^{4/3})$ & $O(n^2)$ & $O(1)$ & No & Yes \\
Counting Sort     & $O(n + k)$ & $O(n + k)$ & $O(n + k)$ & $O(n + k)$ & Yes & No \\
Radix Sort        & $O(d(n + b))$ & $O(d(n + b))$ & $O(d(n + b))$ & $O(n + b)$ & Yes & No \\
Bucket Sort       & $O(n + k)$ & $O(n + k)$ & $O(n^2)$ & $O(n + k)$ & Yes & No \\
TimSort           & $O(n)$ & $O(n \log n)$ & $O(n \log n)$ & $O(n)$ & Yes & No \\
Introsort         & $O(n \log n)$ & $O(n \log n)$ & $O(n \log n)$ & $O(\log n)$ stack & No & Yes \\
Tree Sort (BST)   & $O(n \log n)$ & $O(n \log n)$ & $O(n^2)$ & $O(n)$ & No & No \\
Pigeonhole Sort   & $O(n + k)$ & $O(n + k)$ & $O(n + k)$ & $O(n + k)$ & Yes & No \\
Comb Sort         & $O(n \log n)$ & $O(n^2 / 2^p)$ & $O(n^2)$ & $O(1)$ & No & Yes \\
\bottomrule
\end{tabularx}
\caption{Master comparative table of 15 sorting algorithms.}
\end{table}

\subsection{Fundamental Comparison-Sorting Theorems}

\begin{tcolorbox}[title=\textbf{Theorem: Lower Bound for Comparison-Based Sorting}]
Any comparison-based sorting algorithm requires at least $\mathbf{\Omega(n \log n)}$ comparisons in the worst case.\\
\textbf{Proof Sketch:} A decision tree sorting $n$ distinct keys must have at least $n!$ leaves (one for each permutation). A binary decision tree of height $h$ has at most $2^h$ leaves:
\[
2^h \geq n! \implies h \geq \log_2(n!)
\]
By Stirling's approximation ($\ln n! \approx n \ln n - n$):
\[
h \geq \sum_{i=1}^n \log_2 i \geq \sum_{i=n/2}^n \log_2(n/2) = \frac{n}{2} \log_2\left(\frac{n}{2}\right) = \Omega(n \log n).
\]
\end{tcolorbox}

\newpage
\begin{center}
\vspace{0.5em}
{\color{primary}\rule{0.45\textwidth}{1.2pt}}\\[4pt]
{\Large\bfseries\color{primary} Section IX: Solved Examination Problems}\\[2pt]
{\color{primary}\rule{0.45\textwidth}{1.2pt}}
\vspace{0.5em}
\end{center}

\section{Solved Examination Problems}

\subsection{Problem 1: Trace Bubble Sort Swaps}
\textbf{Problem Statement:} Sort $A = [7, 3, 5, 2]$ using Bubble Sort. Determine intermediate arrays and total swap count.
\begin{itemize}
    \item \textbf{Pass 1:}
    \begin{itemize}
        \item Swap $7, 3 \implies [3, 7, 5, 2]$ (Swap 1)
        \item Swap $7, 5 \implies [3, 5, 7, 2]$ (Swap 2)
        \item Swap $7, 2 \implies [3, 5, 2, 7]$ (Swap 3)
    \end{itemize}
    \item \textbf{Pass 2:}
    \begin{itemize}
        \item Compare $3, 5 \implies$ no swap
        \item Swap $5, 2 \implies [3, 2, 5, 7]$ (Swap 4)
    \end{itemize}
    \item \textbf{Pass 3:}
    \begin{itemize}
        \item Swap $3, 2 \implies [2, 3, 5, 7]$ (Swap 5)
    \end{itemize}
\end{itemize}
Total Swaps: $\boxed{3 + 1 + 1 = 5 \text{ swaps}}$.

\subsection{Problem 2: Count Inversions in an Array}
An \textbf{inversion} is a pair $(i, j)$ such that $i < j$ and $A[i] > A[j]$.
\begin{itemize}
    \item For $A = [3, 1, 2]$: Inversions are $(3, 1)$ and $(3, 2) \implies \boxed{\text{Inversion Count} = 2}$.
    \item Maximum inversions in an array of size $n$:
    \[
    \boxed{I_{\max} = \frac{n(n-1)}{2}} \quad (\text{occurs on strictly reverse-sorted arrays})
    \]
    For $n = 5$ ($[5, 4, 3, 2, 1]$): $I_{\max} = \frac{5 \times 4}{2} = \boxed{10 \text{ inversions}}$.
\end{itemize}

\subsection{Problem 3: Identify the Algorithm from Its Properties}
\textbf{Question:} An algorithm runs in $O(n \log n)$ worst-case time, operates in $O(1)$ auxiliary memory, and is not stable. Which algorithm is it?\\
\textbf{Answer: Heap Sort.}
\begin{itemize}
    \item Merge Sort requires $O(n)$ space.
    \item Quick Sort has $O(n^2)$ worst case and $O(\log n)$ stack space.
    \item Insertion Sort has $O(n^2)$ worst case.
\end{itemize}

\subsection{Problem 4: Strategic Algorithm Selection}

\begin{table}[htbp]
\centering
\small
\begin{tabularx}{\textwidth}{p{6.5cm} p{3.5cm} X}
\toprule
\textbf{Real-World Operational Scenario} & \textbf{Optimal Algorithm} & \textbf{Justification} \\
\midrule
Tiny or nearly sorted array & \textbf{Insertion Sort} & Adaptive $O(n)$ time, minimal constant overhead \\
Large array needing $O(n \log n)$ guarantee and $O(1)$ space & \textbf{Heap Sort} & Worst-case optimal and in-place \\
Large dataset requiring stable sorting & \textbf{Merge Sort} & Stable and predictable $O(n \log n)$ \\
Bounded integers in small range ($k \approx n$) & \textbf{Counting Sort} & Linear $O(n + k)$ non-comparison speed \\
Fixed-width integers / Strings by digits & \textbf{Radix Sort} & Linear $O(d(n + b))$ time \\
General in-memory sorting & \textbf{Quick Sort / Introsort} & Fastest practical cache execution \\
Dataset larger than physical RAM & \textbf{External Merge Sort} & Optimized sequential disk block transfers \\
Partially ordered real-world workloads & \textbf{TimSort} & Exploits natural ordered runs \\
Strictly minimize physical memory write operations & \textbf{Selection Sort} & Performs at most $n - 1$ swaps \\
\bottomrule
\end{tabularx}
\caption{Strategic sorting algorithm selection chart.}
\end{table}

\subsection{Problem 5: Determine Sorting Stability}
\textbf{Question:} Input records: $(\text{Rahim}, 80), (\text{Karim}, 90), (\text{Salma}, 80)$.\\
Sorted ascending by marks:
\[
\text{Option A: } (\text{Rahim}, 80), (\text{Salma}, 80), (\text{Karim}, 90)
\]
\[
\text{Option B: } (\text{Salma}, 80), (\text{Rahim}, 80), (\text{Karim}, 90)
\]
\textbf{Answer: Option A is stable.} Rahim appeared before Salma in the original input; preserving this relative order confirms stability.

\subsection{Problem 6: C++ Standard Library Guarantees}
\begin{itemize}
    \item \texttt{std::sort()}: Guarantees $\mathbf{O(n \log n)}$ comparisons in the worst case (C++11 standard). Typically implemented via \textbf{Introsort}. It is \textbf{unstable}.
    \item \texttt{std::stable\_sort()}: Guarantees \textbf{stability} in $O(n \log n)$ time with auxiliary memory. Implemented via \textbf{Merge Sort / TimSort}.
\end{itemize}

\newpage
\begin{center}
\vspace{0.5em}
{\color{primary}\rule{0.45\textwidth}{1.2pt}}\\[4pt]
{\Large\bfseries\color{primary} Section X: Examination Traps and Quick Revision}\\[2pt]
{\color{primary}\rule{0.45\textwidth}{1.2pt}}
\vspace{0.5em}
\end{center}

\section{Sorting Examination Traps}

\begin{enumerate}
    \item \textbf{Bubble Sort Best Case:} Only $O(n)$ when the early-termination flag is implemented. Unoptimized Bubble Sort is always $\Theta(n^2)$.
    \item \textbf{Insertion Sort Comparison Operator:} Must use \texttt{A[j] > key} for stability. Using \texttt{>=} ruins stability.
    \item \textbf{Merge Sort Memory:} Conventional array Merge Sort is \textbf{not in-place}; it requires $O(n)$ buffer storage.
    \item \textbf{Quick Sort Worst Case:} Conventional Quick Sort is $\mathbf{O(n^2)}$, not $O(n \log n)$. It occurs on already sorted input with first/last pivot.
    \item \textbf{Heap Sort Stability:} Standard Heap Sort is \textbf{unstable} due to long-distance swaps between root and leaf nodes.
    \item \textbf{Selection Sort Comparisons:} Selection Sort \textbf{always} performs $\frac{n(n-1)}{2}$ comparisons ($\Theta(n^2)$), even if the array is already sorted.
    \item \textbf{Counting Sort Scope:} Runs in $O(n + k)$, not unconditional $O(n)$. It is only linear when range $k \leq O(n)$.
    \item \textbf{Radix Sort Prerequisite:} LSD Radix Sort \textbf{strictly requires} a stable inner sub-sort to preserve prior digit orders.
    \item \textbf{Shell Sort Complexity:} Cannot be stated without specifying the gap sequence! Basic halving sequence is $O(n^2)$.
    \item \textbf{TimSort vs Introsort:} TimSort is \textbf{stable} (Python); Introsort is \textbf{unstable} (C++ \texttt{std::sort}).
\end{enumerate}

\section{Sorting Quick-Revision Sheet}

\begin{table}[htbp]
\centering
\small
\begin{tabularx}{\textwidth}{p{3.5cm} p{2cm} p{2cm} p{2cm} c X}
\toprule
\textbf{Algorithm} & \textbf{Best Time} & \textbf{Average Time} & \textbf{Worst Time} & \textbf{Stable?} & \textbf{Auxiliary Space} \\
\midrule
Bubble Sort (Opt)  & $O(n)$ & $O(n^2)$ & $O(n^2)$ & Yes & $O(1)$ \\
Insertion Sort     & $O(n)$ & $O(n^2)$ & $O(n^2)$ & Yes & $O(1)$ \\
Selection Sort     & $O(n^2)$ & $O(n^2)$ & $O(n^2)$ & No  & $O(1)$ \\
Merge Sort         & $O(n \log n)$ & $O(n \log n)$ & $O(n \log n)$ & Yes & $O(n)$ \\
Quick Sort         & $O(n \log n)$ & $O(n \log n)$ & $O(n^2)$ & No  & $O(\log n)$ \\
Heap Sort          & $O(n \log n)$ & $O(n \log n)$ & $O(n \log n)$ & No  & $O(1)$ \\
Counting Sort      & $O(n + k)$ & $O(n + k)$ & $O(n + k)$ & Yes & $O(n + k)$ \\
Radix Sort         & $O(d(n + b))$ & $O(d(n + b))$ & $O(d(n + b))$ & Yes & $O(n + b)$ \\
TimSort            & $O(n)$ & $O(n \log n)$ & $O(n \log n)$ & Yes & $O(n)$ \\
Introsort          & $O(n \log n)$ & $O(n \log n)$ & $O(n \log n)$ & No  & $O(\log n)$ \\
\bottomrule
\end{tabularx}
\caption{Rapid-fire revision summary of key sorting algorithms.}
\end{table}

\begin{quote}
\textbf{Final Advice for the Examination Hall:} For each algorithm, ensure you can seamlessly:
\begin{enumerate}
    \item \textbf{Trace:} Execute 2--3 passes by hand on a 5-element array.
    \item \textbf{Classify:} State time complexity (best, avg, worst), space complexity, stability, and in-place status.
    \item \textbf{Justify:} Explain \textit{why} the algorithm achieves that complexity (e.g., recurrence relation, tree height, or inversion elimination).
\end{enumerate}
\end{quote}


\newpage
\begin{center}
\vspace{0.5em}
{\color{primary}\rule{0.45\textwidth}{1.2pt}}\\[4pt]
{\Large\bfseries\color{primary} Section XI: Pair-Sum Algorithm with $O(n \log n)$ Complexity}\\[2pt]
{\color{primary}\rule{0.45\textwidth}{1.2pt}}
\vspace{0.5em}
\end{center}

\section{Pair-Sum Algorithm with $O(n \log n)$ Complexity}

\subsection{Problem Definition}

Given an array of $n$ integers and a target value $T$, determine whether there exist two distinct array elements whose sum equals $T$:
\[
A[i] + A[j] = T \quad (i \neq j)
\]
For example, given $A = [2, 7, 11, 15, 3, 6]$ and $T = 9$:
\begin{itemize}
    \item Valid pair: $2 + 7 = 9$
    \item Valid pair: $3 + 6 = 9$
\end{itemize}

\subsection{Method 1: Brute Force ($O(n^2)$)}
Check all distinct pairs using nested loops:
\[
\text{Total candidate pairs} = \binom{n}{2} = \frac{n(n-1)}{2} \implies O(n^2) \text{ time, } O(1) \text{ space.}
\]
While simple, this quadratic approach is unacceptable for large datasets in recruitment assessments.

\subsection{Method 2: Sorting and Two Pointers ($O(n \log n)$)}

\subsubsection*{Core Mechanism}
\begin{enumerate}
    \item Sort the array in ascending order: $O(n \log n)$.
    \item Initialize two pointers: \texttt{left} at index $0$, \texttt{right} at index $n - 1$.
    \item Calculate running sum $S = A[\text{left}] + A[\text{right}]$:
    \begin{itemize}
        \item \textbf{Case 1 ($S == T$):} Target pair found! Return $(A[\text{left}], A[\text{right}])$.
        \item \textbf{Case 2 ($S < T$):} Sum is too small. Since array is sorted, increment \texttt{left} to evaluate a larger element.
        \item \textbf{Case 3 ($S > T$):} Sum is too large. Decrement \texttt{right} to evaluate a smaller element.
    \end{itemize}
    \item Terminate when a pair is found or when $\text{left} \geq \text{right}$.
\end{enumerate}

\subsection{Step-by-Step Example}
Given $A = [8, 3, 2, 7, 11, 6]$ and $T = 9$:
\begin{enumerate}
    \item Sort the array: $[2, 3, 6, 7, 8, 11]$.
    \item Initialize: $\text{left} = 0 \; (A[0] = 2)$, $\text{right} = 5 \; (A[5] = 11)$.
\end{enumerate}

\begin{table}[htbp]
\centering
\small
\begin{tabular}{c c c c l}
\toprule
\textbf{Step} & \textbf{Left ($A[\text{left}]$)} & \textbf{Right ($A[\text{right}]$)} & \textbf{Sum ($S$)} & \textbf{Algorithmic Decision} \\
\midrule
1 & 2 & 11 & 13 & $13 > 9 \implies$ Sum too large; decrement \texttt{right} \\
2 & 2 & 8  & 10 & $10 > 9 \implies$ Sum too large; decrement \texttt{right} \\
3 & 2 & 7  & 9  & $\mathbf{9 == 9 \implies \text{Pair Found: } (2, 7)}$ \\
\bottomrule
\end{tabular}
\caption{Two-pointer convergence trace for Target $T = 9$.}
\end{table}

\subsection{C++ Implementation}

\begin{lstlisting}[style=cppstyle, caption={C++ Solution for Pair-Sum using Sorting and Two Pointers.}]
#include <iostream>
#include <vector>
#include <algorithm>
#include <utility>
using namespace std;

bool findPairSum(const vector<int>& input, long long target, pair<int, int>& result) {
    vector<int> a = input;
    sort(a.begin(), a.end()); // O(n log n)

    int left = 0;
    int right = static_cast<int>(a.size()) - 1;

    while (left < right) {
        long long sum = static_cast<long long>(a[left]) + a[right];

        if (sum == target) {
            result = {a[left], a[right]};
            return true;
        } else if (sum < target) {
            left++;
        } else {
            right--;
        }
    }
    return false;
}

int main() {
    vector<int> a = {8, 3, 2, 7, 11, 6};
    long long target = 9;
    pair<int, int> result;

    if (findPairSum(a, target, result)) {
        cout << "Pair found: (" << result.first << ", " << result.second << ")\n";
    } else {
        cout << "No valid pair exists\n";
    }
    return 0;
}
\end{lstlisting}

\textbf{Complexity Analysis:}
\begin{itemize}
    \item \textbf{Sorting Time:} $O(n \log n)$.
    \item \textbf{Two-Pointer Scan Time:} $O(n)$ (each pointer moves at most $n$ times).
    \item \textbf{Total Time Complexity:} $\boxed{O(n \log n)}$.
    \item \textbf{Auxiliary Space:} $O(1)$ if in-place sorting is allowed; $O(n)$ if sorting a cloned array.
\end{itemize}

\subsection{Important Variations}
\begin{itemize}
    \item \textbf{Return Original Array Indices:} Store elements as pairs $(\text{value}, \text{originalIndex})$. Sort by value in $O(n \log n)$, run two pointers, and retrieve the original indices.
    \item \textbf{Already Sorted Array:} Skips sorting! Two-pointer scan alone runs in $\boxed{O(n)}$ linear time.
    \item \textbf{Hash Set Alternative:} Using a hash table offers expected $O(n)$ time by looking up $T - A[i]$, but requires $O(n)$ auxiliary memory and risks worst-case $O(n^2)$ time under hash collisions.
\end{itemize}

\newpage
\begin{center}
\vspace{0.5em}
{\color{primary}\rule{0.45\textwidth}{1.2pt}}\\[4pt]
{\Large\bfseries\color{primary} Section XII: Binary Search}\\[2pt]
{\color{primary}\rule{0.45\textwidth}{1.2pt}}
\vspace{0.5em}
\end{center}

\section{Binary Search}

\textbf{Binary Search} is an optimal logarithmic searching algorithm that operates on a sorted collection by repeatedly bisecting the search interval into halves. At each step, comparing the target against the middle element allows discarding half of the remaining elements.

\subsection{The Four Essential Prerequisites}
\begin{enumerate}
    \item \textbf{Sorted Sequence:} The underlying data must be strictly sorted according to the comparison rule.
    \item \textbf{Monotonic Comparison:} Comparison logic must match the sort order (ascending vs descending).
    \item \textbf{Random Access $O(1)$:} Direct index calculation is required. Linked lists \textbf{cannot} achieve $O(\log n)$ binary search because accessing the midpoint requires traversing $O(n)$ node links.
    \item \textbf{Boundary Exclusions:} Must update boundaries as $\text{low} = \text{mid} + 1$ and $\text{high} = \text{mid} - 1$ to prevent infinite loops.
\end{enumerate}

\subsection{Algorithmic Formulation}

\begin{lstlisting}[style=pseudocode, caption={Binary Search Algorithm.}]
BINARY_SEARCH(A, target):
    low = 0
    high = length(A) - 1

    while low <= high:
        mid = low + floor((high - low) / 2) // Prevents integer overflow
        if A[mid] == target:
            return mid
        else if A[mid] < target:
            low = mid + 1
        else:
            high = mid - 1

    return -1 // Target not found
\end{lstlisting}

\begin{quote}
\textbf{Integer Overflow Avoidance:} In programming languages with fixed integer widths, writing \texttt{mid = (low + high) / 2} can cause arithmetic overflow when $\text{low} + \text{high} > 2^{31}-1$. The algebraically equivalent form:
\[
\boxed{\text{mid} = \text{low} + \left\lfloor \frac{\text{high} - \text{low}}{2} \right\rfloor}
\]
is completely immune to integer overflow.
\end{quote}

\subsection{Trace of Successful Search}
Search for target $T = 23$ in sorted array:
\[
A = [2, 5, 8, 12, 16, 23, 38, 56, 72, 91], \quad n = 10
\]

\begin{table}[htbp]
\centering
\small
\begin{tabular}{c c c c c l}
\toprule
\textbf{Iteration} & \textbf{Low} & \textbf{High} & \textbf{Mid} & \textbf{$A[\text{mid}]$} & \textbf{Decision} \\
\midrule
1 & 0 & 9 & 4 & 16 & $16 < 23 \implies \text{low} = 4 + 1 = 5$ (Search right half) \\
2 & 5 & 9 & 7 & 56 & $56 > 23 \implies \text{high} = 7 - 1 = 6$ (Search left half) \\
3 & 5 & 6 & 5 & 23 & $\mathbf{23 == 23 \implies \text{Target Found at index 5!}}$ \\
\bottomrule
\end{tabular}
\caption{Trace of binary search locating target 23 in 3 iterations.}
\end{table}

\subsection{Trace When Target is Absent}
Search for $T = 24$ in the same array:
\begin{table}[htbp]
\centering
\small
\begin{tabular}{c c c c c l}
\toprule
\textbf{Iteration} & \textbf{Low} & \textbf{High} & \textbf{Mid} & \textbf{$A[\text{mid}]$} & \textbf{Decision} \\
\midrule
1 & 0 & 9 & 4 & 16 & $16 < 24 \implies \text{low} = 5$ \\
2 & 5 & 9 & 7 & 56 & $56 > 24 \implies \text{high} = 6$ \\
3 & 5 & 6 & 5 & 23 & $23 < 24 \implies \text{low} = 6$ \\
4 & 6 & 6 & 6 & 38 & $38 > 24 \implies \text{high} = 5$ \\
\bottomrule
\end{tabular}
\caption{Trace of binary search when target 24 is absent.}
\end{table}
At iteration 5, $\text{low} = 6 > \text{high} = 5$. The loop terminates and returns $-1$.

\subsection{Binary Search with Duplicate Keys}
When array contains duplicate elements (e.g., $[2, 4, 4, 4, 7, 9]$):
\begin{itemize}
    \item \textbf{Lower Bound ($\texttt{lower\_bound}$):} First index where $A[i] \geq \text{target}$ (Index $1$).
    \item \textbf{Upper Bound ($\texttt{upper\_bound}$):} First index where $A[i] > \text{target}$ (Index $4$).
    \item \textbf{Frequency Calculation:}
    \[
    \boxed{\text{Frequency}(x) = \text{upper\_bound}(x) - \text{lower\_bound}(x) = 4 - 1 = 3}
    \]
\end{itemize}

\subsection{Complexity Summary}

\begin{table}[htbp]
\centering
\small
\begin{tabularx}{\textwidth}{l X}
\toprule
\textbf{Execution Case} & \textbf{Time / Space Complexity} \\
\midrule
Best-Case Time          & $O(1)$ (Target located at initial midpoint) \\
Average-Case Time       & $O(\log n)$ \\
Worst-Case Time         & $\boxed{O(\log n)}$ ($\lfloor \log_2 n \rfloor + 1$ comparisons) \\
Iterative Space         & $\boxed{O(1)}$ auxiliary space \\
Recursive Space         & $O(\log n)$ auxiliary stack memory \\
\bottomrule
\end{tabularx}
\caption{Complexity of Binary Search.}
\end{table}

\newpage
\begin{center}
\vspace{0.5em}
{\color{primary}\rule{0.45\textwidth}{1.2pt}}\\[4pt]
{\Large\bfseries\color{primary} Section XIII: Breadth-First Search (BFS)}\\[2pt]
{\color{primary}\rule{0.45\textwidth}{1.2pt}}
\vspace{0.5em}
\end{center}

\section{Breadth-First Search (BFS)}

\textbf{Breadth-First Search (BFS)} is an uninformed graph traversal algorithm that explores vertices in strictly increasing order of their distance (edge count) from a designated starting vertex. It operates level by level using a \textbf{FIFO Queue}.

\subsection{Algorithmic Steps}
\begin{enumerate}
    \item Initialize all vertices as unvisited.
    \item Mark start vertex $s$ as visited and enqueue it.
    \item While queue is not empty:
    \begin{itemize}
        \item Dequeue vertex $u$ and process it.
        \item For every neighbor $v$ of $u$: if $v$ is unvisited, mark $v$ as visited, record predecessor $\text{parent}[v] = u$, and enqueue $v$.
    \end{itemize}
\end{enumerate}

\begin{quote}
\textbf{Vital Implementation Rule:} Mark a vertex as \textbf{visited at the time of enqueueing}, not when dequeueing! Delaying the visited flag causes duplicate enqueues from shared predecessors, ballooning queue memory.
\end{quote}

\subsection{Step-by-Step BFS Trace}
Given graph adjacency lists:
\begin{lstlisting}[style=pseudocode]
A: B, C
B: A, D, E
C: A, F
D: B
E: B, F
F: C, E
\end{lstlisting}
Traversing starting from $A$:
\begin{itemize}
    \item \textbf{Start:} Enqueue $A \implies \text{Queue} = [A]$, $\text{Visited} = \{A\}$.
    \item \textbf{Process $A$:} Dequeue $A$, enqueue unvisited neighbors $B, C \implies \text{Queue} = [B, C]$.
    \item \textbf{Process $B$:} Dequeue $B$, enqueue unvisited neighbors $D, E \implies \text{Queue} = [C, D, E]$.
    \item \textbf{Process $C$:} Dequeue $C$, enqueue unvisited neighbor $F \implies \text{Queue} = [D, E, F]$.
    \item \textbf{Process $D$:} Dequeue $D$, no unvisited neighbors $\implies \text{Queue} = [E, F]$.
    \item \textbf{Process $E$:} Dequeue $E$, neighbors already visited $\implies \text{Queue} = [F]$.
    \item \textbf{Process $F$:} Dequeue $F$, neighbors already visited $\implies \text{Queue} = []$.
\end{itemize}
\textbf{BFS Traversal Sequence:} $\boxed{A, B, C, D, E, F}$.

\begin{table}[htbp]
\centering
\small
\begin{tabular}{c l}
\toprule
\textbf{Distance from Source $A$} & \textbf{Vertices at Level} \\
\midrule
Level 0 & $A$ \\
Level 1 & $B, C$ \\
Level 2 & $D, E, F$ \\
\bottomrule
\end{tabular}
\caption{BFS distance level decomposition.}
\end{table}

\subsection{Shortest-Path Reconstruction}
Because BFS explores by concentric edge levels, it guarantees the \textbf{shortest path} (minimum edge count) in unweighted graphs. Tracking parent pointers:
\[
\text{parent}[B] = A, \quad \text{parent}[C] = A, \quad \text{parent}[F] = C
\]
Tracing backward from $F$: $F \to C \to A$. Reversing gives:
\[
\boxed{A \longrightarrow C \longrightarrow F \quad (\text{Shortest Path of Length 2})}
\]

\subsection{C++ Implementation}

\begin{lstlisting}[style=cppstyle, caption={C++ Breadth-First Search Implementation.}]
#include <iostream>
#include <vector>
#include <queue>
using namespace std;

vector<int> bfs(const vector<vector<int>>& adj, int start) {
    int n = static_cast<int>(adj.size());
    vector<bool> visited(n, false);
    vector<int> order;
    if (start < 0 || start >= n) return order;

    queue<int> q;
    visited[start] = true;
    q.push(start);

    while (!q.empty()) {
        int u = q.front();
        q.pop();
        order.push_back(u);

        for (int v : adj[u]) {
            if (!visited[v]) {
                visited[v] = true;
                q.push(v);
            }
        }
    }
    return order;
}

int main() {
    int V = 6;
    vector<vector<int>> adj(V);
    auto addEdge = [&](int u, int v) {
        adj[u].push_back(v);
        adj[v].push_back(u);
    };

    addEdge(0, 1); addEdge(0, 2); addEdge(1, 3);
    addEdge(1, 4); addEdge(2, 5); addEdge(4, 5);

    vector<int> order = bfs(adj, 0);
    cout << "BFS Order: ";
    for (int x : order) cout << x << " ";
    cout << '\n';
    return 0;
}
\end{lstlisting}

\textbf{Complexity:} Time $\boxed{O(V + E)}$ with adjacency list; Auxiliary space $\boxed{O(V)}$ for queue and visited vector.

\newpage
\begin{center}
\vspace{0.5em}
{\color{primary}\rule{0.45\textwidth}{1.2pt}}\\[4pt]
{\Large\bfseries\color{primary} Section XIV: Depth-First Search (DFS)}\\[2pt]
{\color{primary}\rule{0.45\textwidth}{1.2pt}}
\vspace{0.5em}
\end{center}

\section{Depth-First Search (DFS)}

\textbf{Depth-First Search (DFS)} is a graph traversal algorithm that explores as deep as possible along each branch before backtracking. It operates on \textbf{LIFO} principles via program recursion or an explicit stack.

\subsection{Recursive C++ Implementation}

\begin{lstlisting}[style=cppstyle, caption={C++ Recursive Depth-First Search.}]
#include <iostream>
#include <vector>
using namespace std;

void dfsUtil(const vector<vector<int>>& adj, int u, vector<bool>& visited, vector<int>& order) {
    visited[u] = true;
    order.push_back(u);

    for (int v : adj[u]) {
        if (!visited[v]) {
            dfsUtil(adj, v, visited, order);
        }
    }
}

vector<int> dfs(const vector<vector<int>>& adj, int start) {
    int n = static_cast<int>(adj.size());
    vector<bool> visited(n, false);
    vector<int> order;
    if (start >= 0 && start < n) {
        dfsUtil(adj, start, visited, order);
    }
    return order;
}

int main() {
    int V = 6;
    vector<vector<int>> adj(V);
    auto addEdge = [&](int u, int v) {
        adj[u].push_back(v);
        adj[v].push_back(u);
    };

    addEdge(0, 1); addEdge(0, 2); addEdge(1, 3);
    addEdge(1, 4); addEdge(2, 5); addEdge(4, 5);

    vector<int> order = dfs(adj, 0);
    cout << "DFS Order: ";
    for (int x : order) cout << x << " ";
    cout << '\n';
    return 0;
}
\end{lstlisting}

\subsection{Iterative C++ Implementation}

\begin{lstlisting}[style=cppstyle, caption={C++ Iterative DFS using Explicit Stack.}]
vector<int> dfsIterative(const vector<vector<int>>& adj, int start) {
    int n = static_cast<int>(adj.size());
    vector<bool> visited(n, false);
    vector<int> order;
    if (start < 0 || start >= n) return order;

    vector<int> st;
    st.push_back(start);

    while (!st.empty()) {
        int u = st.back();
        st.pop_back();

        if (visited[u]) continue;
        visited[u] = true;
        order.push_back(u);

        // Push neighbors in reverse to explore smaller indices first
        for (int i = static_cast<int>(adj[u].size()) - 1; i >= 0; --i) {
            int v = adj[u][i];
            if (!visited[v]) st.push_back(v);
        }
    }
    return order;
}
\end{lstlisting}

\subsection{Applications of DFS}
\begin{itemize}
    \item \textbf{Cycle Detection:} Back-edges in directed graphs indicate cycles.
    \item \textbf{Topological Sorting:} Order vertices of a DAG by finishing times.
    \item \textbf{Strongly Connected Components:} Tarjan's and Kosaraju's algorithms.
    \item \textbf{Maze Solving and Backtracking:} Path existence and Sudoku solvers.
\end{itemize}

\newpage
\begin{center}
\vspace{0.5em}
{\color{primary}\rule{0.45\textwidth}{1.2pt}}\\[4pt]
{\Large\bfseries\color{primary} Section XV: BFS vs DFS --- Comparative Analysis}\\[2pt]
{\color{primary}\rule{0.45\textwidth}{1.2pt}}
\vspace{0.5em}
\end{center}

\section{BFS vs DFS --- Comparative Analysis}

\subsection{Comprehensive Comparison Table}

\begin{table}[htbp]
\centering
\small
\begin{tabularx}{\textwidth}{p{3.5cm} X X}
\toprule
\textbf{Basis of Comparison} & \textbf{Breadth-First Search (BFS)} & \textbf{Depth-First Search (DFS)} \\
\midrule
Full Form & Breadth-First Search & Depth-First Search \\
Core Traversal Mechanism & Explores frontier level by level & Dives deeply along paths; backtracks \\
Primary Data Structure & \textbf{Queue} (FIFO) & \textbf{Stack} (LIFO) or Call Stack \\
Shortest Path Guarantee & \textbf{Yes} (in unweighted graphs) & \textbf{No} (may find longer path first) \\
Time (Adjacency List) & $\boxed{O(V + E)}$ & $\boxed{O(V + E)}$ \\
Time (Adjacency Matrix) & $O(V^2)$ & $O(V^2)$ \\
Auxiliary Space & $O(V)$ (stores broad level frontier) & $O(V)$ (stores recursive call stack) \\
Cycle Detection & Requires color / cross-edge tracking & Natural via back-edge detection \\
Topological Sort & Kahn's in-degree algorithm & Reverse post-order traversal \\
Suitability & Shortest paths, peer-to-peer networks & Decision trees, mazes, reachability \\
\bottomrule
\end{tabularx}
\caption{Core architectural comparison between BFS and DFS.}
\end{table}

\subsection{Why DFS Does Not Guarantee Shortest Path}
Consider reaching node $F$ from $A$ in our example graph:
\begin{itemize}
    \item \textbf{Path 1:} $A \to C \to F$ (Length $= 2$ edges).
    \item \textbf{Path 2:} $A \to B \to E \to F$ (Length $= 3$ edges).
\end{itemize}
BFS explores level 1 ($B, C$), then level 2 ($F$), discovering $F$ along the 2-edge path. DFS dives into $B$, then $D$, backtracks to $E$, and discovers $F$ along the 3-edge path first! Hence, DFS \textbf{cannot be relied on} for shortest paths.

\subsection{Traversing Disconnected Graphs}
If a graph contains multiple disconnected components, a single BFS or DFS from $s$ visits only $s$'s component. To traverse the entire forest:
\begin{lstlisting}[style=pseudocode]
for each vertex v in V:
    if not visited[v]:
        BFS(v)  // or DFS(v)
\end{lstlisting}
The overall complexity remains strictly $\boxed{O(V + E)}$.

\newpage
\begin{center}
\vspace{0.5em}
{\color{primary}\rule{0.45\textwidth}{1.2pt}}\\[4pt]
{\Large\bfseries\color{primary} Section XVI: A* Search Algorithm}\\[2pt]
{\color{primary}\rule{0.45\textwidth}{1.2pt}}
\vspace{0.5em}
\end{center}

\section{A* Search Algorithm}

\textbf{A* (A-Star)} is an informed (heuristic-guided) search algorithm that computes the least-cost path from a start node to a goal node. Unlike uninformed searches (BFS, DFS, Dijkstra), A* uses domain knowledge via a \textbf{heuristic function} to direct exploration toward the goal.

\subsection{The Evaluation Function}
A* selects the next node to expand from a min-priority queue based on the evaluation function:
\[
\boxed{f(n) = g(n) + h(n)}
\]
where:
\begin{itemize}
    \item $g(n)$: The exact cost incurred from the start node to node $n$.
    \item $h(n)$: The heuristic estimate of the remaining cost from node $n$ to the goal.
    \item $f(n)$: The estimated total cost of the cheapest solution passing through node $n$.
\end{itemize}

\begin{table}[htbp]
\centering
\small
\begin{tabular}{c c c c l}
\toprule
\textbf{Candidate Node} & \textbf{Cost Incurred $g(n)$} & \textbf{Heuristic $h(n)$} & \textbf{Total Score $f(n)$} & \textbf{Priority Selection} \\
\midrule
Node $A$ & 3 & 8 & 11 & Secondary candidate \\
Node $B$ & 5 & 2 & \textbf{7} & \textbf{Selected first ($f=7 < 11$)} \\
\bottomrule
\end{tabular}
\caption{A* decision: Node $B$ is prioritized despite higher $g(n)$ because its remaining cost is tiny.}
\end{table}

\subsection{Heuristic Properties for Optimality}

\subsubsection{1. Admissible Heuristic (Never Overestimates)}
A heuristic $h(n)$ is \textbf{admissible} if it never overestimates the actual true minimum cost $h^*(n)$ to reach the goal:
\[
\boxed{0 \leq h(n) \leq h^*(n), \quad \text{with } h(\text{goal}) = 0}
\]
\textbf{Guaranteed Optimality:} If $h(n)$ is admissible, A* tree search is guaranteed to return an \textbf{optimal (lowest cost) path}.

\subsubsection{2. Consistent (Monotone) Heuristic}
A heuristic is \textbf{consistent} if for every node $u$ and successor $v$ connected by edge of cost $c(u, v)$:
\[
\boxed{h(u) \leq c(u, v) + h(v)}
\]
This satisfies the triangle inequality. With a consistent heuristic, the $f$-values along any path are non-decreasing, and A* graph search \textbf{never needs to reopen an already closed node}.

\subsection{Why A* Can Be Faster Than BFS (Worked Grid Example)}
Consider navigating an unweighted grid from $(0, 0)$ to $(4, 0)$ with unit movement costs:
\begin{itemize}
    \item \textbf{BFS Approach:} Expands blindly in circular concentric Manhattan diamonds. Before reaching distance 4, BFS expands:
    \[
    \text{Level 0 (1)} + \text{Level 1 (4)} + \text{Level 2 (8)} + \text{Level 3 (12)} = \boxed{25 \text{ expanded states!}}
    \]
    \item \textbf{A* Approach (Manhattan Heuristic $h(x, y) = |4-x| + |y|$):}
    For points along the direct axis:
    \[
    (0,0) \to (1,0) \to (2,0) \to (3,0) \to (4,0) \implies \text{all have } f = g + h = 4!
    \]
    Any point off the direct axis (e.g., $(0, 1)$) has $g=1, h=5 \implies f=6$. A* expands \textbf{only 5 states}, avoiding 20 useless expansions!
\end{itemize}

\subsection{Comparison: BFS vs A*}

\begin{table}[htbp]
\centering
\small
\begin{tabularx}{\textwidth}{p{3.5cm} X X}
\toprule
\textbf{Feature} & \textbf{Breadth-First Search (BFS)} & \textbf{A* Search Algorithm} \\
\midrule
Search Strategy & Uninformed (level-by-level) & Informed (heuristic guided) \\
Data Structure & FIFO Queue & Min-Priority Queue (ordered by $f(n)$) \\
Edge Weights & Unweighted graphs only & Supports arbitrary non-negative weights \\
Heuristic Knowledge & None ($h=0$) & Active heuristic $h(n)$ \\
Shortest Path on Weighted & Fails & \textbf{Guaranteed optimal} (if $h$ is admissible) \\
Special Case $h(n) = 0$ & Equivalent to BFS on unit edges & Becomes \textbf{Dijkstra's / Uniform-Cost Search} \\
\bottomrule
\end{tabularx}
\caption{Key distinctions between BFS and A* Search.}
\end{table}

\newpage
\begin{center}
\vspace{0.5em}
{\color{primary}\rule{0.45\textwidth}{1.2pt}}\\[4pt]
{\Large\bfseries\color{primary} Section XVII: Complete C++ Implementation of A* Search}\\[2pt]
{\color{primary}\rule{0.45\textwidth}{1.2pt}}
\vspace{0.5em}
\end{center}

\section{Complete C++ Implementation of A* Search}

\begin{lstlisting}[style=cppstyle, caption={C++ A* Search Algorithm on Weighted Graph.}]
#include <iostream>
#include <vector>
#include <queue>
#include <limits>
#include <algorithm>
using namespace std;

struct Edge {
    int to;
    long long cost;
};

struct State {
    long long f, g;
    int node;
};

struct CompareState {
    bool operator()(const State& a, const State& b) const {
        return a.f > b.f; // Min-heap: lowest f-score popped first
    }
};

struct AStarResult {
    bool found = false;
    long long cost = 0;
    vector<int> path;
};

AStarResult aStar(const vector<vector<Edge>>& graph, const vector<long long>& h, int start, int goal) {
    int n = static_cast<int>(graph.size());
    const long long INF = numeric_limits<long long>::max() / 4;

    vector<long long> gScore(n, INF);
    vector<int> parent(n, -1);
    priority_queue<State, vector<State>, CompareState> open;

    gScore[start] = 0;
    open.push({h[start], 0, start});

    while (!open.empty()) {
        State cur = open.top();
        open.pop();
        int u = cur.node;

        if (cur.g != gScore[u]) continue; // Stale queue entry
        if (u == goal) break;             // Goal reached

        for (const Edge& edge : graph[u]) {
            int v = edge.to;
            long long tentativeG = gScore[u] + edge.cost;

            if (tentativeG < gScore[v]) {
                gScore[v] = tentativeG;
                parent[v] = u;
                open.push({tentativeG + h[v], tentativeG, v});
            }
        }
    }

    if (gScore[goal] == INF) return {};

    AStarResult res;
    res.found = true;
    res.cost = gScore[goal];
    for (int v = goal; v != -1; v = parent[v]) res.path.push_back(v);
    reverse(res.path.begin(), res.path.end());
    return res;
}

int main() {
    // Vertices: 0=S, 1=A, 2=B, 3=C, 4=D, 5=G
    int V = 6;
    vector<vector<Edge>> graph(V);
    auto addEdge = [&](int u, int v, long long w) { graph[u].push_back({v, w}); };

    addEdge(0, 1, 1); addEdge(0, 2, 2); addEdge(0, 3, 1);
    addEdge(1, 5, 9); addEdge(2, 5, 2); addEdge(3, 4, 1); addEdge(4, 5, 100);

    vector<long long> h = {4, 9, 2, 100, 100, 0}; // Admissible heuristic
    AStarResult res = aStar(graph, h, 0, 5);

    if (res.found) {
        cout << "Path: ";
        for (size_t i = 0; i < res.path.size(); ++i) {
            cout << res.path[i] << (i + 1 < res.path.size() ? " -> " : "\n");
        }
        cout << "Total Cost: " << res.cost << '\n';
    }
    return 0;
}
\end{lstlisting}

\begin{lstlisting}[style=pseudocode, caption={Execution Output}]
Path: 0 -> 2 -> 5
Total Cost: 4
\end{lstlisting}

\subsection{Complexity and Practical Limitations}
\begin{itemize}
    \item \textbf{Time Complexity:} With a consistent heuristic and binary heap, runs in $O((V + E) \log V)$ time. In general implicit search spaces, worst-case time is exponential $O(b^d)$.
    \item \textbf{Space Complexity:} Memory is A*'s greatest operational bottleneck: it stores all generated open and closed nodes in RAM ($O(V)$ on explicit graphs, $O(b^d)$ on implicit trees).
\end{itemize}

\newpage
\begin{center}
\vspace{0.5em}
{\color{primary}\rule{0.45\textwidth}{1.2pt}}\\[4pt]
{\Large\bfseries\color{primary} Section XVIII: Common Recruitment Examination Questions}\\[2pt]
{\color{primary}\rule{0.45\textwidth}{1.2pt}}
\vspace{0.5em}
\end{center}

\section{Common Recruitment Examination Questions}

\subsection{Pair-Sum Questions}
\begin{quote}
\textbf{Q1: How can pair-sum be solved in $O(n \log n)$ time?}\\
\textbf{Answer:} Sort the array in $O(n \log n)$ time, then use a two-pointer scan (\texttt{left} from front, \texttt{right} from rear) in $O(n)$ time.

\textbf{Q2: Why does the two-pointer method work on a sorted array?}\\
\textbf{Answer:} Because elements are monotonic: if the sum is too small, advancing \texttt{left} guarantees an increased candidate sum; if too large, decrementing \texttt{right} guarantees a decreased candidate sum.

\textbf{Q3: What is the complexity if the array is already sorted?}\\
\textbf{Answer:} Exactly $O(n)$ time and $O(1)$ auxiliary space.
\end{quote}

\subsection{Binary Search Questions}
\begin{quote}
\textbf{Q4: What is the essential prerequisite for binary search?}\\
\textbf{Answer:} The collection must be sorted according to the search comparison key, and elements must support $O(1)$ random access.

\textbf{Q5: What is the worst-case complexity of binary search?}\\
\textbf{Answer:} $O(\log n)$ time, performing at most $\lfloor \log_2 n \rfloor + 1$ comparisons.

\textbf{Q6: Why is \texttt{low + (high - low)/2} preferred to \texttt{(low + high)/2}?}\\
\textbf{Answer:} To prevent integer overflow when $\text{low} + \text{high}$ exceeds the maximum signed 32-bit integer limit ($2^{31}-1$).

\textbf{Q7: Does binary search reliably find the first duplicate?}\\
\textbf{Answer:} No; standard binary search can stop at any matching midpoint. A lower-bound search ($\texttt{lower\_bound}$) is required to find the first occurrence.
\end{quote}

\subsection{BFS and DFS Questions}
\begin{quote}
\textbf{Q8: Which algorithm finds shortest paths in an unweighted graph?}\\
\textbf{Answer:} BFS (Breadth-First Search).

\textbf{Q9: What underlying data structures are used by BFS and DFS?}\\
\textbf{Answer:} BFS uses a FIFO Queue; DFS uses a LIFO Stack (or system call stack).

\textbf{Q10: What are the traversal complexities using an adjacency list?}\\
\textbf{Answer:} Both BFS and DFS take $O(V + E)$ time and $O(V)$ auxiliary space.

\textbf{Q11: Does DFS always guarantee shortest paths?}\\
\textbf{Answer:} No; DFS dives deeply along arbitrary branches and can discover a target along a long path before exploring a shorter alternative.
\end{quote}

\subsection{A* Search Questions}
\begin{quote}
\textbf{Q12: What is the A* evaluation formula?}\\
\textbf{Answer:} $f(n) = g(n) + h(n)$, where $g(n)$ is cost incurred from start, and $h(n)$ is estimated remaining cost to goal.

\textbf{Q13: What is an admissible heuristic?}\\
\textbf{Answer:} A heuristic that never overestimates the true minimum remaining cost ($0 \leq h(n) \leq h^*(n)$).

\textbf{Q14: When does A* become Dijkstra's algorithm?}\\
\textbf{Answer:} When $h(n) = 0$ for all nodes, $f(n) = g(n)$, reducing A* to Uniform-Cost Search / Dijkstra's algorithm.

\textbf{Q15: Why can A* outperform BFS on pathfinding?}\\
\textbf{Answer:} The heuristic directs the search towards the goal, pruning large swaths of the search space that lead away from the objective.
\end{quote}

\section{Search and Traversal Master Quick-Revision Sheet}

\begin{table}[htbp]
\centering
\small
\begin{tabularx}{\textwidth}{p{3.2cm} p{2.2cm} p{2.2cm} p{2.2cm} X}
\toprule
\textbf{Algorithm} & \textbf{Data Structure} & \textbf{Time Cost} & \textbf{Aux Space} & \textbf{Core Optimal Use Case} \\
\midrule
Pair-Sum & Two Pointers & $O(n \log n)$ & $O(1)$ & Finding pairs with target sum in sorted array \\
Binary Search & Random Array & $O(\log n)$ & $O(1)$ & Instant logarithmic key lookup in sorted list \\
BFS & FIFO Queue & $O(V + E)$ & $O(V)$ & Shortest path (edge count) in unweighted graph \\
DFS & LIFO Stack & $O(V + E)$ & $O(V)$ & Cycle detection, topological sort, backtracking \\
A* Search & Min-PQ ($g+h$) & $O((V+E)\log V)$ & $O(V)$ & Least-cost pathfinding guided by heuristic \\
\bottomrule
\end{tabularx}
\caption{Master quick-revision comparison of search and traversal algorithms.}
\end{table}

\begin{quote}
\textbf{Final Algorithm Revision Axiom:} In the examination, always verify three invariants:
\begin{enumerate}
    \item \textbf{Preconditions:} Check if the input is sorted (required for Binary Search and Two Pointers).
    \item \textbf{Weights:} For shortest paths, verify if the graph is unweighted (use BFS) or weighted non-negative (use Dijkstra or A*).
    \item \textbf{Heuristics:} For A*, ensure the heuristic is admissible ($h \leq h^*$) to guarantee optimality.
\end{enumerate}
\end{quote}


\begin{center}
\vspace{1em}
{\color{primary}\rule{0.5\textwidth}{1.5pt}}\\[6pt]
{\Large\bfseries\color{primary} Part II: Greedy Algorithms, Recursion, and Dynamic Programming}\\[3pt]
{\color{primary}\rule{0.5\textwidth}{1.5pt}}
\vspace{1em}
\end{center}

These three topics are fundamental to algorithm design and frequently appear in technical recruitment examinations. They are especially important because many questions test not merely whether an algorithm can be implemented, but \textbf{why it works, when it is applicable, and how its time and space complexities are derived}.

This chapter provides rigorous, exhaustive coverage of definitions, design principles, step-by-step execution traces, pseudocode, production-ready C++ implementations, mathematical complexity analysis, correctness proofs, and common examination pitfalls.


\begin{center}
\vspace{0.5em}
{\color{primary}\rule{0.45\textwidth}{1.2pt}}\\[4pt]
{\Large\bfseries\color{primary} Section I: Greedy Algorithms}\\[2pt]
{\color{primary}\rule{0.45\textwidth}{1.2pt}}
\vspace{0.5em}
\end{center}


\section{What Is a Greedy Algorithm?}

A \textbf{greedy algorithm} constructs a solution step by step by selecting the option that appears best at the current stage.

Once a choice is made, the algorithm generally does not reconsider or undo that choice. The objective is to make locally optimal choices that lead to a globally optimal solution.

A greedy algorithm is typically appropriate when a problem has the \textbf{greedy-choice property} and \textbf{optimal substructure}.

\subsubsection{Real-life example}

Suppose you need to schedule as many non-overlapping meetings as possible in a single meeting room.

A greedy strategy is to select the meeting that finishes earliest, then select the next meeting that starts after the previous selection finishes.

This strategy maximizes the number of meetings under the standard activity-selection assumptions.

The important point is that choosing the meeting that finishes earliest leaves as much remaining time as possible for subsequent meetings.

However, not every locally optimal choice leads to a globally optimal solution. The correctness of a greedy strategy must be established for the particular problem.

\subsection{1.1 Main characteristics}

Greedy algorithms generally have the following characteristics:

\begin{enumerate}
    \item \textbf{Local choice:} Select the currently most promising option according to a specified rule.
    \item \textbf{Irrevocable decisions:} Once a choice is made, the algorithm typically retains it.
    \item \textbf{Feasibility:} Choices must respect the constraints of the problem.
    \item \textbf{Optimal substructure:} An optimal solution can be formed using optimal solutions to appropriate smaller subproblems.
    \item \textbf{Greedy-choice property:} There exists an optimal solution that includes the locally selected choice.
\end{enumerate}

The last property is particularly important. Optimal substructure alone is not sufficient to establish that a greedy algorithm works.

\subsection{1.2 Greedy-choice property versus optimal substructure}

These concepts are often confused in examinations.

\textbf{Greedy-choice property:} A locally optimal choice can be included in at least one globally optimal solution without preventing optimal completion.

\textbf{Optimal substructure:} An optimal solution to a problem can be constructed from optimal solutions to smaller, appropriately defined subproblems.

For example, both the Fractional Knapsack problem and the 0/1 Knapsack problem have useful optimal-substructure formulations. However, the ratio-based greedy algorithm is optimal for Fractional Knapsack and is not generally optimal for 0/1 Knapsack.

This distinction is a common examination question.

\subsection{1.3 General greedy algorithm template}

\begin{lstlisting}[style=pseudocode]
GREEDY(problem):
    solution = empty

    while solution is incomplete:
        candidate = select_best_available_candidate()

        if candidate keeps solution feasible:
            add candidate to solution

    return solution
\end{lstlisting}

This is a conceptual template, not a complete algorithm. The selection rule, feasibility condition, and stopping condition must be defined for each problem.

\subsection{1.4 How to prove a greedy algorithm is correct}

Several proof techniques are frequently used.

\subsubsection{Exchange argument}

Show that any optimal solution can be modified to include the greedy choice without making the solution worse.

This is commonly used to prove Activity Selection and Minimum Spanning Tree algorithms.

\subsubsection{Stays-ahead argument}

Show that after each step, the greedy solution is at least as good as a competing solution according to a carefully defined measure.

This technique is useful when comparing progress across a sequence of choices.

\subsubsection{Cut property}

For a weighted, undirected graph, a minimum-weight edge crossing a suitable cut is safe for constructing a Minimum Spanning Tree when the cut respects the edges already chosen. This is a core correctness principle for Prim's and Kruskal's algorithms.

\textbf{Exam principle:} To justify a greedy algorithm, prove more than the fact that it makes the smallest or largest choice. Establish why that choice cannot prevent an optimal complete solution.

\section{Activity Selection Problem}

\subsection{2.1 Problem definition}

Given a set of activities, each with a start time $s_i$ and finish time $f_i$, select the maximum number of activities that do not overlap.

Assume that an activity beginning exactly when another ends is permitted. Thus, an activity is compatible with the last selected activity if:


\[
s_i\geq f_{\text{last}}
\]


Each activity occupies one resource, such as a meeting room, and cannot overlap another selected activity.

\subsection{2.2 Greedy strategy}

Select the available activity with the \textbf{earliest finish time}.

Then repeatedly select the earliest-finishing activity whose start time is at or after the finish time of the last selected activity.

The algorithm is:

\begin{enumerate}
    \item Sort activities by finish time in ascending order.
    \item Select the first activity.
    \item Consider the remaining activities in sorted order.
    \item Select an activity if it does not overlap the last selected activity.
    \item Continue until every activity has been considered.
\end{enumerate}

\subsection{2.3 Step-by-step example}

Consider these activities:

\begin{table}[htbp]
\centering
\caption{Activity Selection Input Dataset}
\label{tab:activity_selection_data}
\begin{tabular}{ccc}
\toprule
\textbf{Activity} & \textbf{Start Time ($s_i$)} & \textbf{Finish Time ($f_i$)} \\
\midrule
A & 1 & 4 \\
B & 3 & 5 \\
C & 0 & 6 \\
D & 5 & 7 \\
E & 3 & 9 \\
F & 5 & 9 \\
G & 6 & 10 \\
H & 8 & 11 \\
I & 8 & 12 \\
J & 2 & 14 \\
K & 12 & 16 \\
\bottomrule
\end{tabular}
\end{table}

The activities are already ordered by finish time.

\textbf{Step 1:} Select A, which finishes at time 4.

\textbf{Step 2:} Consider B. It starts at 3, which is before 4, so reject it.

\textbf{Step 3:} Consider C. It starts at 0, so reject it.

\textbf{Step 4:} Consider D. It starts at 5, which is at or after 4, so select D. Its finish time is 7.

\textbf{Step 5:} Consider E, F, and G. Their start times are before 7, so reject them.

\textbf{Step 6:} Consider H. It starts at 8, so select H. Its finish time is 11.

\textbf{Step 7:} Reject I because it starts at 8, before 11.

\textbf{Step 8:} Reject J because it starts at 2.

\textbf{Step 9:} Select K because it starts at 12, after H finishes.

The selected activities are:


\[
\boxed{A,\ D,\ H,\ K}
\]


Four compatible activities have been selected.

\subsection{2.4 C++ implementation}

\begin{lstlisting}[style=cppstyle]
#include <algorithm>
#include <iostream>
#include <vector>
using namespace std;

struct Activity {
    int start;
    int finish;
    int id;
};

vector<Activity> activitySelection(vector<Activity> activities) {
    sort(activities.begin(), activities.end(),
         [](const Activity& a, const Activity& b) {
             if (a.finish != b.finish) {
                 return a.finish < b.finish;
             }
             return a.start < b.start;
         });

    vector<Activity> selected;

    for (const Activity& activity : activities) {
        if (selected.empty() ||
            activity.start >= selected.back().finish) {
            selected.push_back(activity);
        }
    }

    return selected;
}

int main() {
    vector<Activity> activities = {
        {1, 4, 1}, {3, 5, 2}, {0, 6, 3},
        {5, 7, 4}, {3, 9, 5}, {5, 9, 6},
        {6, 10, 7}, {8, 11, 8}, {8, 12, 9},
        {2, 14, 10}, {12, 16, 11}
    };

    vector<Activity> selected = activitySelection(activities);

    cout << "Selected activities:\n";

    for (const Activity& activity : selected) {
        cout << "Activity " << activity.id
             << ": " << activity.start
             << "-" << activity.finish << '\n';
    }

    cout << "Total: " << selected.size() << '\n';

    return 0;
}
\end{lstlisting}

\subsubsection{Complexity}

Sorting takes $O(n\log n)$, and scanning the activities takes $O(n)$.

Therefore, the total time complexity is:


\[
\boxed{O(n\log n)}
\]


If the activities are already sorted by finish time, the selection phase takes $O(n)$.

\subsection{2.5 Why does the greedy strategy work?}

The earliest-finishing compatible activity leaves the largest possible amount of time for subsequent activities.

Any optimal solution that begins with a different activity can replace that first activity with the earliest-finishing one. The replacement cannot make any subsequent selected activity incompatible, because the greedy activity finishes no later than the replaced activity.

This exchange argument establishes that there is an optimal solution beginning with the greedy choice. The remaining problem is an activity-selection problem on the time after the chosen activity finishes, so the method can be applied repeatedly.

\textbf{Important variation:} If activities have different profits or weights and the goal is to maximize total value rather than the number of activities, this simple greedy rule is not generally optimal. Weighted Interval Scheduling is commonly solved using Dynamic Programming.

\section{Fractional Knapsack Problem}

\subsection{3.1 Problem definition}

Given a knapsack with capacity $W$, each item has a weight $w_i$ and value $v_i$. The objective is to maximize the total value without exceeding capacity.

In the \textbf{Fractional Knapsack problem}, an item may be divided into fractions. If an item weighs 10 kg and has value 100, taking half the item gives weight 5 kg and value 50.

This divisibility property makes a greedy strategy optimal.

\subsection{3.2 Greedy strategy}

Calculate each item's value per unit of weight:


\[
\boxed{r_i=\frac{v_i}{w_i}}
\]


Sort items by this ratio in descending order.

Then:

\begin{enumerate}
    \item Take as much as possible of the item with the highest ratio.
    \item If the remaining capacity is too small for the whole item, take the fraction that fits.
    \item Continue with the next highest ratio until the knapsack is full or all useful items are exhausted.
\end{enumerate}

\subsection{3.3 Worked example}

Knapsack capacity:


\[
W=50
\]


Items:

\begin{table}[htbp]
\centering
\caption{Fractional Knapsack Item Set ($W = 50$)}
\label{tab:fractional_knapsack_data}
\begin{tabular}{cccc}
\toprule
\textbf{Item} & \textbf{Weight ($w_i$)} & \textbf{Value ($v_i$)} & \textbf{Value per Unit Weight ($r_i = v_i/w_i$)} \\
\midrule
A & 10 & 60 & 6.0 \\
B & 20 & 100 & 5.0 \\
C & 30 & 120 & 4.0 \\
\bottomrule
\end{tabular}
\end{table}

The ratio order is A, B, C.

\textbf{Take item A fully.}

Weight used: 10.

Value: 60.

Capacity remaining: 40.

\textbf{Take item B fully.}

Weight used: 30.

Total value: 160.

Capacity remaining: 20.

\textbf{Take two-thirds of item C.}

The remaining capacity is 20, while C weighs 30.

Fraction taken:


\[
\frac{20}{30}=\frac23
\]


Value obtained from C:


\[
120\times\frac23=80
\]


Total value:


\[
60+100+80=\boxed{240}
\]


The optimal value is 240.

\subsection{3.4 C++ implementation}

\begin{lstlisting}[style=cppstyle]
#include <algorithm>
#include <iostream>
#include <vector>
using namespace std;

struct Item {
    double weight;
    double value;
};

double fractionalKnapsack(vector<Item> items, double capacity) {
    sort(items.begin(), items.end(),
         [](const Item& a, const Item& b) {
             return a.value / a.weight > b.value / b.weight;
         });

    double totalValue = 0.0;

    for (const Item& item : items) {
        if (capacity <= 0) {
            break;
        }

        double takenWeight = min(capacity, item.weight);

        totalValue += takenWeight * item.value / item.weight;
        capacity -= takenWeight;
    }

    return totalValue;
}

int main() {
    vector<Item> items = {
        {10, 60}, {20, 100}, {30, 120}
    };

    cout << fractionalKnapsack(items, 50) << '\n';

    return 0;
}
\end{lstlisting}

Output:

\begin{lstlisting}[style=pseudocode]
240
\end{lstlisting}

This implementation assumes strictly positive item weights. Items with zero weight require separate treatment because their value-to-weight ratio is undefined.

\subsubsection{Complexity}

Sorting items takes $O(n\log n)$; scanning them takes $O(n)$.


\[
\boxed{O(n\log n)}
\]


\subsection{3.5 Why does the greedy strategy work?}

The item with the highest value per unit of weight provides the most value for every unit of capacity consumed.

If a proposed solution uses some capacity on an item with a lower ratio while an item with a higher ratio remains available, replacing some weight from the lower-ratio item with the higher-ratio item does not decrease the total value.

Repeated exchanges produce a solution that fills capacity in descending ratio order.

The argument depends on being able to take fractions of items.

\section{Why Greedy Fails for the 0/1 Knapsack Problem}

In the \textbf{0/1 Knapsack problem}, each item must be taken completely or excluded completely. Fractions are not allowed.

The problem is often confused with Fractional Knapsack, but the difference changes the optimal algorithm.

\subsection{4.1 Counterexample}

Suppose the capacity is:


\[
W=50
\]


The items are:

\begin{table}[htbp]
\centering
\caption{0/1 Knapsack Counterexample Items ($W = 50$)}
\label{tab:knapsack_counterexample}
\begin{tabular}{cccc}
\toprule
\textbf{Item} & \textbf{Weight ($w_i$)} & \textbf{Value ($v_i$)} & \textbf{Value-to-Weight Ratio} \\
\midrule
A & 10 & 60 & 6.0 \\
B & 20 & 100 & 5.0 \\
C & 30 & 120 & 4.0 \\
\bottomrule
\end{tabular}
\end{table}

The ratio-based greedy method selects A and B:

\begin{itemize}
    \item Weight: $10+20=30$.
    \item Value: $60+100=160$.
    \item Remaining capacity: 20, which is insufficient for C.
\end{itemize}

However, selecting B and C gives:

\begin{itemize}
    \item Weight: $20+30=50$.
    \item Value: $100+120=220$.
\end{itemize}

Therefore:


\[
\boxed{\text{Optimal 0/1 value}=220}
\]


The greedy ratio approach obtains only 160, so it is not optimal.

\subsubsection{Exam conclusion}

\begin{itemize}
    \item Fractional Knapsack: greedy by value-to-weight ratio is optimal.
    \item 0/1 Knapsack: ratio-based greedy is not generally optimal.
    \item 0/1 Knapsack: Dynamic Programming is a standard solution for integer capacities.
\end{itemize}

The 0/1 version can also be solved using branch-and-bound, backtracking, or other approaches, but a simple ratio-greedy strategy does not guarantee the optimum.

\section{Coin Change Problem}

\subsection{5.1 Problem definition}

Given a set of coin denominations and a target amount, determine the coins needed to produce that amount.

One common version asks for the minimum number of coins.

The greedy approach repeatedly chooses the largest coin that does not exceed the remaining amount.

However, \textbf{greedy coin selection is not optimal for arbitrary denomination systems}.

\subsection{5.2 Example where greedy works}

Suppose the available denominations are:


\[
\\{1,5,10,25\\}
\]


and the target amount is 63.

A greedy solution chooses:

\begin{itemize}
    \item Two 25-unit coins: remaining amount 13.
    \item One 10-unit coin: remaining amount 3.
    \item Three 1-unit coins: remaining amount 0.
\end{itemize}

Total:


\[
2+1+3=\boxed{6\text{ coins}}
\]


For this denomination system, the greedy result is optimal for the example.

\subsection{5.3 Counterexample where greedy fails}

Consider denominations:


\[
\\{1,3,4\\}
\]


and target:


\[
T=6
\]


The greedy method selects:

\begin{itemize}
    \item One 4-unit coin: remaining 2.
    \item Two 1-unit coins.
\end{itemize}

Total:


\[
\boxed{3\text{ coins}}
\]


But the optimal solution is:


\[
3+3=6
\]


which uses only:


\[
\boxed{2\text{ coins}}
\]


Therefore, greedy selection is not optimal for this denomination system.

\subsection{5.4 Complexity}

If the denominations are sorted and there are $m$ coin types, the basic greedy method can run in $O(m)$ time after sorting, or $O(m\log m)$ if sorting is required.

The algorithm is simple and efficient, but its correctness depends on the denomination system.

\textbf{Exam distinction:} A greedy algorithm can be fast and still be wrong if the problem lacks the greedy-choice property.

\section{Minimum Spanning Tree: Kruskal's Algorithm}

\subsection{6.1 What is a Minimum Spanning Tree?}

Given a connected, weighted, undirected graph, a \textbf{Minimum Spanning Tree (MST)} is a spanning tree with the minimum possible total edge weight.

A spanning tree connects every vertex, contains no cycles, and has exactly $V-1$ edges.

For example, a network engineer might need to connect several offices with the minimum total cable length while ensuring every office remains connected.

Two classic greedy algorithms solve the MST problem:

\begin{itemize}
    \item Kruskal's algorithm.
    \item Prim's algorithm.
\end{itemize}

\subsection{6.2 Kruskal's greedy strategy}

Kruskal's algorithm processes edges in ascending order of weight, adding an edge if it does not create a cycle.

Algorithm:

\begin{enumerate}
    \item Sort edges by increasing weight.
    \item Initialize the MST with no edges.
    \item Examine edges in sorted order.
    \item Add an edge if its endpoints are in different connected components.
    \item Stop after selecting $V-1$ edges.
\end{enumerate}

A Disjoint Set Union (DSU), also called Union-Find, efficiently tracks connected components.

\subsection{6.3 Worked example}

Consider this graph:

\begin{table}[htbp]
\centering
\caption{Graph Edges and Weights for Kruskal's MST}
\label{tab:kruskal_edges}
\begin{tabular}{cc}
\toprule
\textbf{Edge} & \textbf{Weight} \\
\midrule
A--B & 1 \\
B--C & 2 \\
A--C & 3 \\
B--D & 4 \\
C--D & 5 \\
\bottomrule
\end{tabular}
\end{table}

Sort the edges:

\texttt{A--B (1), B--C (2), A--C (3), B--D (4), C--D (5)}

\textbf{Step 1:} Select A--B, weight 1.

The MST has one edge.

\textbf{Step 2:} Select B--C, weight 2.

The MST now connects A, B, and C.

\textbf{Step 3:} Examine A--C, weight 3.

A and C are already connected through A--B--C. Adding A--C would create a cycle, so reject it.

\textbf{Step 4:} Select B--D, weight 4.

All four vertices are connected, and the MST has $V-1=3$ edges.

Total weight:


\[
1+2+4=\boxed{7}
\]


\subsection{6.4 Complexity}

Sorting the edges takes:


\[
O(E\log E)
\]


With an efficient DSU implementation, the remaining component operations take approximately linear time with a very small amortized cost per operation.

Therefore, the overall time complexity is:


\[
\boxed{O(E\log E)}
\]


Since the number of edges in a simple graph is at most $O(V^2)$, sorting is often expressed as $O(E\log V)$ for connected simple graphs.

\textbf{Why is Kruskal greedy?} It selects the least-weight edge that can safely be added without forming a cycle. The cut property provides the basis for the correctness argument.

\section{Minimum Spanning Tree: Prim's Algorithm}

Prim's algorithm also constructs a minimum spanning tree, but its greedy choice differs from Kruskal's.

\subsection{7.1 Greedy strategy}

Start from any vertex and repeatedly add the minimum-weight edge connecting a vertex already in the tree to a vertex outside the tree.

Algorithm:

\begin{enumerate}
    \item Choose an initial vertex.
    \item Mark it as part of the growing tree.
    \item Consider edges connecting the tree to outside vertices.
    \item Select the least-weight such edge.
    \item Add the new vertex and edge.
    \item Repeat until all vertices are included.
\end{enumerate}

\subsection{7.2 Kruskal versus Prim}

\begin{table}[htbp]
\centering
\caption{Detailed Comparison between Kruskal's and Prim's MST Algorithms}
\label{tab:kruskal_vs_prim}
\begin{tabularx}{\textwidth}{p{3.5cm} X X}
\toprule
\textbf{Evaluation Basis} & \textbf{Kruskal's Algorithm} & \textbf{Prim's Algorithm} \\
\midrule
\textbf{Main Decision} & Add smallest weight edge not creating cycle & Add smallest weight edge crossing cut to outside \\
\textbf{Starting Point} & No specific start vertex; global edge sorting & Starts from an arbitrary root vertex \\
\textbf{Growth Pattern} & Forest of trees merging into single MST & Grows single connected tree continuously \\
\textbf{Data Structures} & Edge sorting plus Disjoint Set Union (DSU) & Min-priority queue (heap) plus adjacency list \\
\textbf{Complexity} & $O(E \log E)$ & $O(E \log V)$ (binary heap) / $O(V^2)$ (matrix) \\
\textbf{Suitable Graph} & Sparse graphs ($E \ll V^2$) & Dense graphs ($E \approx V^2$) with matrix; sparse with heap \\
\textbf{Output} & Minimum Spanning Tree & Minimum Spanning Tree \\
\bottomrule
\end{tabularx}
\end{table}

Both algorithms are greedy and produce a minimum spanning tree under the standard assumptions of a connected weighted undirected graph. If the graph is disconnected, their forest-producing counterparts yield a minimum spanning forest.

\section{Dijkstra's Shortest-Path Algorithm}

Dijkstra's algorithm finds shortest paths from a chosen source to the other reachable vertices in a weighted graph with \textbf{nonnegative edge weights}.

It uses a greedy strategy: at each step, it finalizes the unprocessed vertex with the smallest known tentative distance.

\subsection{8.1 Algorithm}

\begin{enumerate}
    \item Set the source distance to 0.
    \item Set all other tentative distances to infinity.
    \item Select the unprocessed vertex with the smallest tentative distance.
    \item For every outgoing edge, check whether travelling through the selected vertex improves the neighbor's distance.
    \item Update distances and predecessor pointers when a shorter route is found.
    \item Continue until the priority queue is empty or the required target has been finalized.
\end{enumerate}

The relaxation operation is:


\[
\boxed{\text{if }dist[u]+w(u,v)<dist[v]}
\]


then update:


\[
dist[v]=dist[u]+w(u,v)
\]


\subsection{8.2 Worked example}

Consider the graph with these edge weights:

\begin{table}[htbp]
\centering
\caption{Directed Graph Edge Weights for Dijkstra's Algorithm}
\label{tab:dijkstra_edges}
\begin{tabular}{cc}
\toprule
\textbf{Directed Edge} & \textbf{Weight} \\
\midrule
A $\to$ B & 4 \\
A $\to$ C & 1 \\
C $\to$ B & 2 \\
B $\to$ D & 1 \\
C $\to$ D & 5 \\
\bottomrule
\end{tabular}
\end{table}

Find the shortest distance from A to D.

Initially:


\[
dist(A)=0
\]


All other distances are infinity.

From A:

\begin{itemize}
    \item B gets tentative distance 4.
    \item C gets tentative distance 1.
\end{itemize}

The next vertex with the smallest tentative distance is C.

Through C:

\begin{itemize}
    \item B improves from 4 to $1+2=3$.
    \item D gets tentative distance $1+5=6$.
\end{itemize}

The next vertex is B, with distance 3.

Through B:

\- D improves from 6 to $3+1=4$.

D is then reached with distance 4.

The shortest path is:


\[
\boxed{A\rightarrow C\rightarrow B\rightarrow D}
\]


Total weight:


\[
1+2+1=\boxed{4}
\]


\subsection{8.3 Complexity}

With an adjacency list and a binary heap:


\[
\boxed{O((V+E)\log V)}
\]


is a common bound for a lazy-priority-queue implementation on a graph with a standard number of vertices and edges.

With an adjacency matrix and a simple array-based minimum selection, the complexity is:


\[
\boxed{O(V^2)}
\]


\subsection{8.4 Important limitation}

Dijkstra's algorithm assumes nonnegative edge weights. Negative edges can invalidate its greedy finalization rule.

For graphs containing negative edge weights, Bellman--Ford is a common alternative for single-source shortest paths, provided there is no reachable negative cycle that makes the path cost unbounded below.

\textbf{Exam comparison:}

\begin{itemize}
    \item BFS: shortest paths by number of edges in an unweighted graph.
    \item Dijkstra: shortest paths with nonnegative edge weights.
    \item Bellman--Ford: supports negative edge weights and can detect reachable negative cycles.
    \item Prim and Kruskal: minimum spanning trees, not shortest paths from a source.
\end{itemize}

\section{Huffman Coding}

Huffman Coding is a lossless compression technique that constructs a variable-length, prefix-free code from symbol frequencies.

Symbols that appear more frequently generally receive shorter codes, while less frequent symbols receive longer codes.

\subsection{9.1 Greedy strategy}

Repeatedly combine the two nodes with the smallest frequencies.

Algorithm:

\begin{enumerate}
    \item Create a leaf node for each symbol, with its frequency as its weight.
    \item Insert all nodes into a min-priority queue.
    \item Remove the two nodes with the smallest weights.
    \item Create a parent node whose weight is their sum.
    \item Insert the new parent into the priority queue.
    \item Continue until only one tree remains.
\end{enumerate}

The resulting tree defines the symbol codes: for example, left edges may represent \texttt{0}, and right edges may represent \texttt{1}.

\subsection{9.2 Worked example}

Suppose the symbols have these frequencies:

\begin{table}[htbp]
\centering
\caption{Symbol Frequency Table for Huffman Coding}
\label{tab:huffman_symbols}
\begin{tabular}{cc}
\toprule
\textbf{Symbol} & \textbf{Frequency} \\
\midrule
A & 5 \\
B & 9 \\
C & 12 \\
D & 13 \\
E & 16 \\
F & 45 \\
\bottomrule
\end{tabular}
\end{table}

The two smallest frequencies are 5 and 9.

Combine them:


\[
5+9=14
\]


The frequency list becomes:

\texttt{12, 13, 14, 16, 45}

Combine 12 and 13:


\[
12+13=25
\]


The remaining weights become:

\texttt{14, 16, 25, 45}

Combine 14 and 16:


\[
14+16=30
\]


The remaining weights are:

\texttt{25, 30, 45}

Combine 25 and 30:


\[
25+30=55
\]


Finally, combine 45 and 55:


\[
45+55=100
\]


A single tree remains.

The exact binary code assigned to each symbol depends on the left/right conventions and tie-breaking rules, but the construction repeatedly combines the two lowest frequencies.

\subsection{9.3 Complexity}

Building the tree using a min-priority queue takes:


\[
\boxed{O(n\log n)}
\]


for $n$ distinct symbols.

\textbf{Why is it greedy?} At every step, it combines the two least frequent remaining nodes. Huffman's correctness follows from the fact that an optimal prefix-code tree can be chosen with the two least frequent symbols as deepest siblings, after which they can be combined into one reduced subproblem.

\section{Job Sequencing with Deadlines}

In the standard \textbf{Job Sequencing with Deadlines} problem, each job:

\begin{itemize}
    \item Takes one unit of time.
    \item Has a deadline.
    \item Provides a profit only if it is completed by its deadline.
\end{itemize}

The objective is to maximize total profit.

\subsection{10.1 Greedy strategy}

Sort jobs in descending order of profit.

For each job, assign it to the latest available time slot that does not exceed its deadline.

Placing a job as late as possible leaves earlier slots available for jobs with tighter deadlines.

\subsection{10.2 Worked example}

Consider the following jobs:

\begin{table}[htbp]
\centering
\caption{Job Set with Deadlines and Profits}
\label{tab:job_sequencing}
\begin{tabular}{ccc}
\toprule
\textbf{Job ID} & \textbf{Deadline} & \textbf{Profit} \\
\midrule
A & 2 & 100 \\
B & 1 & 19 \\
C & 2 & 27 \\
D & 1 & 25 \\
E & 3 & 15 \\
\bottomrule
\end{tabular}
\end{table}

Order by profit:

\texttt{A, C, D, B, E}

There are three available time slots because the maximum deadline is 3.

\textbf{Schedule A:} Put A in slot 2, the latest slot at or before its deadline.

Slots:

\texttt{[empty, A, empty]}

\textbf{Schedule C:} Its deadline is 2. Slot 2 is occupied, so put C in slot 1.

Slots:

\texttt{[C, A, empty]}

\textbf{Schedule D:} Its deadline is 1, but slot 1 is occupied. Reject it.

\textbf{Schedule B:} Its deadline is 1, which is also occupied. Reject it.

\textbf{Schedule E:} Put E in slot 3.

Final schedule:

\texttt{[C, A, E]}

Total profit:


\[
27+100+15=\boxed{142}
\]


\subsection{10.3 Complexity}

Sorting takes $O(n\log n)$.

If each job scans backward through the time slots to locate an available slot, the scheduling phase may take $O(nD)$, where $D$ is the maximum deadline. A simple implementation with $D=O(n)$ can therefore have $O(n^2)$ overall time.

Disjoint Set Union can optimize the search for the latest free slot, reducing the scheduling overhead substantially.

\textbf{Important assumption:} The standard greedy algorithm described here applies to unit-duration jobs with deadlines and profit. More general jobs with arbitrary durations require different scheduling methods.

\section{Summary of Greedy Algorithms}

\begin{table}[htbp]
\centering
\caption{Master Summary of Canonical Greedy Algorithms}
\label{tab:greedy_summary}
\begin{tabularx}{\textwidth}{p{3.8cm} X p{4.5cm}}
\toprule
\textbf{Problem} & \textbf{Greedy Choice Strategy} & \textbf{Time Complexity} \\
\midrule
Activity Selection & Earliest finish time of compatible activity & $O(n \log n)$ ($O(n)$ if pre-sorted) \\
Fractional Knapsack & Highest value-to-weight ratio ($v_i/w_i$) first & $O(n \log n)$ \\
Kruskal's MST & Lowest weight edge connecting disjoint sets & $O(E \log E)$ \\
Prim's MST & Lowest weight edge crossing cut boundary & $O(E \log V)$ (heap) / $O(V^2)$ (matrix) \\
Dijkstra's Shortest Path & Lowest tentative distance node & $O((V+E)\log V)$ (heap) \\
Huffman Coding & Merge two lowest frequency trees & $O(n \log n)$ \\
Job Sequencing & Highest profit job placed in latest free slot & $O(n^2)$ standard / $O(n \log n)$ with DSU \\
Coin Change & Largest denomination coin first & Optimal only for canonical coin systems \\
\bottomrule
\end{tabularx}
\end{table}

Greedy algorithms can be simple and efficient, but correctness depends on the problem's properties and the proof for the specific greedy choice.


\newpage
\begin{center}
\vspace{0.5em}
{\color{primary}\rule{0.45\textwidth}{1.2pt}}\\[4pt]
{\Large\bfseries\color{primary} Section II: Recursion and Recurrence Relations}\\[2pt]
{\color{primary}\rule{0.45\textwidth}{1.2pt}}
\vspace{0.5em}
\end{center}


\section{What Is Recursion?}

\textbf{Recursion} is a programming technique in which a function calls itself to solve a problem by reducing it to smaller instances of the same problem.

A recursive solution consists of two essential parts:

\begin{enumerate}
    \item \textbf{Base case:} A condition under which the function returns without making another recursive call.
    \item \textbf{Recursive case:} A rule that reduces the problem and calls the function again.
\end{enumerate}

Without a correct base case or guaranteed progress toward it, recursion may continue indefinitely until the program runs out of call-stack space or otherwise fails.

\subsubsection{Example}

The factorial of a nonnegative integer $n$ is defined as:


\[
n!= \begin{cases} 1, & n=0\\\ n(n-1)!, & n>0 \end{cases}
\]


For example:


\[
5!=5\times4\times3\times2\times1=120
\]


The recursive definition expresses the factorial of a number in terms of the factorial of a smaller number.

\subsection{12.1 How a recursive function executes}

Consider:

\begin{lstlisting}[style=cppstyle]
int factorial(int n) {
    if (n == 0) {
        return 1;
    }

    return n * factorial(n - 1);
}
\end{lstlisting}

When \texttt{factorial(4)} is called:

\begin{enumerate}
    \item \texttt{factorial(4)} calls \texttt{factorial(3)}.
    \item \texttt{factorial(3)} calls \texttt{factorial(2)}.
    \item \texttt{factorial(2)} calls \texttt{factorial(1)}.
    \item \texttt{factorial(1)} calls \texttt{factorial(0)}.
    \item \texttt{factorial(0)} returns 1.
    \item The waiting calls complete in reverse order.
\end{enumerate}

The returns are:


\[
factorial(1)=1\times1=1
\]



\[
factorial(2)=2\times1=2
\]



\[
factorial(3)=3\times2=6
\]



\[
factorial(4)=4\times6=\boxed{24}
\]


Each active function call has its own execution context, commonly represented by a stack frame on the program's call stack.

\subsection{12.2 The call stack}

The call stack tracks active function calls, including information such as local variables, parameters, and return locations.

With recursion, new calls are placed above the calling functions. When a call returns, its frame is removed and control resumes in the calling function.

This naturally follows LIFO order: the most recently invoked unfinished function returns first.

\subsubsection{Why can deep recursion fail?}

If recursion goes too deep, the call stack can be exhausted. This can cause a stack overflow or other runtime failure, depending on the execution environment.

A recursive algorithm should therefore be evaluated for:

\begin{itemize}
    \item Maximum recursion depth.
    \item Correct base conditions.
    \item Guaranteed progress toward the base condition.
    \item Repeated computations.
    \item Whether an iterative implementation would be more suitable.
\end{itemize}

\section{Types of Recursion}

\subsection{13.1 Direct recursion}

A function calls itself directly.

\begin{lstlisting}[style=cppstyle]
void countdown(int n) {
    if (n <= 0) {
        return;
    }

    cout << n << " ";
    countdown(n - 1);
}
\end{lstlisting}

Calling \texttt{countdown(5)} prints:

\begin{lstlisting}[style=pseudocode]
5 4 3 2 1
\end{lstlisting}

\subsection{13.2 Indirect or mutual recursion}

Two or more functions call one another, eventually returning to a function that was already involved in the chain.

For example:

\begin{lstlisting}[style=cppstyle]
bool isEven(int n);

bool isOdd(int n) {
    if (n == 0) return false;
    return isEven(n - 1);
}

bool isEven(int n) {
    if (n == 0) return true;
    return isOdd(n - 1);
}
\end{lstlisting}

This example demonstrates mutual recursion for nonnegative integers. A production implementation should validate the input domain.

\subsection{13.3 Tail recursion}

In \textbf{tail recursion}, the recursive call is the final operation performed by the function.

\begin{lstlisting}[style=cppstyle]
int sumTail(int n, int accumulator = 0) {
    if (n == 0) {
        return accumulator;
    }

    return sumTail(n - 1, accumulator + n);
}
\end{lstlisting}

The recursive call to \texttt{sumTail} is the final action.

Some compilers can optimize tail calls in suitable circumstances, but C++ does not guarantee tail-call optimization. Therefore, a recursive tail call should not automatically be assumed to use constant stack space.

\subsection{13.4 Non-tail recursion}

A recursive function is non-tail recursive when additional work remains after the recursive call returns.

Factorial is a classic example:

\begin{lstlisting}[style=cppstyle]
return n * factorial(n - 1);
\end{lstlisting}

After the recursive call completes, multiplication is still required.

The caller must retain information about the pending multiplication until the call returns.

\subsection{13.5 Linear recursion}

Each function call makes at most one recursive call.

Factorial, GCD using Euclid's algorithm, and recursive binary search are examples.

\subsection{13.6 Tree recursion}

Each function call makes more than one recursive call.

Naive Fibonacci is a common example:

\begin{lstlisting}[style=cppstyle]
int fibonacci(int n) {
    if (n <= 1) {
        return n;
    }

    return fibonacci(n - 1) + fibonacci(n - 2);
}
\end{lstlisting}

This creates a branching recursion tree and repeats many calculations.

For example, both \texttt{fibonacci(5)} and \texttt{fibonacci(4)} eventually call \texttt{fibonacci(3)}, so the same subproblems are calculated repeatedly.

This leads directly to the motivation for Dynamic Programming.

\section{Factorial Using Recursion}

\subsection{14.1 Algorithm}

\begin{lstlisting}[style=cppstyle]
FACTORIAL(n):
    if n == 0:
        return 1

    return n * FACTORIAL(n - 1)
\end{lstlisting}

A negative factorial is not defined in the usual nonnegative-integer factorial problem. A correct program should reject negative inputs rather than recurse indefinitely.

\subsection{14.2 C++ implementation}

\begin{lstlisting}[style=cppstyle]
#include <iostream>
using namespace std;

unsigned long long factorial(int n) {
    if (n < 0) {
        throw invalid_argument("n must be nonnegative");
    }

    if (n == 0 || n == 1) {
        return 1;
    }

    return static_cast<unsigned long long>(n) * factorial(n - 1);
}

int main() {
    int n;
    cin >> n;

    try {
        cout << factorial(n) << '\n';
    } catch (const invalid_argument& e) {
        cout << e.what() << '\n';
    }

    return 0;
}
\end{lstlisting}

This implementation requires the appropriate \texttt{<stdexcept>} header for \texttt{invalid\_argument}. The numeric return type also has a limited range, so large factorials will overflow even though the mathematical result is defined.

\subsection{14.3 Complexity analysis}

For each input value $n$, the algorithm makes a sequence of $n$ recursive calls until reaching the base case.

The recurrence is:


\[
T(n)=T(n-1)+O(1)
\]


with constant-time work per call.

Therefore:


\[
\boxed{T(n)=O(n)}
\]


The maximum number of simultaneously active calls grows linearly with $n$, so the auxiliary call-stack space is:


\[
\boxed{O(n)}
\]


An iterative factorial would use $O(1)$ auxiliary space.

\section{Fibonacci Sequence: Recursion and Its Cost}

The Fibonacci sequence is defined as:


\[
F(0)=0,\qquad F(1)=1
\]


and for $n\geq2$:


\[
\boxed{F(n)=F(n-1)+F(n-2)}
\]


The first values are:

\texttt{0, 1, 1, 2, 3, 5, 8, 13, 21, ...}

\subsection{15.1 Naive recursive implementation}

\begin{lstlisting}[style=cppstyle]
long long fibonacci(int n) {
    if (n <= 1) {
        return n;
    }

    return fibonacci(n - 1) + fibonacci(n - 2);
}
\end{lstlisting}

This implementation repeatedly calculates the same values.

For example, \texttt{fibonacci(5)} calls \texttt{fibonacci(4)} and \texttt{fibonacci(3)}, and the computation of \texttt{fibonacci(4)} itself calls \texttt{fibonacci(3)} again.

The repeated call to \texttt{fibonacci(3)} illustrates overlapping subproblems.

\subsubsection{Complexity}

The recurrence is:


\[
T(n)=T(n-1)+T(n-2)+O(1)
\]


The running time grows exponentially. A commonly quoted upper bound is:


\[
\boxed{O(2^n)}
\]


More precisely, the growth is $\Theta(\varphi^n)$, where $\varphi=(1+\sqrt5)/2$, under the standard unit-cost model.

The maximum recursion depth is $O(n)$, so auxiliary stack space is $O(n)$.

\subsection{15.2 Why is naive Fibonacci inefficient?}

A large amount of work is repeated because the same Fibonacci values are calculated many times.

For example, the value of $F(3)$ is needed when evaluating $F(4)$, but it is also needed separately when evaluating $F(5)$.

The naive recursive method does not remember completed results.

Dynamic Programming solves this by storing already computed Fibonacci values and reusing them.

This leads to memoization and tabulation.

\section{Greatest Common Divisor Using Recursion}

Euclid's algorithm finds the greatest common divisor (GCD) of two integers.

Its key identity is:


\[
\boxed{\gcd(a,b)=\gcd(b,a\bmod b)}
\]


provided $b\neq0$. The base case is:


\[
\boxed{\gcd(a,0)=|a|}
\]


\subsection{16.1 Step-by-step example}

Find:


\[
\gcd(48,18)
\]


First:


\[
48\bmod18=12
\]


Therefore:


\[
\gcd(48,18)=\gcd(18,12)
\]


Next:


\[
18\bmod12=6
\]


So:


\[
\gcd(18,12)=\gcd(12,6)
\]


Finally:


\[
12\bmod6=0
\]


Therefore:


\[
\gcd(12,6)=\gcd(6,0)=\boxed{6}
\]


\subsection{16.2 C++ implementation}

\begin{lstlisting}[style=cppstyle]
#include <iostream>
#include <cstdlib>
using namespace std;

long long gcdRecursive(long long a, long long b) {
    if (b == 0) {
        return llabs(a);
    }

    return gcdRecursive(b, a % b);
}

int main() {
    long long a, b;
    cin >> a >> b;

    cout << gcdRecursive(a, b) << '\n';

    return 0;
}
\end{lstlisting}

\subsubsection{Complexity}

Euclid's algorithm reduces the size of the arguments rapidly.

Its time complexity is:


\[
\boxed{O(\log(\min(|a|,|b|)+1))}
\]


for ordinary nontrivial inputs, stated asymptotically as $O(\log \min(|a|,|b|))$ when both arguments are positive.

The recursive call-stack space is logarithmic in the size of the arguments.

\textbf{Exam point:} Euclid's algorithm is much more efficient than testing every possible divisor from 1 up to the smaller input.

\section{Recursive Binary Search}

Binary Search can be implemented recursively because each step reduces the problem to a smaller search interval.

\subsection{17.1 C++ implementation}

\begin{lstlisting}[style=cppstyle]
#include <iostream>
#include <vector>
using namespace std;

int binarySearchRecursive(
    const vector<int>& a,
    int target,
    int low,
    int high
) {
    if (low > high) {
        return -1;
    }

    int mid = low + (high - low) / 2;

    if (a[mid] == target) {
        return mid;
    }

    if (target < a[mid]) {
        return binarySearchRecursive(
            a, target, low, mid - 1
        );
    }

    return binarySearchRecursive(
        a, target, mid + 1, high
    );
}

int main() {
    vector<int> a = {2, 5, 8, 12, 16, 23, 38, 56};

    cout << binarySearchRecursive(
        a, 23, 0, static_cast<int>(a.size()) - 1
    ) << '\n';

    return 0;
}
\end{lstlisting}

Output:

\begin{lstlisting}[style=pseudocode]
5
\end{lstlisting}

The vector is sorted in ascending order, as required by binary search.

\subsubsection{Complexity}

The search interval is approximately halved in each call:


\[
T(n)=T(n/2)+O(1)
\]


Therefore:


\[
\boxed{T(n)=O(\log n)}
\]


The maximum number of simultaneously active recursive calls is also $O(\log n)$, so auxiliary space is $O(\log n)$.

The iterative implementation achieves the same time complexity with $O(1)$ auxiliary space.

\section{Tower of Hanoi}

The \textbf{Tower of Hanoi} is a classical recursion problem involving three rods and $n$ disks of different sizes.

Initially, all disks are stacked on one rod in decreasing size order, with the largest at the bottom.

The goal is to move the entire stack to another rod, obeying two rules:

\begin{enumerate}
    \item Move only one disk at a time.
    \item Never place a larger disk on top of a smaller disk.
\end{enumerate}

\subsection{18.1 Recursive strategy}

To move $n$ disks from source rod A to destination rod C using auxiliary rod B:

\begin{enumerate}
    \item Move $n-1$ disks from A to B, using C.
    \item Move the largest disk from A to C.
    \item Move $n-1$ disks from B to C, using A.
\end{enumerate}

This produces the recurrence:


\[
\boxed{T(n)=2T(n-1)+1}
\]


with:


\[
T(0)=0
\]


Solving the recurrence gives:


\[
\boxed{T(n)=2^n-1}
\]


This is the minimum number of moves needed.

\subsection{18.2 Example for three disks}

The minimum number of moves is:


\[
2^3-1=7
\]


\subsubsection{C++ implementation}

\begin{lstlisting}[style=cppstyle]
#include <iostream>
using namespace std;

void towerOfHanoi(int n, char source,
                  char auxiliary, char destination) {
    if (n == 0) {
        return;
    }

    towerOfHanoi(n - 1, source, destination, auxiliary);

    cout << "Move disk " << n
         << " from " << source
         << " to " << destination << '\n';

    towerOfHanoi(n - 1, auxiliary, source, destination);
}

int main() {
    int n = 3;

    towerOfHanoi(n, 'A', 'B', 'C');

    return 0;
}
\end{lstlisting}

\subsubsection{Complexity}

\begin{itemize}
    \item Time: $O(2^n)$.
    \item Auxiliary stack space: $O(n)$.
\end{itemize}

\textbf{Exam point:} The recursive depth is linear, even though the total number of operations is exponential.

This distinction is essential when analyzing recursive algorithms: total work and maximum call-stack depth are not the same quantity.

\section{Recurrence Relations and Complexity Analysis}

A \textbf{recurrence relation} expresses the running time of a recursive algorithm in terms of smaller input sizes.

For example, if an algorithm divides an input into two halves, recursively processes both halves, and spends linear time merging the results, its recurrence may be:


\[
\boxed{T(n)=2T(n/2)+O(n)}
\]


The recurrence describes the algorithm's structure. Solving it gives the asymptotic time complexity.

\subsection{19.1 Common recurrence patterns}

\begin{table}[htbp]
\centering
\caption{Common Recurrence Relations and Asymptotic Complexity}
\label{tab:recurrence_patterns}
\begin{tabularx}{\textwidth}{p{4.8cm} X p{3.2cm}}
\toprule
\textbf{Recurrence Relation} & \textbf{Algorithmic Paradigm / Example} & \textbf{Time Complexity} \\
\midrule
$T(n) = T(n-1) + O(1)$ & Factorial recursion, linear countdown & $O(n)$ \\
$T(n) = T(n-1) + O(n)$ & Selection Sort recursive, worst Quick Sort & $O(n^2)$ \\
$T(n) = T(n/2) + O(1)$ & Binary Search, divide by two & $O(\log n)$ \\
$T(n) = 2T(n/2) + O(n)$ & Merge Sort & $O(n \log n)$ \\
$T(n) = 2T(n-1) + O(1)$ & Tower of Hanoi & $O(2^n)$ \\
$T(n) = T(n-1) + T(n-2) + O(1)$ & Naive Fibonacci recursion & $\Theta(\varphi^n)$ \\
\bottomrule
\end{tabularx}
\end{table}

The table assumes the usual base cases and input-size interpretation.

\subsection{19.2 Master Theorem}

The Master Theorem is a method for solving many divide-and-conquer recurrences of the form:


\[
\boxed{T(n)=aT(n/b)+f(n)}
\]


where:

\begin{itemize}
    \item $a\geq1$ is the number of recursive subproblems.
    \item $b>1$ is the factor by which the input size is reduced.
    \item $f(n)$ is the non-recursive work at each call.
\end{itemize}

Define:


\[
c=\log_b a
\]


The standard three cases are:

\subsubsection{Case 1: Recursive work dominates}

If:


\[
f(n)=O(n^{c-\varepsilon})
\]


for some $\varepsilon>0$, then:


\[
\boxed{T(n)=\Theta(n^c)}
\]


\subsubsection{Case 2: Recursive and non-recursive work are balanced}

If:


\[
f(n)=\Theta(n^c\log^k n)
\]


for a constant $k\geq0$, then:


\[
\boxed{T(n)=\Theta(n^c\log^{k+1}n)}
\]


\subsubsection{Case 3: Non-recursive work dominates}

If:


\[
f(n)=\Omega(n^{c+\varepsilon})
\]


for some $\varepsilon>0$, together with the required regularity condition, then:


\[
\boxed{T(n)=\Theta(f(n))}
\]


The regularity condition typically requires that the recursive contribution be a fixed fraction smaller than $f(n)$ for sufficiently large $n$.

The theorem does not apply directly to every recurrence. For example, $T(n)=T(n-1)+O(1)$ and naive Fibonacci's recurrence are not in the standard divide-by-$b$ form.

\subsection{19.3 Worked Master Theorem problem}

\textbf{Question:} Solve the recurrence for Merge Sort:


\[
T(n)=2T(n/2)+O(n)
\]


Identify:


\[
a=2,\qquad b=2,\qquad f(n)=O(n)
\]


Calculate:


\[
c=\log_2 2=1
\]


Thus:


\[
n^c=n
\]


The non-recursive work has the same order as $n^c$, so the second case applies.

Therefore:


\[
\boxed{T(n)=\Theta(n\log n)}
\]


\subsection{19.4 When should the Master Theorem be used?}

It is appropriate when the recurrence follows the required divide-and-conquer form, such as:


\[
T(n)=aT(n/b)+f(n)
\]


It is not a universal method for every recursive function. For recurrences with irregular subproblem sizes, unequal branching, or dependencies on multiple previous terms, other techniques such as recursion trees, substitution, or characteristic equations may be necessary.

\section{Advantages and Disadvantages of Recursion}

\subsubsection{Advantages}

\begin{itemize}
    \item Expresses divide-and-conquer algorithms naturally.
    \item Makes tree and graph traversals straightforward.
    \item Can produce concise solutions for mathematical definitions.
    \item Supports backtracking and recursive decomposition.
    \item Can simplify algorithms such as Merge Sort and Quick Sort.
\end{itemize}

\subsubsection{Disadvantages}

\begin{itemize}
    \item Uses additional call-stack space.
    \item Can fail with excessive recursion depth.
    \item Can repeat the same computation many times.
    \item May have higher function-call overhead than an iterative solution.
    \item Requires careful base cases and progress toward termination.
\end{itemize}

\textbf{Practical rule:} Recursion is a design technique, not an automatic performance improvement. When a recursive solution repeatedly computes overlapping subproblems, memoization or Dynamic Programming may be necessary.


\newpage
\begin{center}
\vspace{0.5em}
{\color{primary}\rule{0.45\textwidth}{1.2pt}}\\[4pt]
{\Large\bfseries\color{primary} Section III: Dynamic Programming}\\[2pt]
{\color{primary}\rule{0.45\textwidth}{1.2pt}}
\vspace{0.5em}
\end{center}


\section{What Is Dynamic Programming?}

\textbf{Dynamic Programming (DP)} is an algorithm-design technique that solves a problem by dividing it into smaller subproblems and storing their results so that each subproblem does not need to be solved repeatedly.

It is especially useful when a problem has:

\begin{enumerate}
    \item \textbf{Optimal substructure:} The optimal solution can be constructed from optimal solutions to appropriately defined smaller subproblems.
    \item \textbf{Overlapping subproblems:} The same subproblems occur repeatedly during the computation.
\end{enumerate}

Many recursive problems have these properties. Dynamic Programming improves efficiency by recording the results of subproblems and reusing them.

\subsubsection{Practical example}

Suppose a program needs to calculate the Fibonacci sequence up to $F(100)$.

A naive recursive solution repeatedly calculates the same smaller Fibonacci values and has exponential running time.

Dynamic Programming computes each Fibonacci value once and stores it. The same problem can then be solved in linear time.

The main idea is not simply to use recursion. It is to \textbf{avoid redundant work by preserving the results of smaller computations}.

\section{Two Main Approaches to Dynamic Programming}

Dynamic Programming is commonly implemented using either memoization or tabulation.

\subsection{22.1 Memoization: Top-down DP}

Memoization starts with the original problem and solves smaller subproblems recursively.

Whenever a subproblem is solved, its result is stored. If the same subproblem is encountered again, the stored result is returned instead of recalculating it.

\subsubsection{Typical procedure}

\begin{enumerate}
    \item Define a recursive function.
    \item Identify the base cases.
    \item Before solving a subproblem, check whether its answer is already stored.
    \item If stored, return the saved answer.
    \item Otherwise, compute the answer recursively, save it, and return it.
\end{enumerate}

\subsubsection{Advantages}

\begin{itemize}
    \item Often resembles the natural recursive definition of the problem.
    \item Computes only the subproblems that are actually reached.
    \item Can be easier to implement when the state space is large but sparsely used.
\end{itemize}

\subsubsection{Disadvantages}

\begin{itemize}
    \item Uses recursion-stack space.
    \item May encounter stack-depth limitations.
    \item Has some overhead for checking and updating the memoization table.
\end{itemize}

\subsection{22.2 Tabulation: Bottom-up DP}

Tabulation solves smaller subproblems first, storing results in a table that is filled in a valid dependency order.

Instead of recursively calling itself, the algorithm builds up the answer through iteration.

\subsubsection{Typical procedure}

\begin{enumerate}
    \item Identify the smallest subproblems.
    \item Initialize the table with their known answers.
    \item Determine which previously computed states are needed for each new state.
    \item Fill the table in a valid order.
    \item Return the entry representing the original problem.
\end{enumerate}

\subsubsection{Advantages}

\begin{itemize}
    \item Avoids recursive call-stack overhead.
    \item Often provides predictable evaluation order.
    \item Can enable space optimization when only a few previous states are required.
\end{itemize}

\subsubsection{Disadvantages}

\begin{itemize}
    \item May calculate states that are not needed for a particular query.
    \item Requires careful selection of the table dimensions and iteration order.
    \item Can be less intuitive when the recurrence has complicated dependencies.
\end{itemize}

\subsection{22.3 Memoization versus tabulation}

\begin{table}[htbp]
\centering
\caption{Systematic Comparison: Memoization vs. Tabulation}
\label{tab:memoization_vs_tabulation}
\begin{tabularx}{\textwidth}{p{3.2cm} X X}
\toprule
\textbf{Dimension} & \textbf{Memoization (Top-Down)} & \textbf{Tabulation (Bottom-Up)} \\
\midrule
\textbf{Direction} & Top-down from target state & Bottom-up from base cases \\
\textbf{Implementation} & Recursion + memo cache & Iterative loops + DP table \\
\textbf{Evaluation Order} & Evaluates states on-demand as needed & Evaluates states in topological order \\
\textbf{Call Stack} & Uses system recursion call stack & No recursion stack; iterative memory only \\
\textbf{Unneeded States} & Skips states unreachable from target & Computes all planned states in table \\
\textbf{State Optimization} & Difficult to compress memory & Easy to reduce dimension (e.g. 2 variables) \\
\textbf{Common Example} & Memoized Fibonacci, recursive LCS & 1D Knapsack, Matrix Tabulation \\
\bottomrule
\end{tabularx}
\end{table}

Both approaches implement Dynamic Programming. Neither is universally better; the appropriate choice depends on the problem's state space, dependencies, and implementation constraints.

\section{A General Workflow for Designing a DP Algorithm}

For examination questions, use the following sequence.

\subsection{Step 1: Define the state}

Specify exactly what each DP entry represents.

For example:


\[
dp[i]=F(i)
\]


means that the entry \texttt{dp[i]} stores the $i$-th Fibonacci number.

In a knapsack problem, a state might be:


\[
dp[i][w]
\]


representing the maximum value achievable using the first $i$ items with capacity $w$.

A well-defined state makes the recurrence easier to derive and verify.

\subsection{Step 2: Identify the recurrence}

Express the answer to a larger state in terms of smaller states.

For Fibonacci:


\[
dp[n]=dp[n-1]+dp[n-2]
\]


For a 0/1 Knapsack state, compare the cases in which an item is excluded or included.

\subsection{Step 3: Define base cases}

Base cases are the smallest states whose answers are known directly.

Incorrect base cases are one of the most common causes of DP errors.

\subsection{Step 4: Determine the transition order}

Ensure that every state is calculated only after the states it depends on have been computed.

For example, if \texttt{dp[i]} depends on \texttt{dp[i-1]} and \texttt{dp[i-2]}, calculate the states in increasing order of $i$.

\subsection{Step 5: Calculate time and space complexity}

Determine:

\begin{itemize}
    \item Number of states.
    \item Cost of computing each state.
    \item Total number of transitions.
    \item Memory required to store the state information.
\end{itemize}

A common estimate is:


\[
\boxed{ \text{Total time} \approx \text{Number of states} \times \text{Work per state} }
\]


This estimate assumes that each state can be computed from previously available states without additional hidden costs.

\subsection{Step 6: Check the result}

Validate that the recurrence handles:

\begin{itemize}
    \item Minimum input sizes.
    \item Impossible or unreachable states.
    \item Boundary conditions.
    \item Integer overflow, when applicable.
    \item The exact quantity requested by the problem.
\end{itemize}

\section{Dynamic Programming Example 1: Fibonacci Numbers}

\subsection{24.1 Recursive recurrence}

The Fibonacci sequence follows:


\[
F(0)=0,\qquad F(1)=1
\]



\[
F(n)=F(n-1)+F(n-2)
\]


Naive recursion is inefficient because the same values are repeatedly calculated.

\subsection{24.2 Top-down DP using memoization}

\begin{lstlisting}[style=cppstyle]
#include <iostream>
#include <vector>
using namespace std;

long long fibonacciMemo(int n, vector<long long>& memo) {
    if (n <= 1) {
        return n;
    }

    if (memo[n] != -1) {
        return memo[n];
    }

    memo[n] = fibonacciMemo(n - 1, memo)
            + fibonacciMemo(n - 2, memo);

    return memo[n];
}

int main() {
    int n;
    cin >> n;

    if (n < 0) {
        cout << "n must be nonnegative\n";
        return 0;
    }

    vector<long long> memo(n + 1, -1);

    cout << fibonacciMemo(n, memo) << '\n';

    return 0;
}
\end{lstlisting}

This example assumes the requested Fibonacci value fits in \texttt{long long}. Larger values require a suitable arbitrary-precision representation.

\subsubsection{Complexity}

Each state from $0$ to $n$ is calculated at most once.

\begin{itemize}
    \item Time: $O(n)$.
    \item Memoization space: $O(n)$.
    \item Recursion-stack space: $O(n)$.
\end{itemize}

The overall auxiliary space is $O(n)$.

\subsection{24.3 Bottom-up DP using tabulation}

\begin{lstlisting}[style=cppstyle]
#include <iostream>
using namespace std;

long long fibonacciDP(int n) {
    if (n <= 1) {
        return n;
    }

    long long previous2 = 0;
    long long previous1 = 1;

    for (int i = 2; i <= n; i++) {
        long long current = previous1 + previous2;

        previous2 = previous1;
        previous1 = current;
    }

    return previous1;
}

int main() {
    int n;
    cin >> n;

    if (n < 0) {
        cout << "n must be nonnegative\n";
        return 0;
    }

    cout << fibonacciDP(n) << '\n';

    return 0;
}
\end{lstlisting}

This version stores only the two previous Fibonacci values because those are the only states needed to calculate the next one.

\subsubsection{Complexity}

\begin{itemize}
    \item Time: $O(n)$.
    \item Auxiliary space: $O(1)$.
\end{itemize}

This demonstrates \textbf{space optimization} in Dynamic Programming.

\subsection{24.4 Comparison of Fibonacci approaches}

\begin{table}[htbp]
\centering
\caption{Performance Comparison of Fibonacci Algorithms}
\label{tab:fibonacci_comparison}
\begin{tabularx}{\textwidth}{p{5.5cm} c c}
\toprule
\textbf{Algorithmic Approach} & \textbf{Time Complexity} & \textbf{Auxiliary Space} \\
\midrule
Naive Tree Recursion & $\Theta(\varphi^n) \approx 1.618^n$ (Exponential) & $O(n)$ Call Stack \\
Top-Down DP (Memoization) & $O(n)$ & $O(n)$ Cache + Stack \\
Bottom-Up DP (Full Array) & $O(n)$ & $O(n)$ Table \\
Bottom-Up DP (Two Variables) & $O(n)$ & $O(1)$ Space Optimized \\
\bottomrule
\end{tabularx}
\end{table}

The key reason DP is faster is that each Fibonacci value is calculated once rather than repeatedly.

\section{Dynamic Programming Example 2: 0/1 Knapsack}

\subsection{25.1 Problem definition}

The 0/1 Knapsack problem asks us to choose items with given weights and values so that their total weight does not exceed a capacity $W$, while maximizing total value.

Each item can either be selected once or not selected. Fractions are not permitted.

This is a classic Dynamic Programming problem because the optimal decision depends on smaller subproblems involving fewer items and smaller capacities.

\subsection{25.2 Define the state}

Let:


\[
\boxed{dp[i][w]}
\]


represent the maximum achievable value using the first $i$ items with a knapsack capacity of $w$.

Suppose item $i$ has weight $wt_i$ and value $val_i$.

There are two possible decisions.

\subsubsection{Case 1: Exclude the item}

The value remains:


\[
dp[i-1][w]
\]


\subsubsection{Case 2: Include the item}

The item can be selected only if:


\[
wt_i\leq w
\]


The resulting value is:


\[
val_i+dp[i-1][w-wt_i]
\]


Therefore, for a feasible item:


\[
\boxed{ dp[i][w]= \max\left( dp[i-1][w], val_i+dp[i-1][w-wt_i] \right) }
\]


If the item is too heavy, the state is:


\[
\boxed{dp[i][w]=dp[i-1][w]}
\]


Base cases:


\[
dp[0][w]=0
\]


and:


\[
dp[i][0]=0
\]


With no items or zero capacity, the maximum achievable value is zero.

\subsection{25.3 Worked example}

Suppose the capacity is:


\[
W=5
\]


The items are:

\begin{table}[htbp]
\centering
\caption{0/1 Knapsack DP Example Items ($W = 5$)}
\label{tab:knapsack_dp_items}
\begin{tabular}{ccc}
\toprule
\textbf{Item ($i$)} & \textbf{Weight ($wt_i$)} & \textbf{Value ($val_i$)} \\
\midrule
1 & 2 & 12 \\
2 & 1 & 10 \\
3 & 3 & 20 \\
4 & 2 & 15 \\
\bottomrule
\end{tabular}
\end{table}

Consider the combinations:

\begin{itemize}
    \item Items 1 and 2: weight $3$, value $22$.
    \item Items 1 and 4: weight $4$, value $27$.
    \item Items 2, 3, and 4: weight $6$, which exceeds capacity.
    \item Items 2 and 3: weight $4$, value $30$.
    \item Items 1 and 3: weight $5$, value $32$.
\end{itemize}

The combination of items 1 and 3 has value 32, which is the maximum.

Therefore:


\[
\boxed{\text{Optimal value}=32}
\]


\subsection{25.4 C++ implementation with one-dimensional DP}

The full two-dimensional table can be compressed to a one-dimensional array.

\begin{lstlisting}[style=cppstyle]
#include <algorithm>
#include <iostream>
#include <vector>
using namespace std;

int knapsack01(
    const vector<int>& weights,
    const vector<int>& values,
    int capacity
) {
    int n = static_cast<int>(weights.size());
    vector<int> dp(capacity + 1, 0);

    for (int i = 0; i < n; i++) {
        // Iterate downward to ensure each item is used once.
        for (int w = capacity; w >= weights[i]; w--) {
            dp[w] = max(
                dp[w],
                values[i] + dp[w - weights[i]]
            );
        }
    }

    return dp[capacity];
}

int main() {
    vector<int> weights = {2, 1, 3, 2};
    vector<int> values = {12, 10, 20, 15};

    cout << knapsack01(weights, values, 5) << '\n';

    return 0;
}
\end{lstlisting}

Output:

\begin{lstlisting}[style=pseudocode]
32
\end{lstlisting}

\subsubsection{Why iterate downward?}

This is a critical implementation detail.

The state \texttt{dp[w]} represents the best value achievable at capacity \texttt{w}.

When processing one item, iterating downward ensures that the value from the preceding item iteration is used, rather than a value updated earlier in the same iteration. Therefore, the current item cannot be selected multiple times.

If the inner loop iterates upward, the newly updated values may be reused during the same item iteration. That produces the recurrence for \textbf{unbounded Knapsack}, where an item can be taken multiple times, rather than 0/1 Knapsack.

\subsection{25.5 Complexity}

With $n$ items and integer capacity $W$:

\begin{itemize}
    \item Time: $O(nW)$.
    \item Auxiliary space: $O(W)$ for the one-dimensional implementation.
    \item Full two-dimensional table space: $O(nW)$.
\end{itemize}

The $O(nW)$ time is called \textbf{pseudo-polynomial}, not polynomial in the length of the input representation in general. The numerical value of $W$ can be exponentially larger than the number of bits needed to represent $W$.

\textbf{Exam distinction:} 0/1 Knapsack DP is a standard solution when capacity is reasonably bounded. For large numeric capacities, other methods or approximations may be necessary.

\section{Dynamic Programming Example 3: Coin Change with Minimum Coins}

This problem differs from the earlier greedy coin-change example.

Given a set of coin denominations and an amount $A$, find the minimum number of coins required to make the amount, assuming unlimited availability of each denomination.

\subsection{26.1 Define the state}

Let:


\[
dp[x]
\]


be the minimum number of coins required to make amount $x$.

Base case:


\[
dp[0]=0
\]


For each positive amount $x$, consider every coin denomination $c$ such that $c\leq x$.

If $x-c$ can be formed, then choosing coin $c$ gives:


\[
dp[x-c]+1
\]


Therefore:


\[
\boxed{ dp[x]=1+\min_{\substack{c\leq x\\ c\text{ is a coin}}}dp[x-c] }
\]


If no denomination can form the amount, the state remains unreachable.

\subsection{26.2 Example}

Coins:


\[
\\{1,3,4\\}
\]


Amount:


\[
A=6
\]


The minimum counts begin:

\begin{table}[htbp]
\centering
\caption{DP State Table for Coin Change ($\{1, 3, 4\}$, Target 6)}
\label{tab:coin_change_dp}
\begin{tabular}{ccc}
\toprule
\textbf{Target Amount ($x$)} & \textbf{Minimum Coins ($dp[x]$)} & \textbf{Optimal Coin Selection} \\
\midrule
0 & 0 & Empty selection \\
1 & 1 & $1$ \\
2 & 2 & $1 + 1$ \\
3 & 1 & $3$ \\
4 & 1 & $4$ \\
5 & 2 & $4 + 1$ \\
6 & 2 & $3 + 3$ \\
\bottomrule
\end{tabular}
\end{table}

The optimal answer is:


\[
\boxed{2\text{ coins}}
\]


This example also demonstrates why taking the largest coin greedily is not always optimal: greedy selects $4+1+1$, which uses three coins.

\subsection{26.3 C++ implementation}

\begin{lstlisting}[style=cppstyle]
#include <algorithm>
#include <iostream>
#include <vector>
using namespace std;

int minimumCoins(const vector<int>& coins, int amount) {
    const int INF = amount + 1;

    vector<int> dp(amount + 1, INF);
    dp[0] = 0;

    for (int x = 1; x <= amount; x++) {
        for (int coin : coins) {
            if (coin <= x && dp[x - coin] != INF) {
                dp[x] = min(dp[x], dp[x - coin] + 1);
            }
        }
    }

    return dp[amount] == INF ? -1 : dp[amount];
}

int main() {
    vector<int> coins = {1, 3, 4};

    cout << minimumCoins(coins, 6) << '\n';

    return 0;
}
\end{lstlisting}

Output:

\begin{lstlisting}[style=pseudocode]
2
\end{lstlisting}

This implementation assumes all coin denominations are positive integers.

\subsection{26.4 Complexity}

Let $m$ be the number of coin denominations and $A$ the target amount.

\begin{itemize}
    \item Time: $O(Am)$.
    \item Auxiliary space: $O(A)$.
\end{itemize}

If the target cannot be represented, the function returns \texttt{-1}.

\textbf{Exam distinction:} The recurrence is similar to 0/1 Knapsack in spirit, but unlimited reuse of coin denominations is allowed. This changes the state transitions and implementation details.

\section{Dynamic Programming Example 4: Longest Common Subsequence (LCS)}

\subsection{27.1 Problem definition}

The \textbf{Longest Common Subsequence} problem asks for the longest sequence of elements that appears in the same relative order in two sequences, though the elements need not be adjacent.

For example:

\begin{lstlisting}[style=pseudocode]
String A: ABCDGH
String B: AEDFHR
\end{lstlisting}

One common subsequence is \texttt{ADH}.

The characters appear in the same order in both strings, so the subsequence has length 3.

A subsequence differs from a substring: a subsequence may skip elements, whereas a substring must consist of consecutive characters.

\subsection{27.2 Define the DP state}

Let:


\[
dp[i][j]
\]


represent the length of the LCS of the first $i$ characters of string A and the first $j$ characters of string B.

Base cases:


\[
dp[0][j]=0
\]



\[
dp[i][0]=0
\]


If either prefix is empty, the longest common subsequence has length zero.

\subsubsection{Case 1: Matching characters}

If:


\[
A[i-1]=B[j-1]
\]


then:


\[
\boxed{dp[i][j]=1+dp[i-1][j-1]}
\]


The matching character extends a common subsequence of the preceding prefixes.

\subsubsection{Case 2: Different characters}

If:


\[
A[i-1]\neq B[j-1]
\]


then:


\[
\boxed{dp[i][j]=\max(dp[i-1][j],dp[i][j-1])}
\]


We compare the solutions obtained by excluding the final character from one sequence or the other.

\subsection{27.3 C++ implementation}

\begin{lstlisting}[style=cppstyle]
#include <algorithm>
#include <iostream>
#include <string>
#include <vector>
using namespace std;

int lcsLength(const string& a, const string& b) {
    int m = static_cast<int>(a.size());
    int n = static_cast<int>(b.size());

    vector<vector<int>> dp(
        m + 1,
        vector<int>(n + 1, 0)
    );

    for (int i = 1; i <= m; i++) {
        for (int j = 1; j <= n; j++) {
            if (a[i - 1] == b[j - 1]) {
                dp[i][j] = 1 + dp[i - 1][j - 1];
            } else {
                dp[i][j] = max(
                    dp[i - 1][j],
                    dp[i][j - 1]
                );
            }
        }
    }

    return dp[m][n];
}

int main() {
    string a = "ABCBDAB";
    string b = "BDCABA";

    cout << lcsLength(a, b) << '\n';

    return 0;
}
\end{lstlisting}

Output:

\begin{lstlisting}[style=pseudocode]
4
\end{lstlisting}

The strings have an LCS of length four. For example, \texttt{BCBA} is a valid common subsequence.

\subsection{27.4 Complexity}

For strings of lengths $m$ and $n$:

\begin{itemize}
    \item Time: $O(mn)$.
    \item Space: $O(mn)$.
\end{itemize}

If only the length is required, the memory can be reduced to $O(\min(m,n))$ using two rows, because each state depends only on the previous row and the current row's previous entry.

If the actual subsequence must be reconstructed, the implementation usually needs to retain sufficient information to trace the decisions, or recompute portions of the table.

\textbf{Exam comparison:} LCS can be solved by naive recursion, but the repeated subproblems make Dynamic Programming substantially more efficient.

\section{Dynamic Programming Example 5: Weighted Interval Scheduling}

This problem illustrates a situation where a greedy strategy based only on finish time does not necessarily maximize total profit.

\subsection{28.1 Problem definition}

Given intervals with start times, finish times, and associated profits, select a set of mutually compatible intervals that maximizes total profit.

Unlike the ordinary Activity Selection problem, each interval has a different value.

\subsection{28.2 Why does the simple greedy method fail?}

Suppose the intervals are:

\begin{table}[htbp]
\centering
\caption{Weighted Interval Scheduling Example Dataset}
\label{tab:weighted_intervals}
\begin{tabular}{cccc}
\toprule
\textbf{Interval} & \textbf{Start Time} & \textbf{Finish Time} & \textbf{Profit ($v_i$)} \\
\midrule
A & 1 & 10 & 100 \\
B & 1 & 4 & 40 \\
C & 4 & 7 & 40 \\
D & 7 & 10 & 40 \\
\bottomrule
\end{tabular}
\end{table}

A finish-time greedy algorithm chooses B, C, and D, earning:


\[
40+40+40=120
\]


In this example, that happens to outperform selecting A, but it demonstrates how weights change the objective. The ordinary earliest-finish rule is not itself an optimality proof for the weighted variant.

To give a clear counterexample, change A's profit to 200. The finish-time greedy selection still returns B, C, and D with total profit 120, whereas selecting A alone gives 200.

Thus, when the objective is total profit, a different algorithm is needed.

\subsection{28.3 DP state and recurrence}

Sort intervals by finish time.

Let $p(j)$ be the index of the last interval that finishes at or before the start time of interval $j$.

Let:


\[
dp[j]
\]


be the maximum profit obtainable from the first $j$ sorted intervals.

There are two possibilities:

\textbf{Exclude interval }$j$\textbf{:}


\[
dp[j-1]
\]


\textbf{Include interval }$j$\textbf{:}


\[
profit_j+dp[p(j)]
\]


Therefore:


\[
\boxed{ dp[j]=\max(dp[j-1],profit_j+dp[p(j)]) }
\]


The base case is:


\[
dp[0]=0
\]


The function $p(j)$ can be computed using binary search once the intervals are sorted by finish time.

\subsection{28.4 Complexity}

\begin{itemize}
    \item Sorting intervals: $O(n\log n)$.
    \item Computing compatible predecessors by binary search: $O(n\log n)$.
    \item Filling the DP table: $O(n)$.
\end{itemize}

Total:


\[
\boxed{O(n\log n)}
\]


Auxiliary space is $O(n)$.

\textbf{Exam takeaway:} Greedy algorithms are effective when a locally optimal choice has a correctness guarantee. When selection decisions interact through profits and compatibility constraints, Dynamic Programming may be required to compare alternatives systematically.

\section{Matrix Chain Multiplication}

Matrix Chain Multiplication is a classic Dynamic Programming problem. The objective is to find the best parenthesization of a sequence of matrix multiplications so that the number of scalar multiplications is minimized.

The order of the matrices cannot be changed, because matrix multiplication is not generally commutative. However, the order of parenthesization can be changed because matrix multiplication is associative.

For example:


\[
(AB)C=A(BC)
\]


when the matrix dimensions are compatible, but the number of scalar multiplications can differ.

\subsection{29.1 Worked example}

Suppose the dimensions are:

\begin{itemize}
    \item $A_1: 10\times30$
    \item $A_2: 30\times5$
    \item $A_3: 5\times60$
\end{itemize}

There are two parenthesizations.

\subsubsection{\textbf{Option 1: }$(A_1A_2)A_3$}

First multiply $A_1$ and $A_2$:


\[
10\times30\times5=1500
\]


The resulting matrix has dimensions $10\times5$.

Then multiply the result by $A_3$:


\[
10\times5\times60=3000
\]


Total cost:


\[
1500+3000=\boxed{4500}
\]


\subsubsection{\textbf{Option 2: }$A_1(A_2A_3)$}

First multiply $A_2$ and $A_3$:


\[
30\times5\times60=9000
\]


The resulting matrix has dimensions $30\times60$.

Then multiply it by $A_1$:


\[
10\times30\times60=18000
\]


Total:


\[
9000+18000=\boxed{27000}
\]


Therefore, the first parenthesization is much cheaper:


\[
\boxed{4500\text{ scalar multiplications}}
\]


The order of matrices stays the same. Only the grouping changes.

\subsection{29.2 DP recurrence}

Suppose matrix $A_i$ has dimensions:


\[
p_{i-1}\times p_i
\]


Let $m[i][j]$ represent the minimum number of scalar multiplications required to calculate:


\[
A_iA_{i+1}\cdots A_j
\]


The base case is:


\[
m[i][i]=0
\]


A single matrix requires no multiplication.

For $i<j$, try every possible split $k$ between $i$ and $j-1$:


\[
\boxed{ m[i][j]= \min_{i\leq k<j} \left( m[i][k]+m[k+1][j]+p_{i-1}p_kp_j \right) }
\]


The first two terms are the costs of solving the two subchains. The final term is the cost of multiplying the resulting matrices.

\subsection{29.3 Complexity}

For $n$ matrices:

\begin{itemize}
    \item Time: $O(n^3)$.
    \item Space: $O(n^2)$.
\end{itemize}

The DP table contains $O(n^2)$ subproblems, and evaluating a state may examine $O(n)$ split positions.

\textbf{Exam point:} Matrix Chain Multiplication is a classic example of optimal substructure and overlapping subproblems. It is not solved optimally by simply choosing the next multiplication with the smallest immediate cost.


\newpage
\begin{center}
\vspace{0.5em}
{\color{primary}\rule{0.45\textwidth}{1.2pt}}\\[4pt]
{\Large\bfseries\color{primary} Section IV: Comparative Analysis and Selection Framework}\\[2pt]
{\color{primary}\rule{0.45\textwidth}{1.2pt}}
\vspace{0.5em}
\end{center}


These three concepts are related, but they are not interchangeable. Understanding their differences is especially important for conceptual MCQs and algorithm-selection questions.

\section{Comparative Matrix: Greedy vs. Recursion vs. Dynamic Programming}

\begin{table}[htbp]
\centering
\caption{Comprehensive Comparison: Greedy vs. Recursion vs. Dynamic Programming}
\label{tab:master_paradigms}
\begin{tabularx}{\textwidth}{p{2.8cm} X X X}
\toprule
\textbf{Evaluation Dimension} & \textbf{Greedy Algorithms} & \textbf{Recursion} & \textbf{Dynamic Programming} \\
\midrule
\textbf{Core Strategy} & Make irrevocable locally optimal choices & Reduce problem to smaller identical instances & Store and reuse solutions to overlapping subproblems \\
\textbf{Decision Reconsideration} & Never reconsiders prior decisions & Explores recursive branching paths & Considers all valid transitions via recurrence \\
\textbf{Subproblem Overlap} & Subproblems do not need to overlap & Recalculates subproblems unless cached & Explicitly exploits overlapping subproblems \\
\textbf{Key Prerequisites} & Greedy-choice property + optimal substructure & Well-defined base case + terminating reduction & Optimal substructure + overlapping subproblems \\
\textbf{Implementation} & Sorting + iterative scans & Recursive function calls & Top-down memoization or bottom-up tabulation \\
\textbf{Canonical Examples} & Activity Selection, Kruskal, Dijkstra & Factorial, Tower of Hanoi & 0/1 Knapsack, LCS, Matrix Chain \\
\textbf{Primary Limitation} & May produce suboptimal local traps & Risk of call stack overflow & High memory/state space requirements \\
\textbf{Memory Footprint} & Usually $O(1)$ or $O(n)$ auxiliary & $O(\text{depth})$ system stack frames & Cache array or multidimensional DP table \\
\bottomrule
\end{tabularx}
\end{table}

\subsubsection{Important relationship}

A recursive algorithm can be implemented using Dynamic Programming if it contains overlapping subproblems whose results can be reused.

Dynamic Programming can be implemented iteratively without recursion.

A greedy algorithm may itself use sorting, a priority queue, or another data structure. The defining characteristic is the greedy selection strategy, not the particular data structure.

\section{How to Decide Which Technique to Use}

Consider the question:

\textbf{"How can I find the minimum or maximum cost under a set of constraints?"}

The fact that a problem asks for an optimum does not automatically mean Greedy or Dynamic Programming is appropriate.

Use the following reasoning.

\subsubsection{Case 1: A local choice is provably safe}

Use a greedy algorithm when a choice can be made irrevocably without preventing an optimal solution.

Examples:

\begin{itemize}
    \item Fractional Knapsack.
    \item Ordinary Activity Selection.
    \item Minimum Spanning Tree.
    \item Dijkstra's shortest-path algorithm with nonnegative edge weights.
\end{itemize}

\subsubsection{Case 2: The problem naturally reduces to smaller instances}

Use recursion when the problem can be defined in terms of smaller instances and the resulting depth and running time are acceptable.

Examples:

\begin{itemize}
    \item Factorial.
    \item Binary search.
    \item Tree traversal.
    \item Tower of Hanoi.
\end{itemize}

Recursion alone does not prevent duplicate work. Naive Fibonacci is a classic example of a recursive algorithm with severe repeated computation.

\subsubsection{Case 3: The same subproblems recur}

Use Dynamic Programming when the problem has overlapping subproblems and the relevant answers can be stored and reused.

Examples:

\begin{itemize}
    \item Fibonacci.
    \item 0/1 Knapsack.
    \item Coin Change.
    \item Longest Common Subsequence.
    \item Matrix Chain Multiplication.
    \item Weighted Interval Scheduling.
\end{itemize}

A DP formulation typically explores multiple alternatives and uses a recurrence to select the best result.

\subsubsection{A practical checklist}

Before implementing a solution, ask:

\begin{enumerate}
    \item What exactly does the problem want to maximize, minimize, count, or decide?
    \item Can I define the solution using smaller subproblems?
    \item Do the smaller subproblems repeat?
    \item Is a locally optimal choice provably safe?
    \item What is the state, recurrence, and base case?
    \item How many states exist, and how much work does each state require?
    \item Can the memory requirement be reduced without losing needed information?
\end{enumerate}

This checklist helps identify the algorithmic structure rather than choosing a technique solely because it is familiar.


\newpage
\begin{center}
\vspace{0.5em}
{\color{primary}\rule{0.45\textwidth}{1.2pt}}\\[4pt]
{\Large\bfseries\color{primary} Section V: Critical Examination Traps and Pitfalls}\\[2pt]
{\color{primary}\rule{0.45\textwidth}{1.2pt}}
\vspace{0.5em}
\end{center}


\section{Examination Traps: Greedy Algorithms}

\textbf{1. Assuming greedy always gives the optimum.}

It does not. Ratio-based greedy is optimal for Fractional Knapsack but not for general 0/1 Knapsack.

\textbf{2. Confusing Activity Selection with Weighted Interval Scheduling.}

Selecting the earliest-finishing interval is optimal for maximizing the number of compatible intervals. It is not generally optimal when each interval has a different profit and the objective is maximum total profit.

\textbf{3. Forgetting Dijkstra's constraint.}

Standard Dijkstra's algorithm relies on nonnegative edge weights. It is not generally correct when negative-weight edges are present.

\textbf{4. Confusing minimum spanning trees with shortest paths.}

Prim's and Kruskal's algorithms minimize the total weight of a spanning tree. Dijkstra's algorithm minimizes path costs from a source. These are different optimization objectives.

\textbf{5. Assuming any locally optimal coin choice works.}

Greedy coin selection may fail for arbitrary denomination systems, as demonstrated by the denominations $\\{1,3,4\\}$ and target 6.

\section{Examination Traps: Recursion}

\textbf{1. Missing the base case.} Without a terminating condition, recursion may continue indefinitely.

\textbf{2. Not making progress.} A recursive call must reduce the problem in a way that eventually reaches the base case.

\textbf{3. Confusing recursion depth with total work.} Tower of Hanoi has $O(2^n)$ time but $O(n)$ maximum recursive depth.

\textbf{4. Ignoring overlapping subproblems.} Naive Fibonacci is exponential because it repeatedly calculates the same values.

\textbf{5. Assuming tail recursion is always optimized.} C++ does not guarantee tail-call elimination.

\textbf{6. Forgetting numeric limits.} Mathematical factorial and Fibonacci values grow rapidly; ordinary integer types hold only a finite range.

\section{Examination Traps: Dynamic Programming}

\textbf{1. Defining a vague state.} The state must precisely describe the subproblem whose answer is being stored.

\textbf{2. Using an incorrect base case.} An incorrect initialization can affect every state that depends on it.

\textbf{3. Filling the table in the wrong order.} A DP state must be calculated after its dependencies, unless the implementation specifically handles otherwise.

\textbf{4. Using the wrong loop direction in Knapsack.} In one-dimensional 0/1 Knapsack, process capacity in descending order. Ascending order can accidentally allow an item to be reused.

\textbf{5. Confusing polynomial and pseudo-polynomial complexity.} $O(nW)$ for Knapsack depends on the numeric capacity $W$, not just its encoded input length.

\textbf{6. Assuming every recursive problem requires DP.} Recursion is useful even when subproblems do not overlap, as in ordinary tree traversal or Tower of Hanoi.

\textbf{7. Assuming every optimization problem requires DP.} Some problems have a valid greedy strategy that is simpler and more efficient.

\textbf{8. Confusing the optimal value with the selected solution.} A DP table may return the optimal profit or length, but reconstructing the actual chosen items or sequence may require additional bookkeeping or backtracking through the table.


\newpage
\begin{center}
\vspace{0.5em}
{\color{primary}\rule{0.45\textwidth}{1.2pt}}\\[4pt]
{\Large\bfseries\color{primary} Section VI: Formulas, Viva Q\&A, and Quick Revision}\\[2pt]
{\color{primary}\rule{0.45\textwidth}{1.2pt}}
\vspace{0.5em}
\end{center}


\section{Master Complexity and Formula Reference Sheet}

\begin{table}[htbp]
\centering
\caption{Master Algorithmic Complexity Reference Sheet}
\label{tab:master_complexity_sheet}
\begin{tabularx}{\textwidth}{p{4.2cm} p{3.2cm} p{3.2cm} X}
\toprule
\textbf{Algorithm / Problem} & \textbf{Time Complexity} & \textbf{Auxiliary Space} & \textbf{Key Design Characteristic} \\
\midrule
Factorial (Recursive) & $O(n)$ & $O(n)$ Stack & One recursive frame per decrement \\
Fibonacci (Naive) & $\Theta(\varphi^n)$ Exponential & $O(n)$ Stack & Redundant overlapping call tree \\
Fibonacci (Memoized) & $O(n)$ & $O(n)$ Cache + Stack & Cache hits eliminate redundant work \\
Fibonacci (2 Variables) & $O(n)$ & $O(1)$ Optimized & Bottom-up sliding state window \\
Binary Search (Recursive) & $O(\log n)$ & $O(\log n)$ Stack & Logarithmic search space bisection \\
Binary Search (Iterative) & $O(\log n)$ & $O(1)$ Optimal & In-place two-pointer interval halving \\
Tower of Hanoi & $O(2^n)$ & $O(n)$ Stack & Recurrence $T(n) = 2T(n-1) + 1$ \\
Activity Selection & $O(n \log n)$ & $O(1)$ to $O(n)$ & Greedy sort by earliest finish time \\
Fractional Knapsack & $O(n \log n)$ & $O(1)$ to $O(n)$ & Greedy sort by value-to-weight ratio \\
Kruskal's MST & $O(E \log E)$ & $O(V + E)$ DSU & Edge sorting with Union-Find disjoint sets \\
Dijkstra's Algorithm & $O((V+E)\log V)$ & $O(V + E)$ Adjacency & Min-heap; strictly nonnegative edge weights \\
0/1 Knapsack DP & $O(nW)$ & $O(W)$ Optimized & Pseudo-polynomial; reverse inner loop \\
Coin Change (Min Coins) & $O(Am)$ & $O(A)$ Table & Amount $A$, $m$ coin denominations \\
Longest Common Subseq. & $O(mn)$ & $O(mn)$ or $O(\min(m,n))$ & 2D grid character matching recurrence \\
Matrix Chain Mult. & $O(n^3)$ & $O(n^2)$ Table & Catalan bracketings, parenthesization DP \\
Weighted Interval Sched. & $O(n \log n)$ & $O(n)$ Array & Predecessor search $p(j)$ + 1D DP \\
\bottomrule
\end{tabularx}
\end{table}

The table assumes the standard implementations described in this chapter. In particular, graph algorithms' storage figures include graph representation when indicated, and DP space can vary if only an optimal value is required instead of reconstructing the solution.


\section{Likely Viva Questions: Greedy Algorithms}

\begin{quote}
\textbf{Question 1: What are the two important properties associated with greedy optimization algorithms?}\\[4pt]
Answer: The greedy-choice property and optimal substructure.
\end{quote}
\begin{quote}
\textbf{Question 2: Which Knapsack problem is optimally solved using the value-to-weight ratio greedy strategy?}\\[4pt]
Answer: Fractional Knapsack.
\end{quote}
\begin{quote}
\textbf{Question 3: Why is a ratio-based greedy algorithm not generally optimal for 0/1 Knapsack?}\\[4pt]
Answer: Items cannot be divided. A locally attractive item may prevent selecting a combination of other items with a higher total value.
\end{quote}
\begin{quote}
\textbf{Question 4: Which edge does Kruskal's algorithm consider first?}\\[4pt]
Answer: The minimum-weight edge, provided adding it does not form a cycle in the selected forest.
\end{quote}
\begin{quote}
\textbf{Question 5: What is the key greedy decision in Prim's algorithm?}\\[4pt]
Answer: Add a minimum-weight edge connecting the current tree to a vertex outside it.
\end{quote}
\begin{quote}
\textbf{Question 6: Under what conditions is Dijkstra's algorithm appropriate?}\\[4pt]
Answer: For shortest paths in a weighted graph with nonnegative edge weights.
\end{quote}
\begin{quote}
\textbf{Question 7: What is the greedy strategy used by Huffman Coding?}\\[4pt]
Answer: Repeatedly combine the two nodes with the smallest frequencies.
\end{quote}
\begin{quote}
\textbf{Question 8: Does a greedy coin-change algorithm always minimize the number of coins?}\\[4pt]
Answer: No. Its correctness depends on the denomination system.
\end{quote}
\section{Likely Viva Questions: Recursion}

\begin{quote}
\textbf{Question 9: What are the two essential parts of a recursive function?}\\[4pt]
Answer: A base case and a recursive case.
\end{quote}
\begin{quote}
\textbf{Question 10: What causes excessive recursion depth?}\\[4pt]
Answer: A large number of active recursive calls, potentially exhausting call-stack space.
\end{quote}
\begin{quote}
\textbf{Question 11: What is the recurrence relation for factorial?}\\[4pt]
Answer: 
\[
T(n)=T(n-1)+O(1)
\]

\end{quote}
\begin{quote}
\textbf{Question 12: What is the minimum number of moves for Tower of Hanoi with }$n$\textbf{ disks?}\\[4pt]
Answer: 
\[
2^n-1
\]

\end{quote}
\begin{quote}
\textbf{Question 13: What is the time complexity of naive recursive Fibonacci?}\\[4pt]
Answer: Exponential, with $\Theta(\varphi^n)$ growth under the usual model.
\end{quote}
\begin{quote}
\textbf{Question 14: What is tail recursion?}\\[4pt]
Answer: Recursion in which the recursive call is the final operation performed by the function.
\end{quote}
\begin{quote}
\textbf{Question 15: What are the time and auxiliary-space complexities of recursive binary search?}\\[4pt]
Answer: $O(\log n)$ time and $O(\log n)$ call-stack space.
\end{quote}
\section{Likely Viva Questions: Dynamic Programming}

\begin{quote}
\textbf{Question 16: What are the two main properties commonly used to identify DP problems?}\\[4pt]
Answer: Optimal substructure and overlapping subproblems.
\end{quote}
\begin{quote}
\textbf{Question 17: What is memoization?}\\[4pt]
Answer: A top-down approach that caches answers to previously solved subproblems.
\end{quote}
\begin{quote}
\textbf{Question 18: What is tabulation?}\\[4pt]
Answer: A bottom-up approach that fills a table in a valid dependency order.
\end{quote}
\begin{quote}
\textbf{Question 19: Why is DP Fibonacci faster than naive recursive Fibonacci?}\\[4pt]
Answer: It computes each required Fibonacci state once and reuses the stored result.
\end{quote}
\begin{quote}
\textbf{Question 20: What is the time complexity of the 0/1 Knapsack DP using a one-dimensional table?}\\[4pt]
Answer: $O(nW)$, where $n$ is the number of items and $W$ is the capacity.
\end{quote}
\begin{quote}
\textbf{Question 21: Why is }$O(nW)$\textbf{ called pseudo-polynomial?}\\[4pt]
Answer: It depends on the numerical value $W$, which may be exponentially larger than the number of bits used to represent $W$.
\end{quote}
\begin{quote}
\textbf{Question 22: What is the time complexity of the standard LCS DP algorithm?}\\[4pt]
Answer: $O(mn)$, where $m$ and $n$ are the lengths of the two sequences.
\end{quote}
\begin{quote}
\textbf{Question 23: What is the main difference between memoization and tabulation?}\\[4pt]
Answer: Memoization is top-down and solves states as needed; tabulation is bottom-up and fills states in a planned order.
\end{quote}
\begin{quote}
\textbf{Question 24: Does every problem that can be solved recursively require Dynamic Programming?}\\[4pt]
Answer: No. DP is especially useful when subproblems overlap and their results can be reused.
\end{quote}
\section{Final Revision Summary}

The essential distinctions to remember are:

\begin{itemize}
    \item \textbf{Greedy:} Makes locally optimal choices. Correctness depends on a valid greedy-choice argument. Examples include Activity Selection, Fractional Knapsack, Kruskal's MST, Prim's MST, Dijkstra under nonnegative weights, and Huffman Coding.
    \item \textbf{Recursion:} Solves a problem through smaller instances of itself. It requires a base case and a reduction that guarantees termination. Its performance depends on the recurrence and the number of repeated calculations.
    \item \textbf{Dynamic Programming:} Stores and reuses answers to overlapping subproblems. It can be implemented using top-down memoization or bottom-up tabulation.
\end{itemize}

When solving an examination question, identify the problem's state or decision structure first, then select the appropriate algorithmic method. Write the recurrence or greedy rule explicitly, justify why it is correct, and derive time and space complexity from the actual operations.

For a final practice session, prioritize four comparisons: \textbf{Fractional versus 0/1 Knapsack, naive recursion versus memoized Fibonacci, greedy Activity Selection versus Weighted Interval Scheduling, and recursive versus iterative Binary Search.} These examples demonstrate the essential differences between the three paradigms.


\begin{center}
\vspace{1em}
{\color{primary}\rule{0.5\textwidth}{1.5pt}}\\[6pt]
{\Large\bfseries\color{primary} Part III: High-Value Programming Problems and Solutions}\\[3pt]
{\color{primary}\rule{0.5\textwidth}{1.5pt}}
\vspace{1em}
\end{center}

This chapter brings together essential practical algorithm problems and complete implementations frequently tested in Bangladesh government technical recruitment examinations (including Assistant Programmer, Assistant Director (ICT), Assistant Network Engineer, Assistant Maintenance Engineer, and Assistant Database Administrator). It synthesizes mathematical programming, number theory, numerical series, array algorithms, recursion, Dynamic Programming, greedy heuristics, and weighted-graph algorithms.

Each problem includes a precise specification, core mathematical and algorithmic invariants, complete C++17 implementations, execution traces, formal time and space complexity derivations, and critical examination takeaways.


\begin{center}
\vspace{0.5em}
{\color{primary}\rule{0.45\textwidth}{1.2pt}}\\[4pt]
{\Large\bfseries\color{primary} Section I: Number Theory and Mathematical Programming}\\[2pt]
{\color{primary}\rule{0.45\textwidth}{1.2pt}}
\vspace{0.5em}
\end{center}

\section{Number Theory and Mathematical Programming}


\subsection{Problem 1. Calculate the LCM of two integers}

\textbf{Question:} Write a program to calculate the Least Common Multiple (LCM) of two integers.

The \textbf{Least Common Multiple} of two integers is the smallest positive integer divisible by both numbers, for positive nonzero inputs.

For example:


\[
\operatorname{LCM}(12,18)=36
\]


Multiples of 12 are $12,24,36,48,\ldots$, while multiples of 18 are $18,36,54,\ldots$. The first positive common multiple is 36.

\subsubsection{Mathematical relationship}

The GCD and LCM satisfy:


\[
\boxed{\operatorname{LCM}(a,b)=\frac{|a|}{\operatorname{GCD}(a,b)}|b|}
\]


This formula applies when $a$ and $b$ are not both zero. When either operand is zero, the LCM is conventionally zero.

The formula should be calculated as:


\[
\boxed{\operatorname{LCM}(a,b)=\left(\frac{|a|}{\operatorname{GCD}(a,b)}\right)|b|}
\]


rather than multiplying $a$ and $b$ before dividing, because dividing first reduces the risk of integer overflow.

\subsubsection{C++ implementation}

\begin{lstlisting}[style=cppstyle]
#include <iostream>
#include <numeric>
using namespace std;

long long calculateLCM(long long a, long long b) {
    if (a == 0 || b == 0) {
        return 0;
    }

    long long g = gcd(a, b);

    return (a / g) * b;
}

int main() {
    long long a, b;
    cin >> a >> b;

    cout << "LCM = " << calculateLCM(a, b) << '\n';

    return 0;
}
\end{lstlisting}

Input:

\begin{lstlisting}[style=pseudocode]
12 18
\end{lstlisting}

Output:

\begin{lstlisting}[style=pseudocode]
LCM = 36
\end{lstlisting}

This implementation assumes nonnegative inputs and a result representable in \texttt{long long}. For general signed values, normalize the signs carefully and account for the fact that the absolute value of the minimum representable signed integer may not fit in the same type.

\subsubsection{Complexity}

Euclid's GCD algorithm takes $O(\log\min(a,b))$ time for positive inputs. The LCM calculation then requires a constant number of arithmetic operations.

Therefore:

\begin{itemize}
    \item \textbf{Time:} $O(\log\min(a,b))$.
    \item \textbf{Auxiliary space:} $O(1)$ for the iterative GCD implementation used internally by the standard library.
\end{itemize}

\textbf{Exam point:} LCM is efficiently calculated using GCD; generating multiples one by one is unnecessary.

\subsection{Problem 2. Calculate GCD using Euclid's algorithm}

\textbf{Question:} Find the Greatest Common Divisor (GCD) of two positive integers.

The GCD is the largest positive integer that divides both inputs without a remainder.

For example:


\[
\operatorname{GCD}(48,18)=6
\]


\subsubsection{Algorithm}

Euclid's algorithm uses the identity:


\[
\boxed{\gcd(a,b)=\gcd(b,a\bmod b)}
\]


until the second operand becomes zero. The remaining first operand is the GCD.

\subsubsection{Step-by-step calculation}

Find $\gcd(48,18)$.


\[
48\bmod18=12
\]


Thus:


\[
\gcd(48,18)=\gcd(18,12)
\]


Next:


\[
18\bmod12=6
\]


So:


\[
\gcd(18,12)=\gcd(12,6)
\]


Finally:


\[
12\bmod6=0
\]


Therefore:


\[
\boxed{\gcd(48,18)=6}
\]


\subsubsection{C++ implementation}

\begin{lstlisting}[style=cppstyle]
#include <iostream>
using namespace std;

long long gcdEuclid(long long a, long long b) {
    while (b != 0) {
        long long remainder = a % b;
        a = b;
        b = remainder;
    }

    return a;
}

int main() {
    long long a, b;
    cin >> a >> b;

    cout << "GCD = " << gcdEuclid(a, b) << '\n';

    return 0;
}
\end{lstlisting}

For positive inputs, the algorithm works by repeatedly replacing the larger problem with a smaller remainder problem.

\textbf{Complexity:} $O(\log\min(a,b))$ time and $O(1)$ auxiliary space.

\textbf{Related examination question:} How can the GCD be used to determine whether two integers are relatively prime?

Two positive integers are relatively prime, or coprime, if:


\[
\boxed{\gcd(a,b)=1}
\]


\subsection{Problem 3. Determine whether a number is prime}

\textbf{Question:} Write a program to determine whether a positive integer is a prime number.

A \textbf{prime number} is an integer greater than 1 that has exactly two positive divisors: 1 and itself.

Examples: $2,3,5,7,11,13,17,19$.

The numbers 0 and 1 are not prime. A number greater than 1 that is not prime is called composite.

\subsubsection{Efficient algorithm}

To test an integer $n$:

\begin{enumerate}
    \item If $n<2$, it is not prime.
    \item Check whether any integer from 2 through $\sqrt n$ divides it.
    \item If such a divisor exists, it is composite.
    \item Otherwise, it is prime.
\end{enumerate}

Why only test through $\sqrt n$?

If $n=a\times b$ is composite, at least one of its factors must be less than or equal to $\sqrt n$. If both factors were greater than $\sqrt n$, their product would exceed $n$.

\subsubsection{C++ implementation}

\begin{lstlisting}[style=cppstyle]
#include <iostream>
using namespace std;

bool isPrime(long long n) {
    if (n < 2) {
        return false;
    }

    if (n == 2) {
        return true;
    }

    if (n % 2 == 0) {
        return false;
    }

    for (long long i = 3; i <= n / i; i += 2) {
        if (n % i == 0) {
            return false;
        }
    }

    return true;
}

int main() {
    long long n;
    cin >> n;

    cout << (isPrime(n) ? "Prime" : "Not prime") << '\n';

    return 0;
}
\end{lstlisting}

The condition \texttt{i <= n / i} avoids the possible overflow of \texttt{i * i} when \texttt{n} is a large integer.

Example:

\begin{lstlisting}[style=pseudocode]
Input: 97
Output: Prime
\end{lstlisting}

\subsubsection{Complexity}

The trial-division method tests approximately half the integers up to $\sqrt n$, so:

\begin{itemize}
    \item \textbf{Time:} $O(\sqrt n)$.
    \item \textbf{Auxiliary space:} $O(1)$.
\end{itemize}

The check for even numbers reduces unnecessary work but does not change the asymptotic complexity.

\subsection{\textbf{Problem 4. Generate all prime numbers up to }$N$}

\textbf{Question:} Print every prime number between 2 and a given integer $N$.

For example, if $N=30$, the output should be:

\begin{lstlisting}[style=pseudocode]
2 3 5 7 11 13 17 19 23 29
\end{lstlisting}

Testing each number separately with trial division works, but there is a more efficient method for generating many primes.

\subsubsection{The Sieve of Eratosthenes}

The \textbf{Sieve of Eratosthenes} finds all primes up to $N$ by marking multiples of each discovered prime as composite.

Algorithm:

\begin{enumerate}
    \item Initially mark every integer from 2 to $N$ as potentially prime.
    \item Start with $p=2$.
    \item Mark $2p,3p,4p,\ldots$ as composite.
    \item Find the next unmarked number.
    \item Repeat until $p^2>N$.
    \item The remaining unmarked integers are prime.
\end{enumerate}

Marking begins at $p^2$, because smaller multiples of $p$ have already been processed as multiples of smaller primes.

\subsubsection{C++ implementation}

\begin{lstlisting}[style=cppstyle]
#include <iostream>
#include <vector>
using namespace std;

vector<int> generatePrimes(int n) {
    vector<int> primes;

    if (n < 2) {
        return primes;
    }

    vector<bool> isPrime(n + 1, true);
    isPrime[0] = false;
    isPrime[1] = false;

    for (int p = 2; p <= n / p; ++p) {
        if (isPrime[p]) {
            for (long long multiple = 1LL * p * p;
                 multiple <= n;
                 multiple += p) {
                isPrime[static_cast<int>(multiple)] = false;
            }
        }
    }

    for (int i = 2; i <= n; ++i) {
        if (isPrime[i]) {
            primes.push_back(i);
        }
    }

    return primes;
}

int main() {
    int n;
    cin >> n;

    vector<int> primes = generatePrimes(n);

    for (int p : primes) {
        cout << p << " ";
    }

    cout << '\n';
    return 0;
}
\end{lstlisting}

\subsubsection{Complexity}

\begin{itemize}
    \item \textbf{Time:} $O(N\log\log N)$.
    \item \textbf{Space:} $O(N)$.
\end{itemize}

This is generally more efficient than testing every integer independently when a large collection of primes is needed.

\textbf{Exam distinction:} Prime testing for one number is typically done with trial division up to its square root. Generating all primes up to a limit is commonly done using the Sieve of Eratosthenes.

\subsection{Problem 5. Find the prime factorization of an integer}

\textbf{Question:} Express a positive integer as a product of prime factors.

For example:


\[
360=2^3\times3^2\times5
\]


\subsubsection{Algorithm}

\begin{enumerate}
    \item Start with the smallest prime candidate, 2.
    \item While it divides the number, record the factor and divide the number by it.
    \item Move to the next possible factor.
    \item Continue while the factor does not exceed the square root of the remaining value.
    \item If the remaining value is greater than 1, it is itself a prime factor.
\end{enumerate}

\subsubsection{C++ implementation}

\begin{lstlisting}[style=cppstyle]
#include <iostream>
using namespace std;

void primeFactorization(long long n) {
    if (n < 2) {
        cout << "Prime factorization requires n >= 2\n";
        return;
    }

    for (long long p = 2; p <= n / p; ++p) {
        while (n % p == 0) {
            cout << p << " ";
            n /= p;
        }
    }

    if (n > 1) {
        cout << n;
    }

    cout << '\n';
}

int main() {
    long long n;
    cin >> n;

    primeFactorization(n);

    return 0;
}
\end{lstlisting}

Input:

\begin{lstlisting}[style=pseudocode]
360
\end{lstlisting}

Output:

\begin{lstlisting}[style=pseudocode]
2 2 2 3 3 5
\end{lstlisting}

\subsubsection{Complexity}

\begin{itemize}
    \item \textbf{Time:} $O(\sqrt n)$ in the worst case under trial division.
    \item \textbf{Auxiliary space:} $O(1)$, excluding printed output.
\end{itemize}

The factors are printed in nondecreasing order.

\textbf{Applications:} Finding divisors, simplifying fractions, determining whether an integer is square-free, and calculating GCD and LCM using prime exponents.


\newpage
\begin{center}
\vspace{0.5em}
{\color{primary}\rule{0.45\textwidth}{1.2pt}}\\[4pt]
{\Large\bfseries\color{primary} Section II: Numerical Series and Mathematical Loops}\\[2pt]
{\color{primary}\rule{0.45\textwidth}{1.2pt}}
\vspace{0.5em}
\end{center}

\section{Numerical Series and Mathematical Loops}


\subsection{\textbf{Problem 6. Calculate the sum of even numbers from 1 to }$N$}

\textbf{Question:} Calculate:


\[
2+4+6+\cdots
\]


up to the largest even number not exceeding $N$.

For example, if $N=10$:


\[
2+4+6+8+10=30
\]


\subsubsection{Mathematical solution}

Let $k$ be the number of even integers in the range from 1 to $N$.

Then:


\[
k=\left\lfloor\frac N2\right\rfloor
\]


The sum is:


\[
2+4+\cdots+2k =2(1+2+\cdots+k)
\]


Using:


\[
1+2+\cdots+k=\frac{k(k+1)}2
\]


we obtain:


\[
\boxed{S_{\text{even}}=k(k+1)}
\]


where:


\[
\boxed{k=\left\lfloor\frac N2\right\rfloor}
\]


Therefore:


\[
\boxed{ S_{\text{even}} = \left\lfloor\frac N2\right\rfloor \left(\left\lfloor\frac N2\right\rfloor+1\right) }
\]


\subsubsection{C++ implementation}

\begin{lstlisting}[style=cppstyle]
#include <iostream>
using namespace std;

long long sumEvenNumbers(long long n) {
    long long k = n / 2;
    return k * (k + 1);
}

int main() {
    long long n;
    cin >> n;

    if (n < 0) {
        cout << "N must be nonnegative\n";
        return 0;
    }

    cout << sumEvenNumbers(n) << '\n';

    return 0;
}
\end{lstlisting}

Input:

\begin{lstlisting}[style=pseudocode]
10
\end{lstlisting}

Output:

\begin{lstlisting}[style=pseudocode]
30
\end{lstlisting}

For $N=9$, the largest even number is 8, so:


\[
2+4+6+8=20
\]


The formula correctly gives:


\[
k=\lfloor9/2\rfloor=4
\]



\[
S=4(5)=20
\]


\subsubsection{Complexity}

\begin{itemize}
    \item \textbf{Formula-based solution:} $O(1)$ time and $O(1)$ auxiliary space.
    \item \textbf{Loop-based solution:} $O(N)$ time and $O(1)$ auxiliary space.
\end{itemize}

For very large values of $N$, the result may exceed the range of \texttt{long long}. A formula can be constant-time without necessarily being safe for arbitrarily large integers.

\textbf{Exam point:} When a numerical series has a closed-form formula, it can often replace a loop and reduce the running time from linear to constant.

\subsection{Problem 7. Calculate the sum of the first **$N$** natural numbers}

\textbf{Question:} Find:


\[
1+2+3+\cdots+N
\]


The formula is:


\[
\boxed{S=\frac{N(N+1)}2}
\]


\subsubsection{Example}

For $N=100$:


\[
S=\frac{100\times101}{2}
\]



\[
\boxed{S=5050}
\]


\subsubsection{C++ implementation}

\begin{lstlisting}[style=cppstyle]
#include <iostream>
using namespace std;

long long sumNaturalNumbers(long long n) {
    return n * (n + 1) / 2;
}

int main() {
    long long n;
    cin >> n;

    if (n < 0) {
        cout << "N must be nonnegative\n";
        return 0;
    }

    cout << sumNaturalNumbers(n) << '\n';
    return 0;
}
\end{lstlisting}

\textbf{Complexity:} $O(1)$ time and $O(1)$ auxiliary space, assuming the intermediate arithmetic fits the selected type.

A loop that adds every integer from 1 to $N$ would take $O(N)$ time.

\subsubsection{Related formula: Sum of the first **$N$** odd numbers}

The first $N$ positive odd numbers are:


\[
1,3,5,\ldots,(2N-1)
\]


Their sum is:


\[
\boxed{S_{\text{odd}}=N^2}
\]


For example:


\[
1+3+5+7+9=5^2=\boxed{25}
\]


\subsection{Problem 8. Calculate the sum of squares}

\textbf{Question:} Calculate:


\[
1^2+2^2+3^2+\cdots+N^2
\]


The formula is:


\[
\boxed{S=\frac{N(N+1)(2N+1)}6}
\]


\subsubsection{Example}

For $N=5$:


\[
S=1+4+9+16+25=55
\]


Using the formula:


\[
S=\frac{5(6)(11)}6=\boxed{55}
\]


\subsubsection{C++ implementation}

\begin{lstlisting}[style=cppstyle]
#include <iostream>
using namespace std;

long long sumSquares(long long n) {
    return n * (n + 1) * (2 * n + 1) / 6;
}

int main() {
    long long n;
    cin >> n;

    if (n < 0) {
        cout << "N must be nonnegative\n";
        return 0;
    }

    cout << sumSquares(n) << '\n';
    return 0;
}
\end{lstlisting}

\textbf{Complexity:} $O(1)$ time and $O(1)$ auxiliary space, provided the arithmetic does not overflow.

For large inputs, intermediate multiplication may overflow even when the final result fits. A more robust implementation can divide factors by 2 and 3 before multiplication or use a wider/arbitrary-precision integer type.

\subsection{Problem 9. Generate an arithmetic progression}

An \textbf{Arithmetic Progression (AP)} is a numerical sequence in which the difference between consecutive terms is constant.

For example:


\[
3,7,11,15,19,\ldots
\]


The first term is $a=3$, and the common difference is $d=4$.

The $n$-th term is:


\[
\boxed{a_n=a+(n-1)d}
\]


The sum of the first $n$ terms is:


\[
\boxed{S_n=\frac n2\left[2a+(n-1)d\right]}
\]


\subsubsection{Example}

Find the 20th term and the sum of the first 20 terms of:


\[
3,7,11,15,\ldots
\]


Given:


\[
a=3,\qquad d=4,\qquad n=20
\]


The 20th term is:


\[
a_{20}=3+(20-1)4
\]



\[
a_{20}=3+76=\boxed{79}
\]


The sum is:


\[
S_{20}=\frac{20}{2}[2(3)+19(4)]
\]



\[
S_{20}=10(6+76)
\]



\[
\boxed{S_{20}=820}
\]


\subsubsection{C++ implementation}

\begin{lstlisting}[style=cppstyle]
#include <iostream>
using namespace std;

int main() {
    long long a, d, n;
    cin >> a >> d >> n;

    if (n <= 0) {
        cout << "Number of terms must be positive\n";
        return 0;
    }

    long long nth = a + (n - 1) * d;
    long long sum = n * (2 * a + (n - 1) * d) / 2;

    cout << "Nth term: " << nth << '\n';
    cout << "Sum: " << sum << '\n';

    return 0;
}
\end{lstlisting}

\textbf{Complexity:} $O(1)$ time and $O(1)$ auxiliary space, assuming intermediate calculations fit the integer type.

\textbf{Exam variation:} If a question asks for all terms, printing $n$ values necessarily takes at least $O(n)$ time, even though the $n$-th term can be calculated in $O(1)$.

\subsection{Problem 10. Generate a geometric progression}

A \textbf{Geometric Progression (GP)} is a numerical sequence in which each term is obtained by multiplying the previous term by a constant ratio $r$.

For example:


\[
2,6,18,54,162,\ldots
\]


Here:


\[
a=2,\qquad r=3
\]


The $n$-th term is:


\[
\boxed{a_n=ar^{n-1}}
\]


For $r\neq1$, the sum of the first $n$ terms is:


\[
\boxed{S_n=a\frac{r^n-1}{r-1}}
\]


For $r=1$, every term is $a$, so:


\[
\boxed{S_n=na}
\]


\subsubsection{Example}

Find the 5th term and the sum of the first 5 terms of:


\[
2,6,18,54,162
\]


The fifth term is:


\[
a_5=2(3)^4=\boxed{162}
\]


The sum is:


\[
S_5=2\frac{3^5-1}{3-1}
\]



\[
S_5=2\frac{243-1}{2}
\]



\[
\boxed{S_5=242}
\]


\subsubsection{Exam note}

For integer sequences, powers may grow very rapidly and overflow ordinary integer types. The mathematical formula is valid, but the implementation must consider numerical limits.


\newpage
\begin{center}
\vspace{0.5em}
{\color{primary}\rule{0.45\textwidth}{1.2pt}}\\[4pt]
{\Large\bfseries\color{primary} Section III: Pair-Sum and Two-Element Selection}\\[2pt]
{\color{primary}\rule{0.45\textwidth}{1.2pt}}
\vspace{0.5em}
\end{center}

\section{Pair-Sum and Two-Element Selection Problems}


\subsection{Problem 11. Find two elements whose sum equals a target}

\textbf{Question:} Given an unsorted array and a target $T$, find two elements at distinct positions whose sum is $T$.

Input:

\begin{lstlisting}[style=pseudocode]
Array:  [8, 3, 2, 7, 11, 6]
Target: 9
\end{lstlisting}

One valid answer is:


\[
2+7=9
\]


Another valid answer is:


\[
3+6=9
\]


This problem is often called \textbf{Two Sum} or the pair-sum problem.

\subsubsection{Approach A: Brute force}

Test every pair of positions $i<j$.

\begin{lstlisting}[style=pseudocode]
For i from 0 to n-2:
    For j from i+1 to n-1:
        If A[i] + A[j] == target:
            Return the pair
\end{lstlisting}

\textbf{Time:} $O(n^2)$.

\textbf{Auxiliary space:} $O(1)$.

\subsubsection{Approach B: Sorting and two pointers}

Sort the array, then use a pointer at each end.

\begin{itemize}
    \item If the sum is too small, increment the left pointer.
    \item If the sum is too large, decrement the right pointer.
    \item If the sum matches the target, return the pair.
\end{itemize}

The total time is $O(n\log n)$, dominated by sorting.

\subsubsection{C++ implementation with original indices}

This version returns the original indices rather than merely the values. Sorting pairs preserves the correspondence between each value and its original position.

\begin{lstlisting}[style=cppstyle]
#include <algorithm>
#include <iostream>
#include <utility>
#include <vector>
using namespace std;

bool twoSum(
    const vector<int>& a,
    long long target,
    pair<int, int>& answer
) {
    vector<pair<int, int>> items;

    for (int i = 0; i < static_cast<int>(a.size()); ++i) {
        items.push_back({a[i], i});
    }

    sort(items.begin(), items.end());

    int left = 0;
    int right = static_cast<int>(items.size()) - 1;

    while (left < right) {
        long long sum =
            static_cast<long long>(items[left].first)
            + items[right].first;

        if (sum == target) {
            answer = {items[left].second, items[right].second};
            return true;
        }

        if (sum < target) {
            ++left;
        } else {
            --right;
        }
    }

    return false;
}

int main() {
    vector<int> a = {8, 3, 2, 7, 11, 6};
    long long target = 9;

    pair<int, int> answer;

    if (twoSum(a, target, answer)) {
        cout << "Original indices: "
             << answer.first << ", " << answer.second << '\n';
        cout << "Values: "
             << a[answer.first] << ", " << a[answer.second] << '\n';
    } else {
        cout << "No pair found\n";
    }

    return 0;
}
\end{lstlisting}

The returned pair of indices may differ from another valid pair because several solutions can exist.

\textbf{Complexity:} $O(n\log n)$ time and $O(n)$ auxiliary space for the sorted copy and stored indices.

\textbf{Important:} The condition \texttt{left < right} ensures that the same array position cannot be selected twice. Equal values at different positions can still form a valid pair.

\subsubsection{Approach C: Hash set}

A hash-based approach scans the array once. For each value $x$, check whether $T-x$ has already appeared. If it has, a valid pair is found; otherwise, record $x$.

This can achieve expected $O(n)$ time and $O(n)$ space under suitable hash-table performance assumptions.

\textbf{Exam comparison:}

\begin{table}[htbp]
\centering
\caption{Algorithmic Approaches to the Pair-Sum (Two Sum) Problem}
\label{tab:two_sum_comparison}
\begin{tabularx}{\textwidth}{p{4.2cm} X X}
\toprule
\textbf{Method} & \textbf{Time Complexity} & \textbf{Auxiliary Space} \\
\midrule
Brute Force & $O(n^2)$ & $O(1)$ \\
Sorting + Two Pointers & $O(n \log n)$ & $O(n)$ when preserving original indices \\
Hash Set / Map & Expected $O(n)$ & $O(n)$ \\
\bottomrule
\end{tabularx}
\end{table}

For a question specifically requesting an $O(n\log n)$ solution, sorting and two pointers is the standard method.

\subsection{Problem 12. Find the two largest distinct elements in an array}

\textbf{Question:} Find the largest and second-largest distinct elements without sorting.

Input:

\begin{lstlisting}[style=pseudocode]
10 25 8 25 17 30 30
\end{lstlisting}

Output:

\begin{lstlisting}[style=pseudocode]
Largest: 30
Second largest: 25
\end{lstlisting}

The second-largest value must be distinct from the largest. A duplicate \texttt{30} must not be counted as the second-largest distinct value.

\subsubsection{Algorithm}

Maintain two values:

\begin{itemize}
    \item \texttt{largest}: maximum encountered so far.
    \item \texttt{secondLargest}: maximum encountered value that is strictly smaller than \texttt{largest}.
\end{itemize}

For every element:

\begin{enumerate}
    \item If it is larger than \texttt{largest}, move the previous largest into \texttt{secondLargest}, then update \texttt{largest}.
    \item Otherwise, if it is smaller than \texttt{largest} but greater than \texttt{secondLargest}, update \texttt{secondLargest}.
    \item Ignore duplicates of \texttt{largest}.
\end{enumerate}

\subsubsection{C++ implementation}

\begin{lstlisting}[style=cppstyle]
#include <iostream>
#include <vector>
#include <optional>
using namespace std;

int main() {
    vector<int> a = {10, 25, 8, 25, 17, 30, 30};

    optional<int> largest;
    optional<int> secondLargest;

    for (int x : a) {
        if (!largest || x > *largest) {
            secondLargest = largest;
            largest = x;
        } else if (x < *largest &&
                   (!secondLargest || x > *secondLargest)) {
            secondLargest = x;
        }
    }

    if (!largest) {
        cout << "Array is empty\n";
    } else if (!secondLargest) {
        cout << "No second-largest distinct element\n";
    } else {
        cout << "Largest: " << *largest << '\n';
        cout << "Second largest: " << *secondLargest << '\n';
    }

    return 0;
}
\end{lstlisting}

This uses \texttt{std::optional}, available in C++17, to represent whether a value has been found.

\subsubsection{Complexity}

\begin{itemize}
    \item \textbf{Time:} $O(n)$, one pass through the array.
    \item \textbf{Auxiliary space:} $O(1)$.
\end{itemize}

\textbf{Exam variation:} Find the two smallest distinct elements, or find the two largest elements when duplicates are permitted. The latter problem allows the largest value to be selected twice if it appears at two distinct positions.

\subsection{Problem 13. Select two elements with the minimum absolute difference}

\textbf{Question:} Given an array of integers, find two distinct elements with the smallest absolute difference.

Example:

\begin{lstlisting}[style=pseudocode]
Array: [10, 3, 20, 12, 8]
\end{lstlisting}

Sort the array:

\texttt{[3, 8, 10, 12, 20]}

Compare adjacent differences:

\begin{itemize}
    \item $8-3=5$
    \item $10-8=2$
    \item $12-10=2$
    \item $20-12=8$
\end{itemize}

The minimum difference is:


\[
\boxed{2}
\]


The pairs \texttt{(8, 10)} and \texttt{(10, 12)} both achieve this minimum.

\subsubsection{Why do adjacent elements suffice?}

After sorting, if two values are not adjacent, there is at least one value between them. Their difference is at least as large as one of the adjacent differences within that interval. Therefore, a minimum-difference pair can be found among adjacent sorted values.

\subsubsection{C++ implementation}

\begin{lstlisting}[style=cppstyle]
#include <algorithm>
#include <cstdlib>
#include <iostream>
#include <limits>
#include <vector>
using namespace std;

int main() {
    vector<int> a = {10, 3, 20, 12, 8};

    if (a.size() < 2) {
        cout << "At least two elements are required\n";
        return 0;
    }

    sort(a.begin(), a.end());

    long long bestDifference = LLONG_MAX;
    int first = a[0];
    int second = a[1];

    for (size_t i = 1; i < a.size(); ++i) {
        long long difference =
            static_cast<long long>(a[i]) - a[i - 1];

        if (difference < bestDifference) {
            bestDifference = difference;
            first = a[i - 1];
            second = a[i];
        }
    }

    cout << "Pair: " << first << ", " << second << '\n';
    cout << "Minimum difference: "
         << bestDifference << '\n';

    return 0;
}
\end{lstlisting}

Add \texttt{\#include <climits>} for \texttt{LLONG\_MAX} in this exact implementation.

\subsubsection{Complexity}

Sorting costs $O(n\log n)$, and scanning adjacent elements costs $O(n)$.

Therefore:


\[
\boxed{O(n\log n)}
\]


The sorting step dominates the total complexity.


\newpage
\begin{center}
\vspace{0.5em}
{\color{primary}\rule{0.45\textwidth}{1.2pt}}\\[4pt]
{\Large\bfseries\color{primary} Section IV: Weighted Graphs and Negative Cycles}\\[2pt]
{\color{primary}\rule{0.45\textwidth}{1.2pt}}
\vspace{0.5em}
\end{center}

\section{Weighted Graphs and Negative-Cycle Detection}


\subsection{Problem 14. Determine whether a weighted directed graph contains a negative cycle}

This is a high-value graph problem because it combines graph theory, shortest-path algorithms, edge relaxation, and careful handling of disconnected components.

\subsubsection{Definition}

A \textbf{negative cycle} is a directed cycle whose total edge weight is negative.

For example, consider:


\[
A\rightarrow B \quad (2)
\]



\[
B\rightarrow C \quad (-5)
\]



\[
C\rightarrow A \quad (1)
\]


The total cycle weight is:


\[
2+(-5)+1=-2
\]


Since the total is negative:


\[
\boxed{\text{The graph contains a negative cycle}}
\]


\subsubsection{Why is a negative cycle important?}

If a reachable negative cycle exists, repeatedly traversing that cycle can reduce the path cost without bound.

Consequently, a finite shortest-path distance does not exist for vertices reachable from that cycle, under the usual shortest-walk definition.

This is why standard shortest-path algorithms must account for negative cycles when negative edge weights are permitted.

\subsection{Bellman--Ford Algorithm and Edge Relaxation}

The \textbf{Bellman--Ford algorithm} computes shortest-path distances from a source in a weighted directed graph and can detect negative cycles reachable from that source.

It uses repeated \textbf{edge relaxation}.

For an edge $u\rightarrow v$ with weight $w$, if the path through $u$ improves the distance to $v$, update:


\[
\boxed{ \text{if }dist[u]+w<dist[v], \quad dist[v]=dist[u]+w }
\]


The algorithm repeats this process.

In a graph with $V$ vertices and no reachable negative cycle, a shortest path can be chosen without repeating vertices and therefore uses at most $V-1$ edges. Consequently, $V-1$ rounds of relaxation suffice.

If any distance can still be improved in an additional round, a reachable negative cycle exists.

\subsubsection{Bellman--Ford algorithm}

\begin{enumerate}
    \item Initialize the source distance to zero.
    \item Initialize other distances to infinity.
    \item Relax all edges up to $V-1$ times, stopping early if no update occurs.
    \item Perform one additional pass over the edges.
    \item If an edge can still improve a distance from a reachable vertex, report a negative cycle reachable from the source.
\end{enumerate}

\subsection{Distinction: Reachable vs. Graph-Wide Negative Cycle}

A standard source-based Bellman--Ford implementation detects negative cycles reachable from the chosen source.

However, the question may ask:

\textbf{"Does the graph contain any negative cycle, including one in a disconnected component?"}

For that version, initialize every distance to zero. This is equivalent to adding a hypothetical super-source connected to every vertex by a zero-weight edge.

Then run the Bellman--Ford relaxation process.

This ensures that every component is considered, even if the original graph is disconnected.

\textbf{For the problem as worded here, "contains a negative cycle", the implementation below checks the entire graph.}

\subsection{C++ Implementation: Graph-Wide Negative Cycle Detection}

\begin{lstlisting}[style=cppstyle]
#include <iostream>
#include <vector>
using namespace std;

struct Edge {
    int from;
    int to;
    long long weight;
};

bool hasNegativeCycle(int V, const vector<Edge>& edges) {
    // Initializing all distances to 0 is equivalent to
    // adding a super-source connected to every vertex
    // by a zero-weight edge.
    vector<long long> dist(V, 0);

    for (int i = 0; i < V; ++i) {
        bool changed = false;

        for (const Edge& e : edges) {
            long long candidate = dist[e.from] + e.weight;

            if (candidate < dist[e.to]) {
                dist[e.to] = candidate;
                changed = true;

                // An update during the Vth pass indicates
                // a negative cycle somewhere in the graph.
                if (i == V - 1) {
                    return true;
                }
            }
        }

        if (!changed) {
            return false;
        }
    }

    return false;
}

int main() {
    int V = 4;
    vector<Edge> edges = {
        {0, 1, 1},
        {1, 2, 1},
        {2, 1, -3},
        {2, 3, 2}
    };

    if (hasNegativeCycle(V, edges)) {
        cout << "Negative cycle exists\n";
    } else {
        cout << "No negative cycle\n";
    }

    return 0;
}
\end{lstlisting}

Output:

\begin{lstlisting}[style=pseudocode]
Negative cycle exists
\end{lstlisting}

The graph contains the cycle:


\[
1\rightarrow2\rightarrow1
\]


Its total weight is:


\[
1+(-3)=-2
\]


Therefore, it is a negative cycle.

\subsubsection{Complexity}

Each relaxation pass examines all $E$ edges, and at most $V$ passes are required.

Therefore:


\[
\boxed{\text{Time}=O(VE)}
\]


The distance array occupies:


\[
\boxed{\text{Auxiliary space}=O(V)}
\]


The edge list occupies $O(E)$ storage, but it is input to the algorithm rather than additional working storage.

\subsubsection{Correctness intuition}

If no negative cycle exists, shortest paths can be represented by paths of at most $V-1$ edges. After sufficiently many relaxation rounds, no further improvement is possible.

A negative cycle permits continued improvements as its negative total cost is traversed repeatedly. An improvement in the $V$-th pass therefore signals a negative cycle under the stated initialization and graph assumptions.

\subsubsection{Common examination traps}

\begin{itemize}
    \item Checking only $V-1$ passes is not sufficient to detect a negative cycle; an additional check is necessary.
    \item Source-based Bellman--Ford does not detect negative cycles in unreachable components.
    \item Initializing all distances to zero is a standard technique for detecting a negative cycle anywhere in a graph.
    \item An undirected negative-weight edge, when modeled as two opposite directed edges with the same negative weight, produces a negative closed walk by traversing the edge back and forth. Negative-cycle questions commonly assume directed graphs unless otherwise specified.
    \item Bellman--Ford supports negative edge weights, but no finite shortest-path solution exists for vertices reachable from a relevant negative cycle.
\end{itemize}


\newpage
\begin{center}
\vspace{0.5em}
{\color{primary}\rule{0.45\textwidth}{1.2pt}}\\[4pt]
{\Large\bfseries\color{primary} Section V: Recursion and Dynamic Programming Problems}\\[2pt]
{\color{primary}\rule{0.45\textwidth}{1.2pt}}
\vspace{0.5em}
\end{center}

\section{Recursion and Dynamic Programming Problems}


\subsection{Problem 15. Calculate the sum of digits of an integer using recursion}

\textbf{Question:} Write a recursive program to calculate the sum of digits of a nonnegative integer.

For example:


\[
4827\rightarrow4+8+2+7=\boxed{21}
\]


\subsubsection{Algorithm}

The last digit of an integer $n$ is:


\[
n\bmod10
\]


Removing that digit gives:


\[
\left\lfloor\frac n{10}\right\rfloor
\]


Therefore, the sum of digits can be defined recursively:


\[
\boxed{ S(n)= \begin{cases} n,&0\leq n<10\\\ (n\bmod10)+S(\lfloor n/10\rfloor),&n\geq10 \end{cases} }
\]


\subsubsection{C++ implementation}

\begin{lstlisting}[style=cppstyle]
#include <iostream>
using namespace std;

long long sumDigits(long long n) {
    if (n < 10) {
        return n;
    }

    return n % 10 + sumDigits(n / 10);
}

int main() {
    long long n;
    cin >> n;

    if (n < 0) {
        cout << "Enter a nonnegative integer\n";
        return 0;
    }

    cout << sumDigits(n) << '\n';

    return 0;
}
\end{lstlisting}

Input:

\begin{lstlisting}[style=pseudocode]
4827
\end{lstlisting}

Output:

\begin{lstlisting}[style=pseudocode]
21
\end{lstlisting}

\subsubsection{Complexity}

An integer with $d$ decimal digits requires one recursive call per digit.

\begin{itemize}
    \item \textbf{Time:} $O(d)$.
    \item \textbf{Auxiliary space:} $O(d)$ for the recursion stack.
\end{itemize}

Since $d=O(\log n)$ for $n>0$, the time complexity can also be written as $O(\log n)$.

\textbf{Exam variation:} Reverse the digits of a number, count its digits, or determine whether its digits form a palindrome.

\subsection{Problem 16. Calculate Fibonacci using Dynamic Programming}

\textbf{Question:} Find the $N$-th Fibonacci number using Dynamic Programming.

The Fibonacci sequence is defined by:


\[
F(0)=0,\qquad F(1)=1
\]



\[
\boxed{F(n)=F(n-1)+F(n-2)}
\]


For $N=10$:


\[
F(10)=\boxed{55}
\]


\subsubsection{Iterative DP solution}

\begin{lstlisting}[style=cppstyle]
#include <iostream>
using namespace std;

long long fibonacci(int n) {
    if (n <= 1) {
        return n;
    }

    long long previous2 = 0;
    long long previous1 = 1;

    for (int i = 2; i <= n; ++i) {
        long long current = previous1 + previous2;
        previous2 = previous1;
        previous1 = current;
    }

    return previous1;
}

int main() {
    int n;
    cin >> n;

    if (n < 0) {
        cout << "N must be nonnegative\n";
        return 0;
    }

    cout << fibonacci(n) << '\n';

    return 0;
}
\end{lstlisting}

\subsubsection{Complexity}

\begin{itemize}
    \item \textbf{Time:} $O(n)$.
    \item \textbf{Auxiliary space:} $O(1)$.
\end{itemize}

This is a space-optimized Dynamic Programming implementation. It stores only the two previous values because the recurrence requires no earlier state.

\textbf{Important:} The \texttt{long long} type cannot represent arbitrarily large Fibonacci numbers. An implementation intended for very large $N$ needs arbitrary-precision arithmetic or another representation.

\subsubsection{Comparison of approaches}

\begin{table}[htbp]
\centering
\caption{Comparison of Approaches for Fibonacci Calculation}
\label{tab:fibonacci_approaches}
\begin{tabularx}{\textwidth}{p{4.5cm} X X}
\toprule
\textbf{Method} & \textbf{Time Complexity} & \textbf{Auxiliary Space} \\
\midrule
Naive Recursion & Exponential, $\Theta(\varphi^n)$ & $O(n)$ Stack \\
Memoization (Top-down) & $O(n)$ & $O(n)$ Cache + Stack \\
Full-table Tabulation & $O(n)$ & $O(n)$ Array \\
Two-variable DP & $O(n)$ & $O(1)$ Space Optimized \\
\bottomrule
\end{tabularx}
\end{table}

This is a standard example for comparing recursion and Dynamic Programming.

\subsection{Problem 17. Calculate factorial iteratively}

\textbf{Question:} Find $N!$ without recursion.

The factorial of a nonnegative integer $N$ is:


\[
\boxed{N!=1\times2\times3\times\cdots\times N}
\]


with:


\[
\boxed{0!=1}
\]


\subsubsection{C++ implementation}

\begin{lstlisting}[style=cppstyle]
#include <iostream>
using namespace std;

unsigned long long factorial(int n) {
    unsigned long long result = 1;

    for (int i = 2; i <= n; ++i) {
        result *= static_cast<unsigned long long>(i);
    }

    return result;
}

int main() {
    int n;
    cin >> n;

    if (n < 0) {
        cout << "N must be nonnegative\n";
        return 0;
    }

    cout << factorial(n) << '\n';

    return 0;
}
\end{lstlisting}

For $N=5$, the result is:


\[
5!=\boxed{120}
\]


\subsubsection{Complexity}

\begin{itemize}
    \item \textbf{Time:} $O(n)$.
    \item \textbf{Auxiliary space:} $O(1)$.
\end{itemize}

The iterative implementation avoids recursion-stack growth, but its numerical result can still overflow the chosen integer type.


\newpage
\begin{center}
\vspace{0.5em}
{\color{primary}\rule{0.45\textwidth}{1.2pt}}\\[4pt]
{\Large\bfseries\color{primary} Section VI: Additional High-Value Programming Problems}\\[2pt]
{\color{primary}\rule{0.45\textwidth}{1.2pt}}
\vspace{0.5em}
\end{center}

\section{Additional High-Value Programming Problems}


\subsection{Problem 18. Check whether a number is a palindrome}

A \textbf{palindrome number} reads the same forward and backward.

Examples:

\begin{itemize}
    \item \texttt{121} is a palindrome.
    \item \texttt{1221} is a palindrome.
    \item \texttt{123} is not a palindrome.
\end{itemize}

\subsubsection{Algorithm}

\begin{enumerate}
    \item Store the original number.
    \item Reverse its digits using the remainder and integer-division operations.
    \item Compare the reversed number with the original.
\end{enumerate}

\subsubsection{C++ implementation}

\begin{lstlisting}[style=cppstyle]
#include <iostream>
using namespace std;

bool isPalindrome(long long n) {
    if (n < 0) {
        return false;
    }

    long long original = n;
    long long reversed = 0;

    while (n > 0) {
        int digit = static_cast<int>(n % 10);
        reversed = reversed * 10 + digit;
        n /= 10;
    }

    return original == reversed;
}

int main() {
    long long n;
    cin >> n;

    cout << (isPalindrome(n) ? "Palindrome" : "Not palindrome")
         << '\n';

    return 0;
}
\end{lstlisting}

Input:

\begin{lstlisting}[style=pseudocode]
1221
\end{lstlisting}

Output:

\begin{lstlisting}[style=pseudocode]
Palindrome
\end{lstlisting}

Under this implementation, zero is considered a palindrome, while negative numbers are not.

\textbf{Complexity:} $O(d)$ time, where $d$ is the number of digits, and $O(1)$ auxiliary space.

\textbf{Important limitation:} The reversed number can overflow for sufficiently large input values. A production implementation should check overflow during reversal or use a different approach.

\subsection{Problem 19. Check whether a number is an Armstrong number}

An \textbf{Armstrong number} in base 10 is a nonnegative integer equal to the sum of its digits, each raised to the power of the number of digits.

For example:


\[
153=1^3+5^3+3^3
\]



\[
153=1+125+27=153
\]


Therefore, 153 is an Armstrong number.

Other examples include 0, 1, 370, 371, and 407.

\subsubsection{Algorithm}

\begin{enumerate}
    \item Count the number of digits.
    \item Extract each digit.
    \item Raise the digit to the required power.
    \item Sum these powered digits.
    \item Compare the sum with the original number.
\end{enumerate}

\subsubsection{C++ implementation}

\begin{lstlisting}[style=cppstyle]
#include <iostream>
using namespace std;

long long integerPower(int base, int exponent) {
    long long result = 1;

    for (int i = 0; i < exponent; ++i) {
        result *= base;
    }

    return result;
}

bool isArmstrong(long long n) {
    if (n < 0) {
        return false;
    }

    long long original = n;
    long long temp = n;
    int digits = 0;

    do {
        ++digits;
        temp /= 10;
    } while (temp > 0);

    long long sum = 0;
    temp = n;

    do {
        int digit = static_cast<int>(temp % 10);
        sum += integerPower(digit, digits);
        temp /= 10;
    } while (temp > 0);

    return sum == original;
}

int main() {
    long long n;
    cin >> n;

    cout << (isArmstrong(n) ? "Armstrong" : "Not Armstrong")
         << '\n';

    return 0;
}
\end{lstlisting}

Input:

\begin{lstlisting}[style=pseudocode]
153
\end{lstlisting}

Output:

\begin{lstlisting}[style=pseudocode]
Armstrong
\end{lstlisting}

\subsubsection{Complexity}

For a number with $d$ digits, the algorithm performs $d$ digit-processing steps, with each digit raised to a power of $d$. The simple exponentiation loop takes $O(d)$ work per digit.

Therefore:

\begin{itemize}
    \item \textbf{Time:} $O(d^2)$ under this implementation.
    \item \textbf{Auxiliary space:} $O(1)$.
\end{itemize}

The implementation assumes intermediate values fit in \texttt{long long}. For larger inputs, overflow checks or a big-integer representation would be necessary.

\subsection{Problem 20. Calculate **$a^n$** using fast exponentiation}

\textbf{Question:} Calculate an integer power:


\[
a^n
\]


efficiently.

A straightforward algorithm multiplies $a$ by itself $n$ times, requiring $O(n)$ multiplications.

\textbf{Fast exponentiation}, also called exponentiation by squaring, reduces the number of multiplications to $O(\log n)$.

\subsubsection{Mathematical principle}

If $n$ is even:


\[
\boxed{a^n=(a^{n/2})^2}
\]


If $n$ is odd:


\[
\boxed{a^n=a(a^{\lfloor n/2\rfloor})^2}
\]


Thus, the exponent can be repeatedly divided by two.

\subsubsection{Example}

Calculate:


\[
2^{10}
\]


Using repeated squaring:


\[
2^2=4
\]



\[
2^4=16
\]



\[
2^8=256
\]


Since $10=8+2$:


\[
2^{10}=2^8\times2^2
\]



\[
\boxed{2^{10}=1024}
\]


\subsubsection{C++ implementation}

\begin{lstlisting}[style=cppstyle]
#include <iostream>
using namespace std;

long long fastPower(long long base, unsigned long long exponent) {
    long long result = 1;

    while (exponent > 0) {
        if (exponent & 1ULL) {
            result *= base;
        }

        exponent >>= 1ULL;

        if (exponent > 0) {
            base *= base;
        }
    }

    return result;
}

int main() {
    long long base;
    unsigned long long exponent;

    cin >> base >> exponent;

    cout << fastPower(base, exponent) << '\n';

    return 0;
}
\end{lstlisting}

Input:

\begin{lstlisting}[style=pseudocode]
2 10
\end{lstlisting}

Output:

\begin{lstlisting}[style=pseudocode]
1024
\end{lstlisting}

\subsubsection{Complexity}

The exponent is halved at every iteration:


\[
\boxed{\text{Time}=O(\log n)}
\]


The algorithm uses a constant number of variables:


\[
\boxed{\text{Auxiliary space}=O(1)}
\]


The code assumes that the final result and intermediate multiplications fit in \texttt{long long}. For large powers, use modular exponentiation or arbitrary-precision arithmetic.

\subsubsection{Related problem: Modular exponentiation}

When only $a^n\bmod M$ is needed, reduce the result and base modulo $M$ during the algorithm.

This avoids unnecessarily constructing the enormous value $a^n$. With careful arithmetic, modular exponentiation also runs in $O(\log n)$ multiplications.

It is widely useful in number theory and algorithmic problems involving large integer calculations.


\newpage
\section{Negative-Cycle Detection: Source-Specific Variation}


A common examination variation asks:

\textbf{"Determine whether there is a negative-weight cycle reachable from a particular source vertex."}

This is different from checking whether any negative cycle exists anywhere in the graph.

The earlier implementation initialized every distance to zero and therefore considered all graph components. For a source-specific test, use the following initialization and reachability condition.

\begin{lstlisting}[style=pseudocode]
dist[source] = 0
dist[all other vertices] = infinity

Relax all edges up to V - 1 times,
but only if dist[edge.from] is finite.

For every edge u -> v with weight w:
    If dist[u] is finite AND dist[u] + w < dist[v]:
        A negative cycle is reachable from the source.
\end{lstlisting}

The corresponding extra-pass condition is:


\[
\boxed{ dist[u]\neq\infty \quad\land\quad dist[u]+w(u,v)<dist[v] }
\]


This guard is essential: an edge from an unreachable vertex must not be treated as evidence of a reachable negative cycle.

\subsubsection{Why the distinction matters}

Suppose the graph contains two disconnected components:

\begin{itemize}
    \item Component 1 has the source vertex and no negative cycle.
    \item Component 2 contains a negative cycle.
\end{itemize}

A source-specific Bellman--Ford check for a source in Component 1 should report no \textit{reachable} negative cycle. The global check, which initializes all distances to zero, should report that the graph does contain a negative cycle.

\textbf{Exam rule:} Read the wording carefully. "A negative cycle reachable from source $s$" and "any negative cycle in the graph" are different questions.


\newpage
\begin{center}
\vspace{0.5em}
{\color{primary}\rule{0.45\textwidth}{1.2pt}}\\[4pt]
{\Large\bfseries\color{primary} Section VII: Formulas, Pitfalls, and Final Review}\\[2pt]
{\color{primary}\rule{0.45\textwidth}{1.2pt}}
\vspace{0.5em}
\end{center}

\section{Master Complexity Reference for Examination Problems}


\begin{table}[htbp]
\centering
\caption{Master Algorithmic Complexity Reference for High-Value Problems}
\label{tab:ch09_complexity_reference}
\begin{tabularx}{\textwidth}{p{4.2cm} p{4.2cm} p{3.2cm} X}
\toprule
\textbf{Problem} & \textbf{Recommended Algorithm} & \textbf{Time Complexity} & \textbf{Auxiliary Space} \\
\midrule
GCD of two integers & Euclid's algorithm & $O(\log\min(a,b))$ & $O(1)$ iteratively \\
LCM of two integers & GCD + division first & $O(\log\min(a,b))$ & $O(1)$ \\
Check whether number is prime & Trial division up to $\sqrt{n}$ & $O(\sqrt{n})$ & $O(1)$ \\
Generate all primes up to $N$ & Sieve of Eratosthenes & $O(N\log\log N)$ & $O(N)$ \\
Prime factorization & Trial division & $O(\sqrt{n})$ worst case & $O(1)$ \\
Sum of even numbers up to $N$ & Closed-form formula & $O(1)$ & $O(1)$ \\
Sum of natural numbers up to $N$ & Closed-form formula & $O(1)$ & $O(1)$ \\
Sum of squares up to $N$ & Closed-form formula & $O(1)$ & $O(1)$ \\
Arithmetic progression sum & Closed-form formula & $O(1)$ & $O(1)$ \\
Geometric progression sum & Closed-form formula & $O(1)$ operations & $O(1)$ \\
Pair-sum in an unsorted array & Sorting + two pointers & $O(n\log n)$ & $O(n)$ with copy \\
Two largest distinct elements & One-pass scan & $O(n)$ & $O(1)$ \\
Minimum-difference pair & Sort + adjacent scan & $O(n\log n)$ & Depends on sort \\
Sum of digits & Recursion or loop & $O(d)$ & $O(d)$ recursive, $O(1)$ iterative \\
Palindrome number & Reverse digits & $O(d)$ & $O(1)$ \\
Fast exponentiation & Binary exponentiation & $O(\log n)$ mults & $O(1)$ iterative \\
Fibonacci ($N$-th) & DP with two variables & $O(n)$ & $O(1)$ \\
Negative cycle anywhere in graph & Bellman--Ford (zero init) & $O(VE)$ & $O(V)$ \\
Negative cycle from source & Source-init Bellman--Ford & $O(VE)$ & $O(V)$ \\
\bottomrule
\end{tabularx}
\end{table}

Here, $d$ is the number of digits in the input integer. The mathematical complexity assumes that integer arithmetic takes constant time and that the required results fit in the selected numeric representation. For arbitrary-precision integers, the cost of arithmetic must also be considered.

\textbf{Implementation note:} In Problem 13, include \texttt{<climits>} to use \texttt{LLONG\_MAX}, or use \texttt{numeric\_limits<long long>::max()} from \texttt{<limits>}.


\section{Common Errors to Avoid in Programming Examinations}


\begin{quote}
\textbf{1. Confusing GCD and LCM.}  GCD is the greatest common divisor; LCM is the least common multiple. The relationship between them provides an efficient way to calculate either.
\end{quote}

\begin{quote}
\textbf{2. Mishandling the inputs 0 and 1 in prime checking.}  Neither is prime. A prime number must be greater than 1.
\end{quote}

\begin{quote}
\textbf{3. Using trial division unnecessarily for large collections of primes.}  Use the Sieve of Eratosthenes when all primes up to a limit are required.
\end{quote}

\begin{quote}
\textbf{4. Forgetting integer overflow.}  A formula can be mathematically correct while its C++ implementation produces an incorrect result because intermediate multiplication exceeds the integer type's range.
\end{quote}

\begin{quote}
\textbf{5. Confusing a loop with a closed-form formula.}  Summing every even number up to $N$ takes $O(N)$ time with a simple loop, but the formula $k(k+1)$, where $k=\lfloor N/2\rfloor$, takes $O(1)$ arithmetic operations.
\end{quote}

\begin{quote}
\textbf{6. Accidentally reusing an element in Pair Sum.}  Use distinct positions. The two-pointer condition \texttt{left < right} ensures that a single position cannot be selected twice.
\end{quote}

\begin{quote}
\textbf{7. Losing original indices when sorting.}  If the question requests the positions of a pair, store each value with its original index before sorting.
\end{quote}

\begin{quote}
\textbf{8. Claiming Bellman--Ford detects all negative cycles from an arbitrary source.}  Source-initialized Bellman--Ford detects negative cycles reachable from that source. Initialize every distance to zero when checking for negative cycles anywhere in the graph.
\end{quote}

\begin{quote}
\textbf{9. Confusing negative edges with negative cycles.}  A negative edge does not automatically imply a negative cycle. A negative cycle requires a closed directed path whose total weight is negative.
\end{quote}

\begin{quote}
\textbf{10. Replacing an algorithm's complexity with a single generic result.}  A two-pointer scan is $O(n)$ on sorted input, while sorting an unsorted array first makes the complete Pair-Sum solution $O(n\log n)$.
\end{quote}

\begin{quote}
\textbf{11. Ignoring the numerical limits of factorial and Fibonacci.}  Their mathematical values grow rapidly, and ordinary C++ integer types eventually overflow.
\end{quote}

\begin{quote}
\textbf{12. Misunderstanding DP state reuse.}  The purpose of memoization and tabulation is to avoid repeatedly solving the same subproblems. Correct state definitions and transitions matter more than merely creating an array called \texttt{dp}.
\end{quote}


\section{Final Practical Competency Checklist}


Before completing the Algorithms chapter, you should be able to implement each of the following without consulting the solutions:

\begin{table}[htbp]
\centering
\caption{Practical Algorithmic Competency Checklist for Recruitment Candidates}
\label{tab:recruitment_checklist}
\begin{tabularx}{\textwidth}{p{6.5cm} X}
\toprule
\textbf{Practice Task} & \textbf{Core Competency / Concept Tested} \\
\midrule
Calculate GCD and LCM & Euclid's algorithm and overflow-safe division \\
Check whether a number is prime & Trial division boundary $\sqrt{n}$ and edge cases ($0, 1$) \\
Generate primes up to $N$ & Sieve of Eratosthenes memory array \\
Find prime factors & Number theory and factorization termination \\
Calculate the sum of even numbers & Closed-form mathematical formula versus $O(N)$ loop \\
Generate arithmetic / geometric progression & Numerical sequences and power overflow mitigation \\
Find a pair with a target sum & Sorting + two pointers ($O(n\log n)$) versus hash set \\
Find the two largest distinct elements & One-pass array scan tracking strict inequality \\
Select the pair with minimum difference & Sorting and adjacent neighbor comparison \\
Calculate factorial and Fibonacci & Iterative state compression vs recursion overhead \\
Calculate powers efficiently & Fast exponentiation by repeated squaring ($O(\log n)$) \\
Detect negative cycle in weighted graph & Global Bellman--Ford edge relaxation and 0-init \\
Find shortest path with nonnegative weights & Dijkstra's algorithm with min-heap priority queue \\
Select non-overlapping activities & Greedy optimization sorting by earliest finish time \\
Solve 0/1 Knapsack & Dynamic Programming state space and reverse loop \\
Evaluate a recurrence & Master Theorem and recurrence tree bounds \\
\bottomrule
\end{tabularx}
\end{table}
